% concatenated from main.tex, sections/abstract.tex, sections/intro.tex, sections/related_work.tex, sections/contributions.tex, sections/algorithm.tex, sections/assumptions_and_preliminaries.tex, sections/convergence.tex, sections/complexity_goldstein.tex, sections/numerical_results.tex, sections/conclusions.tex, sections/appendix.tex, sections/appendix_estimates.tex
\documentclass[11pt,a4paper]{article}

\voffset=-1.5cm \hoffset=-1.4cm \textwidth=16cm \textheight=22.0cm

\usepackage{graphicx}
\usepackage{subcaption} % for subfigures
\usepackage{amsmath, amsfonts, amssymb}
\usepackage{enumerate}
\usepackage{longtable}
\usepackage{multirow}
\usepackage{color}
\usepackage{url}
\usepackage{caption}
\usepackage{subcaption}
\usepackage{enumitem}
\usepackage{dsfont}
\usepackage{algorithm}        % must be before algpseudocode
\usepackage{algpseudocode}
\usepackage{mathtools}
\usepackage{hyperref}
\usepackage[normalem]{ulem}
% \usepackage{bbm}    % or dsfont

\newtheorem{theorem}{Theorem}[section]
\newtheorem{acknowledgement}{Acknowledgement}[section]
% \usepackage[ruled,vlined,noline,linesnumbered]{algorithm2e}
\newtheorem{axiom}{Axiom}[section]
\newtheorem{case}{Case}[section]
\newtheorem{claim}{Claim}[section]

\newtheorem{conclusion}{Conclusion}[section]
\newtheorem{condition}{Condition}[section]
\newtheorem{conjecture}{Conjecture}[section]
\newtheorem{corollary}{Corollary}[section]
\newtheorem{criterion}{Criterion}[section]
\newtheorem{definition}{Definition}[section]
\newtheorem{example}{Example}[section]
\newtheorem{exercise}{Exercise}[section]
\newtheorem{lemma}{Lemma}[section]
\newtheorem{notation}{Notation}[section]
\newtheorem{problem}{Problem}[section]
\newtheorem{proposition}{Proposition}[section]
\newtheorem{remark}{Remark}[section]
\newtheorem{solution}{Solution}[section]
\newtheorem{assumption}{Assumption}[section]
\newtheorem{summary}{Summary}[section]
\newtheorem{note}{Note}[section]
\newtheorem{doubt}{Doubt}[section]
\newtheorem{properties}{Properties}[section]

% \newenvironment{proof}[1][Proof]{\textbf{#1.} }{\ \rule{0.5em}{0.5em} \vspace{1ex}}

\newenvironment{proof}[1][Proof]{%
  \par\noindent\textbf{#1.}\ %
}{%
  \par\vspace{-0.5 ex}            % <-- space of roughly one line‐to‐line
  \noindent\hfill\rule{0.5em}{0.5em}\par
  \vspace{1ex}
}

\newcommand{\Sc}[2]{\langle#1, #2\rangle}
\newcommand{\ud}{\, \mathrm{d}}
\newcommand{\st}{\owns_k^{\prime}}
\newcommand{\djs}{\hat{d}_j}
\newcommand{\gammaj}{\gamma_j}
\newcommand{\sm}{\setminus}
\newcommand{\tx}[1]{\textnormal{#1}}
\newcommand{\PP}{\ensuremath{\mathsf P}}
\newcommand{\OM}{\Omega}
\newcommand{\R}{\mathbb{R}}
\newcommand{\HH}{\ensuremath \mathcal{H}}
\newcommand{\n}[1]{\|#1\|}
\newcommand{\A}{\mathcal{A}}
\newcommand{\y}{y_j}
\newcommand{\MM}{\ensuremath \mathcal{M}}
\newcommand{\N}{\ensuremath \mathbb{N}}
\newcommand{\prox}{\ensuremath{\text{\rm prox}}}
\newcommand{\dom}{\ensuremath{\text{\rm dom}}}
\newcommand{\ef}{\varepsilon_f}
\newcommand{\eq}{\varepsilon_q}
\newcommand{\F}{\mathcal{F}}
\newcommand{\lf}{L_f}
\newcommand{\Pb}{\mathbb{P}}
\newcommand{\mo}{\mathds{1}}
\newcommand{\me}{\mathbb{E}}
\newcommand{\LG}{L_{\Gamma}}
\newcommand{\tG}{\tilde{\Gamma}}
\newcommand{\td}[1]{\tilde{#1}}
\newcommand{\Lr}{L_r}
\newcommand{\Lf}{L_f}
\newcommand{\grad}{\tx{grad}}
\newcommand{\tP}{\Gamma}
\newcommand{\Sb}{\mathbb{S}}
\setlength{\unitlength}{1mm}
\newcommand{\eg}{\varepsilon_g}
\newcommand{\eh}{\varepsilon_h}
\newcommand{\dd}{\varepsilon_f}
\newcommand{\eT}{\eta}
\newcommand{\tD}{\Delta}
\newcommand{\f}{\tilde{f}}
\newcommand{\D}{\Delta}
\newcommand{\Imax}{\bar{\imath}}%I_{\max}
\newcommand{\tPhi}{\tilde{\Phi}}
\newcommand{\tdelta}{\Delta}
\newcommand{\Fp}{\F_{k - 1}^+}
\def\real{\mathbb{R}}

% exponent for Delta random variable
\newcommand{\Dexp}{p}
% exponent for alpha random variable in tails assumption
\newcommand{\aexp}{\Dexp/(\Dexp-1)}
% golstein delta radius
\newcommand{\goldelta}{r}

\DeclareMathOperator{\argmax}{argmax}
\DeclareMathOperator{\argmin}{argmin}

\newcommand{\defeq}{\coloneqq}
\newcommand{\ip}[2]{\left\langle #1, #2 \right\rangle}
\newcommand{\norm}[1]{\left\lVert #1 \right\rVert}
\newcommand{\fik}{f^{(i)}_{k}}
\newcommand{\tfik}{\tilde{f}^{(i)}_{k}}
\newcommand{\Dik}{\Delta^{(i)}_{k}}
\newcommand{\tDik}{\Delta^{(i)}_{k}}
\newcommand{\fks}{f_k(s)}
\newcommand{\fko}{f_k^0}
\newcommand{\dk}{\delta_k}
\newcommand{\xk}{x_k}
\newcommand{\Xk}{X_k}

\newcommand{\ind}{\mathbf{1}}   % “ind” for indicator
\newcommand{\one}{\mathbf{1}}   % “one” for the 1

% \newcommand{\EE}{\ensuremath{\mathsf E}}
\newcommand{\EE}{\mathbb{E}}

\newcommand{\sara}[1]{{\color{red}{\textbf{[SS: #1]}}\color{black}}}
\newcommand{\saraedit}[1]{{\color{red}{#1}\color{red}}}
\newcommand{\Ant}[1]{\textcolor{magenta}{#1}}


\newcounter{countana}


\DeclareMathOperator{\dist}{dist}
\title{Extrapolation-based Direct Search for Nonsmooth Stochastic Zeroth-Order Optimization}

	\author{
    A.~Palmieri \thanks{Dipartimento di Matematica, Universit\`a
		di Padova, Italy
		(\tt{anthony.palmieri@studenti.unipd.it}).}
	\and
	F.~Rinaldi \thanks{Dipartimento di Matematica, Universit\`a
		di Padova, Italy
		(\tt{rinaldi@math.unipd.it}).}
    \and
    S.~Shashaani
    \thanks{Edward P. Fitts Department of Industrial and System Engineering, North Carolina State University, Raleigh, NC, USA
		(\tt{sshasha2@ncsu.edu}).}
}

\begin{document}
	\maketitle

%\begin{abstract}
% \noindent
We propose and analyze a stochastic direct-search method for unconstrained
zeroth-order minimization of locally Lipschitz, possibly nonsmooth, objectives.
The method combines random polling directions with a stochastic extrapolating
line search based on a sufficient-decrease test of order \(\Dexp\). Under
conditional accuracy assumptions on the stochastic estimates, 
% %which can be
% verified for mean-zero finite-higher-moment oracle noise through suitable sample
% averaging, 
we prove almost-sure convergence to Clarke
stationary points. We further establish an expected iteration complexity
bound. Specifically, using a supermartingale stopping-time argument, we prove
that
$\mathcal O\left(
        \max\left\{
            \goldelta^{-\Dexp},
            \varepsilon^{-\Dexp/(\Dexp-1)}
        \right\}
    \right)
$
iterations are sufficient in expectation to reach an
\((\goldelta,\varepsilon)\)-Goldstein stationary point. Moreover, we derive a corresponding expected tested-point complexity bound of order
\(\mathcal O\bigl(\varepsilon^{1-n}
\max\{r^{-p},\varepsilon^{-p/(p-1)}\}\bigr)\). To the best of
our knowledge, this is the first convergence and expected-complexity analysis
for an \emph{extrapolation-based direct-search} method in a \emph{nonsmooth
stochastic} setting. Numerical experiments on a DFO benchmark suite highlight competitive performance against well-established stochastic direct-search methods.
%\end{abstract}

% \AnthonyNew{
% \begin{abstract}
% We propose and analyze a direct-search scheme for unconstrained stochastic zeroth-order minimization of locally Lipschitz, possibly nonsmooth, objectives that combines a stochastic line search with random dense search directions. We impose minimal assumptions on the stochastic oracle, that is, zero-mean and finite-variance noise with no variance-stepsize coupling. Under these conditions, we prove almost-sure convergence to Clarke stationary points, while also establishing some preliminary complexity results. Numerical experiments on a DFO benchmark suite highlight competitive performance against well-established stochastic direct-search methods.
% \end{abstract}
% }

% introduction
\section{Introduction} \label{sec:intro}
We study the unconstrained problem
\begin{equation}\label{gen_prob}
  \min_{x\in\R^n} f(x),
\end{equation}
where $f:\R^n\to\R$ is a \emph{locally} Lipschitz (possibly nonsmooth) function that cannot be evaluated exactly. 
% this entails no loss of generality on the (assumed) compact iterate set $\mathcal{C}$ (Assumption~\ref{ass:compact}), since local Lipschitzness implies the existence of a finite constant $L$ with
% \[
% |f(x)-f(y)| \le L\,\|x-y\|\qquad\text{for all }x,y\in \mathcal{C}.
% \]
% Thus, all bounds in the proofs depend only on such a uniform constant on $\mathcal{C}$.\\
%In our setting, $f$  
Instead, we only have access  to a \emph{stochastic
zeroth-order oracle} that returns a random estimate $\tilde f(x)$ of $f(x)$ at any query point $x$. For example, %In some contexts, 
the stochastic estimate can be modeled as a random variable
parameterized by $x$, namely
\[
  \tilde f(x)\;=\;F(x,\zeta),\qquad \zeta\sim\mathbb P,
\]
i.e., the oracle draws a random seed $\zeta$  from an underlying probability space
$(\Omega,\mathcal F,\mathbb P)$ and outputs $F(x,\zeta)$. 
% No gradient or subgradient information isassumed available. 
Throughout the paper, all algorithmic quantities are random variables adapted to
the natural filtration induced by the oracle and the algorithmic randomness (precise assumptions on
the estimator, e.g., probabilistic accuracy, tails, and independence are stated in
Section~\ref{sec:assumptions}).
Our goal is to design and analyze a \emph{derivative-free} method that combines random dense directions with a
stochastic line search, proving convergence to Clarke stationary points and establishing expected
complexity bounds in terms of oracle calls.

\subsection{Related Work} \label{subsec:related_work}
Derivative-free optimization (DFO) focuses on  problems whose derivative information is unavailable, unreliable, or simply too expensive to obtain, requiring algorithms to work only with (possibly noisy) function values. Two main approaches are common in this context, which are the model-based methods and the direct-search methods. In model-based methods, one builds at each iteration a local surrogate  intended to approximate $f$ on a neighborhood of the current point, proposes a step by (approximately) minimizing the model inside that neighborhood, and accepts or rejects it by comparing predicted and observed decrease. Direct-search methods, by contrast, do not form an explicit model, but rather evaluate $f$ on a structured set of trial points 
%of the form $x_k+\Delta_k d$ (with $d$ taken from a sufficiently rich direction set) 
and update the current point based on a tailored acceptance test (e.g., sufficient-decrease of the objective function).
The deterministic DFO theory is well established; see, e.g., the monograph of Conn, Scheinberg, and Vicente~\cite{MR2487816}, and the survey of Larson, Menickelly, and Wild~\cite{MR3963507} for further details. 
% When $f$ is nonsmooth, the natural optimality notions rely on generalized derivatives (notably Clarke stationarity) together with the associated calculus~\cite{Clarke1990,Jahn1996}. 
Both model-based and direct-search schemes extend to stochastic zeroth-order oracles, but the analysis must then keep track of how estimation error affects acceptance decisions and step-size updates.\\
On the model-based side, trust-region methods based on probabilistic accuracy conditions~\cite{MR3245880} demonstrate that much of the classical trust-region mechanism remains valid when local models and function estimates are sufficiently accurate with %are correct with sufficiently 
high probability, and martingale-based arguments play a central role in the theoretical analysis. This line of work ultimately led to the development of STORM-type frameworks, which systematically integrate randomized local models with stochastic function evaluations within a trust-region mechanism~\cite{MR3800867}. Expected-complexity guarantees for STORM-type methods can be derived through stopping-time and supermartingale techniques~\cite{MR4151319}, and this analysis has been adapted both to random-subspace trust-region schemes for large-scale problems~\cite{MR4777851} and to stochastic line-search mechanisms~\cite{MR4060460}. In a separate framework for derivative-free stochastic trust-region schemes, adaptive sampling allows to %balance
drive down the estimation error according to a measure of stationarity, ensuring that the local models and function estimates are eventually sufficiently accurate almost surely; ASTRO-DF is a representative example of an algorithm built on this principle~\cite{MR3880261}. Subsequent work establishes almost-sure iteration and sample complexity results for ASTRO-DF and quantifies how variance-reduction devices, such as common random numbers (CRN), can improve sample complexity of these methods by replacing stringent per-point accuracy requirements with the difference accuracy requirements~\cite{MR4959972}. A related development %replaces stringent per-point accuracy requirements by means of 
uses similar difference error tail bounds within the supermartingle complexity analysis to obtain %, which allow 
reduced sample sizes (with further benefits under CRN) while preserving convergence guarantees for both trust-region and direct-search methods~\cite{MR4758401}.\\
On the direct-search side, StoMADS extends MADS to stochastic black-box objectives through probabilistic estimates and variance control, and uses a mesh-based strategy to obtain almost-sure convergence to Clarke stationary points~\cite{MR4238147}, although with no complexity analysis. %. No expected-complexity analysis is included in the paper. 
Constrained variants incorporate progressive barriers and probabilistic feasibility control~\cite{MR4550962}. For smooth objectives, the work by Dzahini~\cite{MR4359469} provides an expected-complexity analysis for stochastic direct search under a power-type sufficient-decrease condition, again via a supermartingale-based argument, and a recent survey summarizes current direct-search paradigms and guarantees, including stochastic settings and line-search variants~\cite{MR4919092}. Linesearch-based DFO with extrapolation under noisy oracles has also been studied recently, with convergence and complexity results in smooth regimes~\cite{desantis2025linesearchbasedderivativefreemethodnoisy}. Another zeroth-order paradigm is randomized smoothing, where one replaces the
objective by a smoothed surrogate and suitably build random estimators of the surrogate gradient. This idea is
central to the random gradient-free framework of Nesterov and
Spokoiny~\cite{MR3627456}. For nonsmooth nonconvex objectives, \cite{LinEtAl2022GFM} was among the first works to connect randomized smoothing
with Goldstein stationarity and to analyze gradient-free methods for reaching
\((r,\varepsilon)\)-Goldstein stationary points. This line was later sharpened
by Kornowski and Shamir~\cite{JMLR:kornowskishamir}, who obtained optimal
dimension-dependence for stochastic zero-order nonsmooth nonconvex optimization.
Our work fits into the stochastic direct-search/line-search thread, but targets a fully nonsmooth, derivative-free setting: we couple dense random directions with a run-until-failure line-search policy. We prove almost-sure convergence in the spirit of~\cite{MR4758401}, by relying on Clarke~\cite{Clarke1990} and Goldstein/Lebourg nonsmooth tools~\cite{Goldstein1977,Jahn1996}. Moreover, we derive an expected-complexity result by leveraging the supermartingale framework proposed in~\cite{MR4151319}, while working under minimal oracle assumptions.

% \textit{DFO in general.}
% Derivative-free optimization (DFO) addresses problems where derivatives are unavailable or unreliable and only (possibly noisy) function values can be queried. Comprehensive references include the monographs \cite{MR3729457,MR2487816} and the survey \cite{MR3963507}. Two dominant design philosophies have emerged: \emph{trust-region} (TR) methods, which build local models and regulate steps by a radius, and \emph{direct-search} (DS) methods, which drive progress by directional polling and sufficient-decrease tests without explicit models. Beyond these, zero-order and stochastic-approximation viewpoints quantify baseline rates and dimension effects for convex/nonconvex settings \cite{MR3627456,MR3134439,10.5555/3600270.3602167}. Inexact-oracle regularization frameworks extend complexity guarantees to cubic/ARC schemes with noisy or inexact information \cite{MR4343731,MR4375515,10408598}. Worst-case guarantees for derivative-free TR in nonsmooth regimes further motivate working with minimal smoothness \cite{MR3554884}. Our analysis will rely on martingale tools and nonsmooth calculus, for which standard references are \cite{books/daglib/0073491,clason2025introductionnonsmoothanalysisoptimization}.
% \noindent
% \textit{Stochastic trust-region (model-based) methods.}
% Stochastic TR methods typically assume \emph{probabilistically accurate} models/estimates and establish global convergence or rates under fixed-confidence accuracy. Early probabilistic-model TR and random-model TR frameworks are given by \cite{MR3245880,MR3800867}; ASTRO-DF develops adaptive sampling for interpolation models and proves almost-sure convergence \cite{MR3880261}. A supermartingale/stopping-time framework for expected-complexity of stochastic TR was introduced in \cite{MR4151319} and later adapted to large-scale settings via random subspaces \cite{MR4777851}. Recent analysis sharpens zeroth-/first-order evaluation complexity and clarifies the role of variance reduction and common random numbers (CRN) \cite{MR4959972}. A complementary reduction-centric line replaces per-point accuracy by a \emph{weak tail bound on the reduction estimate}, yielding reduced sample sizing—and further gains under CRN—while retaining global convergence for both TR and DS \cite{MR4758401}. Stochastic line-search variants with fixed-confidence accuracy further connect to this supermartingale viewpoint \cite{MR4060460}. These developments inform the sampling and drift mechanisms we use, although our method is model-free and tailored to nonsmooth settings.
% \noindent
% \textit{Stochastic direct search.}
% In DS, progress is certified via sufficient-decrease tests along directions/patterns. StoMADS handles noisy black-box objectives using probabilistic estimates and variance control, proving almost-sure convergence to Clarke stationary points \cite{MR4238147}; constrained variants adopt progressive barriers with probabilistic feasibility \cite{MR4550962}. For smooth objectives, a stochastic directional DS (SDDS) with a power-type sufficient decrease attains expected-complexity bounds via a supermartingale driver \cite{MR4359469}. A recent survey consolidates guarantees and algorithmic paradigms for modern DS, including stochastic regimes \cite{MR4919092}. Linesearch-based DFO with extrapolation under noise has also been analyzed, providing convergence and iteration bounds \cite{desantis2025linesearchbasedderivativefreemethodnoisy}. Our work belongs to this DS strand: we study a \emph{model-free}, dense-direction line-search in a fully nonsmooth, gradient-free setting; we calibrate sampling through reduction-tail ideas \cite{MR4758401}, verify drift and success probabilities through a Cartis–Scheinberg supermartingale driver \cite{MR4151319}, and use Clarke/Lebourg tools \cite{clason2025introductionnonsmoothanalysisoptimization} to obtain convergence and expected-complexity under minimal assumptions.

% \paragraph{Model-based stochastic DFO via trust regions.}
% A central strand of stochastic derivative-free optimization (SDFO) builds \emph{probabilistically accurate} local models and regulates steps with trust regions. Foundational contributions include trust regions driven by random/probabilistic models and estimation accuracy tests \cite{MR3245880,MR3800867}, and the adaptive-sampling ASTRO-DF framework \cite{MR3880261}. Recent progress sharpens iteration/evaluation complexity for both zeroth- and first-order stochastic trust-region methods, quantifying how sampling and noise regularity affect rates \cite{MR4959972}. These works typically (i) assume smooth or weakly nonsmooth objectives, (ii) construct and validate local models with fixed-probability accuracy, and (iii) couple acceptance with variance/accuracy conditions. \emph{Our algorithm departs from model-based TR: it is model-free, relies on a power-type sufficient-decrease test on \emph{reduction estimates}, and calibrates sampling by \emph{tail bounds} for the reduction error together with an explicit post-selection \(L^1\) control at the selected inner level.}

% \paragraph{Direct-search under noise.}
% On the direct-search side, StoMADS establishes almost-sure convergence to Clarke-stationary points using probabilistic estimates and variance control \cite{MR4238147}. For smooth objectives, Dzahini derives expected-complexity for stochastic directional direct search using a power-type sufficient-decrease test and a supermartingale driver \cite{MR4359469}. Random-subspace variants with trust regions achieve convergence and expected-complexity bounds in large scale \cite{MR4777851}. A recent survey synthesizes theoretical guarantees and algorithmic paradigms for modern direct-search (deterministic and stochastic) \cite{MR4919092}. \emph{Relative to this line, we keep the direct-search flavor but (a) employ \emph{random dense directions} sampled on the sphere, (b) control acceptance through \emph{reduction-tail} calibration (as opposed to fixed-probability accuracy of two separate function values), and (c) add a \emph{post-selection} \(L^1\) bound to handle the selection bias induced by choosing the successful inner level \(H_k\).}

% \paragraph{Tail conditions and reduced sample sizing.}
% Closest to our estimation viewpoint is the weak tail-bound condition of Rinaldi--Vicente--Zeffiro \cite{MR4758401}. By placing a tail bound directly on the \emph{reduction estimate} that appears in the sufficient-decrease test (for both TR and DS), they show that per-iteration sample sizes can be significantly lowered compared with classical finite-variance prescriptions, while retaining global convergence (including nonsmooth settings). \emph{We adopt the same philosophical shift—control the \emph{reduction} rather than the two function values separately—yet operate with random dense directions and explicitly account for post-selection effects at the chosen inner level \(H_k\).}

% \paragraph{Supermartingale and stopping-time analyses.}
% Our complexity analysis leverages the supermartingale/stopping-time framework of Blanchet--Cartis--Menickelly--Scheinberg \cite{MR4151319}, which yields expected stopping-time bounds once a positive drift and a success probability bounded away from \(1/2\) are verified. Related developments for stochastic line search appear in \cite{MR4060460}. \emph{In our setting, the drift is obtained by combining the power-type sufficient-decrease test with a selected-level \(L^1\) control on the reduction-estimate error, while the success probability is ensured through standard inner-loop sampling/capping devices; the resulting Cartis-type driver produces expected evaluation bounds.} For general background on martingale tools we refer to \cite{books/daglib/0073491}.

% \paragraph{Random directions and zeroth-order methods.}
% Random-direction techniques are classical in zero-order optimization. In convex regimes, Gaussian/random directions achieve dimension-dependent rates comparable (up to factors) to gradient-based methods \cite{MR3627456}. These insights motivate our use of \emph{dense} random directions for nonsmooth SDFO, yet—unlike smoothing-based ZO methods—our acceptance is reduction-based with tail-calibrated sampling and post-selection control. We also note connections to random-subspace TR for scalability \cite{MR4777851}.

% \paragraph{Background on nonsmooth analysis (used in stationarity statements).}
% % Our stationarity guarantees are phrased in terms of Clarke subdifferentials; background and references for nonsmooth analysis and optimization can be found in \cite{clason2025introductionnonsmoothanalysisoptimization}.

% Model-based SDFO builds \emph{probabilistically accurate} local models within trust regions and regulates steps via accuracy tests; representative foundations include probabilistic-model TR and random-model TR frameworks \cite{MR3245880,MR3800867} and the adaptive-sampling ASTRO-DF scheme \cite{MR3880261}. Recent analyses sharpen zeroth-/first-order evaluation complexity and clarify how sampling and noise regularity drive rates \cite{MR4959972}. In parallel, direct-search methods under noise obtain a.s.\ convergence and expected-complexity guarantees using probabilistic estimates, notably StoMADS \cite{MR4238147}, the supermartingale-based analysis of stochastic directional direct search \cite{MR4359469}, and scalable random-subspace TR variants \cite{MR4777851}; see also the recent survey of theory and algorithmic paradigms \cite{MR4919092}. A key development for both TR and DS is the weak tail-bound condition, which controls the \emph{reduction estimate} directly (rather than two noisy function values) and thereby reduces sample sizing while preserving global convergence \cite{MR4758401}.

\subsection{Contributions} \label{subsec:contributions}
The proposed approach combines a line search strategy with random directions uniformly generated on the unit sphere. 
%We work with a mean-zero, finite $r$-th moment noise \sara{change from $r$ since it is used in Goldstein stationarity and complexity}, and implement the resulting assumptions through sample-average estimates. This yields a sample complexity $\mathcal O\!\big(\Delta_k^{-2\Dexp}\big),$
%where $\Delta_k$ is the stepsize at iteration $k$. 
% The $\log$ term arises from explicit tail bounds on the random depth of the inner line search, and represents the mild cost of avoiding stronger accuracy hypotheses. Under additional structure commonly assumed in the literature (e.g., fixed-confidence per-evaluation accuracy or weak reduction-tail/CRN conditions), the logarithmic overhead disappears and one recovers the usual algebraic scalings, see, e.g., \cite{MR4060460,MR4758401}.
% We prove almost-sure convergence to Clarke stationary points and, by suitably adapting the analysis given in \cite{MR4151319,MR4060460}, we give an expected iteration complexity bound, and a corresponding tested-point \sara{?} complexity bound, via a supermartingale stopping-time framework. 
Under conditional tail-accuracy assumptions (motivated by the weak tail-bound framework introduced in~\cite{MR4758401}), we prove almost-sure convergence to Clarke stationary points and, by suitably
adapting the analyses in~\cite{MR4151319,MR4060460}, we derive an expected
iteration complexity bound of order
\(\mathcal O(\max\{r^{-p},\varepsilon^{-p/(p-1)}\})\) for reaching an
\((r,\varepsilon)\)-Goldstein stationary point. We also obtain the corresponding
expected tested-point complexity bound
\(\mathcal O(\varepsilon^{1-n}\max\{r^{-p},\varepsilon^{-p/(p-1)}\})\), where
tested-point complexity counts the number of points at which stochastic function
estimates are queried.
The analysis relies exclusively on nonsmooth tools, that is Clarke and Goldstein subdifferentials and the Lebourg mean value theorem (see, e.g., \cite{Clarke1990,Goldstein1977,Jahn1996} for further details on these theoretical tools). To replace gradient-based descent estimates in the complexity argument, we test multiple \emph{dense random directions} per iteration (see, e.g., \ \cite{MR3627456}) and establish a geometric success guarantee: with probability strictly larger than $1/2$, at least one sampled direction falls in a spherical cap aligned with a Goldstein-descent direction. This spherical-cap mechanism is fundamentally different from the positive-spanning arguments used in smooth direct-search analyses (see, e.g., \cite{MR4359469}). %Finally, our assumptions are intentionally light: $f$ is merely locally Lipschitz and the noise is mean-zero with finite $r$-th moment, with no variance-stepsize coupling.
To the best of our knowledge, this is the first convergence and expected-complexity analysis for an \emph{extrapolation-based direct-search} method in a  \emph{nonsmooth} setting.
A recent work by De Santis, Liuzzi and Lucidi~\cite{desantis2025linesearchbasedderivativefreemethodnoisy} studies a stochastic line-search scheme in a smooth setting, assuming continuously differentiable objectives with $L$-Lipschitz gradients and fixed-confidence accuracy for the function estimates; in that framework, the sampling strategy enforces a variance-stepsize coupling (of order $\Delta_k^4$) and the resulting complexity guarantees are stated in terms of the expected gradient norm (e.g., $\mathcal O(\varepsilon^{-2})$).
The present work addresses a complementary regime, focusing on merely locally Lipschitz, fully nonsmooth objectives and purely zeroth-order stochastic information. In contrast to analyses that assume a variance-stepsize coupling directly on the stochastic estimates, we impose conditional polynomial-tail accuracy conditions. %which can be derived from a finite-moment noise model by adaptive sample averaging.} 
 We combine a line search with dense random directions and obtain convergence and expected-complexity results in terms of Clarke/Goldstein stationarity. %\sara{is our work also in terms of expected gradient norm?}. % via a supermartingale stopping-time argument. 
% Under these minimal assumptions, the only additional cost is a mild $\log^2(1/\Delta_k)$ term in the batch schedule, which vanishes under stronger accuracy hypotheses such as fixed-confidence or reduction-tail/CRN conditions.

% We provide, to the best of our knowledge, the first convergence and expected–complexity analysis of a \emph{direct-search, line-search based} algorithm in a fully \emph{nonsmooth, gradient-free} setting. The line search follows a run–until–failure policy with dense random directions, and we introduce a selection-aware estimation scheme that is \emph{implementable} via sample averages under additive, mean-zero, finite-variance noise. This leads to the practical batch schedule
% \[
% m_k=\Theta\!\big(\Delta_k^{-2\Dexp}[1+\log(1/\Delta_k)]^{2}\big).
% \]
% The $\log^2$ factor—stemming from explicit tail bounds we derive for the inner line-search stopping time—is a benign consequence of working under minimal assumptions; under stronger structural conditions used in the literature (e.g., fixed-confidence per-evaluation accuracy or weak reduction-tail/CRN hypotheses \cite{MR4758401}) our batch sizing recovers the standard algebraic scalings exactly (i.e., no $\log^2$ overhead).\\
% On this foundation we establish almost-sure convergence and derive expected evaluation complexity through the Cartis–Scheinberg supermartingale/stopping-time framework \cite{MR4151319} (see also \cite{MR4060460}), in a \emph{purely zeroth-order} manner: both parts of the analysis rely only on nonsmooth tools—Clarke subdifferentials and the Lebourg mean value theorem \cite{clason2025introductionnonsmoothanalysisoptimization}. To compensate for the absence of gradients in the complexity step, we test multiple \emph{dense random directions} per iteration (cf.\ \cite{MR3627456}) and prove a geometric success guarantee (probability $>1/2$ of finding a direction sufficiently aligned with a Clarke-descent direction) via a \emph{spherical-cap} argument on the unit sphere, which appears to be new in the Clarke setting. Our assumptions are intentionally light: $f$ is only locally Lipschitz and the oracle noise is only mean-zero, and finite-variance, with \emph{no} variance–stepsize coupling.

% We analyze a \emph{model-free} direct-search method with dense random directions and an inner \emph{run-until-failure} line-search. This inner policy induces dependence among the random variables used for acceptance and estimation and, in particular, creates \emph{selection bias} in conditional means, which complicates the direct use of sufficient-decrease arguments. Our main contributions are:

% \begin{enumerate}
% \item \textbf{Selection-aware estimation under run-until-failure.}
% Since $H_k$ is a success run-length, the acceptance event depends on the same noise used in the reduction, inducing dependence and \emph{selection bias}. We address this with a single post-selection condition
% \(
% \mathbb{E}\!\left[|E_{k,H_k}|\mid\mathcal F_{k-1}\right]\le c\,\Delta_k^{\Dexp},
% \)
% and show it is implementable via sample averages: tail bounds for $H_k$ as functions of $\Delta_k$ yield a batch schedule $m_k$ that enforces the inequality, with
% \(
% m_k=\Theta\!\left(\Delta_k^{-2\Dexp}\,[1+\log(1/\Delta_k)]^{2}\right)
% \)
% for finite-variance noise (with sharper constants under stronger noise tail assumptions).

% % \item \textbf{Tailored convergence proof for run-until-failure.}
% % We prove that the series of stepsizes satisfies $\sum_k \Delta_k^{\,\Dexp}<\infty$ and $\Delta_k\to 0$ almost surely with an argument tailored to the run-until-failure policy (no inner cap), which avoids altering the algorithm to simplify analysis.

% \item \textbf{Expected evaluation complexity via supermartingales, without gradients.}
% We derive expected complexity bounds using the Cartis–Scheinberg supermartingale/stopping-time framework \cite{MR4151319} (cf.\ also \cite{MR4060460}), but in a purely zeroth-order setting: both convergence and complexity rely only on nonsmooth analysis tools (Clarke subdifferential and the Lebourg mean value theorem) \cite{clason2025introductionnonsmoothanalysisoptimization}, not on gradient information.
% % Once a positive drift and a success probability bounded away from $1/2$ are verified; martingale tools are standard \cite{books/daglib/0073491}.

% \item \textbf{Dense random directions with geometric success guarantees.}
% In complexity analysis, to compensate for the lack of gradients, the algorithm tests multiple \emph{random dense directions} per outer iteration (cf.\ classical random-direction methods \cite{MR3627456}). We show that, with probability strictly larger than $1/2$, one of these directions is sufficiently aligned with a Clarke-descent direction; the proof follows Dzahini’s idea \cite{MR4359469} but replaces positive-spanning arguments with \emph{spherical-cap} geometry on the unit sphere. To the best of our knowledge, this geometric treatment in the Clarke setting is new.

% \item \textbf{Light structural assumptions.}
% We require only \emph{local Lipschitz} continuity of $f$—the minimal regularity for Clarke calculus—and an additive mean-zero finite-variance oracle. We do \emph{not} assume that the noise variance scales with stepsize.

% \item \textbf{Noise-model refinements and sampling complexity.}
% We quantify how stronger noise assumptions (e.g., sub-Gaussian tails or finite $r$-th moments) sharpen constants in the sample-size schedule, and we report the resulting impact on overall evaluation complexity. These refinements are compatible with variance-reduction/CRN mechanisms.
% \end{enumerate}

% Our work adopts this reduction-centric view in a \emph{model-free} dense-direction line search: (i) we employ random directions on the sphere (cf.\ classical random-direction ZO methods \cite{MR3627456}); (ii) we calibrate sampling via a \emph{tail} bound on the reduction error; and (iii) we explicitly control the \emph{post-selection} $L^1$ error at the selected inner level to handle selection bias. Expected-evaluation complexity then follows from the supermartingale/stopping-time framework \cite{MR4151319} (see also stochastic line search \cite{MR4060460}), once a positive drift and a success probability bounded away from $1/2$ are verified; martingale tools are standard \cite{books/daglib/0073491}. For stationarity statements in the nonsmooth setting, we follow the Clarke-calculus background in \cite{clason2025introductionnonsmoothanalysisoptimization}.

% Summarizing, in relation to (i) model-based SDFO with probabilistically accurate models \cite{MR3245880,MR3800867,MR3880261,MR4959972} and (ii) direct-search schemes with fixed-probability estimate accuracy \cite{MR4238147,MR4359469,MR4919092,MR4777851}, we contribute:
% \begin{enumerate}
% \item a \emph{random dense directions} direct-search variant with a power-type sufficient-decrease test on \emph{reduction} estimates;
% \item a \emph{tail} calibration of the reduction error, in the spirit of \cite{MR4758401}, that decouples sampling from strong variance assumptions and interacts cleanly with common-random-numbers or variance reduction;
% \item an explicit, implementable \emph{post-selection} \(L^1\) control at the selected inner level \(H_k\), needed to close the drift in the presence of selection bias; and
% \item an expected-complexity analysis via a Cartis-type supermartingale/stopping-time driver \cite{MR4151319,MR4060460}, together with standard inner-loop sampling/capping to ensure a uniform success probability.
% \end{enumerate}

\subsection{Paper Structure} \label{subsec:paper_structure}
The paper proceeds as follows. 
Section~2 presents the algorithm, the stochastic line search, and notation.
Section~3 states the standing assumptions (compactness, conditional independence, and $\Dexp$-tail).
Section~4 proves a uniformly positive conditional expected merit-function decrease proportional  to $\Delta_k^p$, shows $\sum_k \Delta_k^{\Dexp}<\infty$ for $\Dexp >1$ almost surely and hence $\Delta_k\to 0$ almost surely. This establishes convergence to Clarke stationary point almost surely. Section 5 proves that the algorithm reaches an $(\goldelta,\varepsilon)$-Goldstein stationary point in expected $\mathcal O\left(\max\left\{\goldelta^{-\Dexp}, \varepsilon^{-\frac{\Dexp}{\Dexp-1}}\right\}\right)$ iterations %oracle evaluations 
using a supermartingale stopping-time analysis. Furthermore, it derives
the corresponding expected tested-point complexity.
% similar to the one given in \cite{MR4151319}. 
Section~6 presents numerical experiments showing that the proposed method is competitive with state-of-the-art stochastic derivative-free algorithms on a standard nonsmooth benchmark.
Section 7 finally summarizes the theoretical results and discusses possible future extensions of the method. The appendix gathers auxiliary proofs.


% % literature overview
% \input{sections/related_work}

% algorithm presented
\section{Algorithm: Direct Search with Extrapolation for Nonsmooth Objectives} \label{sec:algo}

We consider a direct-search method with extrapolation for possibly nonsmooth
stochastic zeroth-order optimization. The outer scheme is given in
Algorithm~\ref{alg:dse2-cartis}, and the stochastic line-search subroutine is given in
Algorithm~\ref{alg:lsep}. We fix $\Dexp\in (1,2]$, $\theta>0$, $\gamma\in(0,1)$, 
$\bar{m}\in\mathbb N$,  and $\Imax\in\mathbb N$.
%We work on a probability space $(\Omega,\mathcal F,\mathbb P)$. 
Random
quantities are denoted by uppercase letters and their realizations by lowercase
letters. At the beginning of iteration $k$, the current iterate and stepsize are
$X_k\in\mathbb R^n$ and $\Delta_k>0$. The algorithm then samples
$\bar{m}$ candidate directions
$D_{k,1},\ldots,D_{k,\bar{m}}\in\mathbb S^{n-1},
$
where we denote by
\(
    \mathbb S^{n-1}
    :=
    \{d\in\mathbb R^n:\|d\|=1\}
\)
the unit sphere in \(\mathbb R^n\). We use two sigma-algebras. Let $\mathcal G_{k-1}$ denote the information
available before sampling the directions at iteration $k$, so that $X_k$ and
$\Delta_k$ are $\mathcal G_{k-1}$-measurable. Conditionally on
$\mathcal G_{k-1}$, the directions $D_{k,1},\ldots,D_{k,\bar{m}}$ are sampled
independently and uniformly on $\mathbb S^{n-1}$. After sampling them, but
before drawing any stochastic function estimates at iteration $k$, we set
\[
    \mathcal F_{k-1}
    :=
    \mathcal G_{k-1}
    \vee
    \sigma(D_{k,1},\ldots,D_{k,\bar{m}}).
\]
Thus $X_k$, $\Delta_k$, and all directions $D_{k,m}$ are
$\mathcal F_{k-1}$-measurable, while the estimates generated during iteration
$k$ are not. Oracle assumptions are stated conditionally on
$\mathcal F_{k-1}$. For $m=1,\ldots,\bar{m}$ and $i\ge0$, define
\[
    \Delta_{k,i}:=\gamma^{-i}\Delta_k,
    \qquad
    X_{k,m}^{(i)}:=X_k+\Delta_{k,i}D_{k,m}.
\]
We also set $X_k^{(-1)}:=X_k$ and $\Delta_{k,-1}:=0$. The corresponding true and
estimated values are
\[
    f_{k,m}^{(i)}:=f(X_{k,m}^{(i)}),
    \qquad
    \tilde f_{k,m}^{(i)}:=\tilde f(X_{k,m}^{(i)}),
    \qquad i\ge0,
\]
and
\[
    f_k^{(-1)}:=f(X_k),
    \qquad
    \tilde f_k^{(-1)}:=\tilde f(X_k).
\]
At iteration $k$, the algorithm first samples $\bar{m}$ search directions independently, then it observes $\tilde f_k^{(-1)}$ and finally
tests the directions $D_{k,1},\ldots,D_{k,\bar{m}}$ sequentially. Direction $m$ is
successful at depth $i\ge0$ if
\[
    \tilde f_k^{(-1)}-\tilde f_{k,m}^{(i)}
    \ge
    \theta\Delta_{k,i}^{\Dexp}.
\]
The algorithm accepts the first direction for which the line search succeeds.
We denote its index by $m_k^*$ and the last successful extrapolation depth by
$H_k\ge0$. If no direction succeeds, we set $H_k=-1$. Thus, on a successful iteration,
\[
    D_k:=D_{k,m_k^*},
    \qquad
    X_k^{(i)}:=X_{k,m_k^*}^{(i)},
    \qquad
    f_k^{(i)}:=f_{k,m_k^*}^{(i)},
    \qquad
    \tilde f_k^{(i)}:=\tilde f_{k,m_k^*}^{(i)}.
\]
The accepted step length is
\[
    \Delta_k^{\rm acc}
    :=
    \begin{cases}
        0, & H_k=-1,\\[1mm]
        \gamma^{-H_k}\Delta_k, & H_k\ge0.
    \end{cases}
\]
Hence $X_{k+1}=X_k+\Delta_k^{\rm acc}D_k,$
with the convention that $X_{k+1}=X_k$ when $H_k=-1$. The stepsize update is
\[
\Delta_{k+1}
=
\begin{cases}
    \gamma\Delta_k, & H_k=-1,\\[1mm]
    \gamma^{-1}\Delta_k, & H_k=0,\\[1mm]
    \gamma^{-H_k}\Delta_k, & H_k\ge1.
\end{cases}
\] %\sara{so same stepsize when $H_k=0$ and $H_k=1$?}
Equivalently, on successful iterations,
$\Delta_{k+1}
    =
    \max\{\Delta_k^{\rm acc},\gamma^{-1}\Delta_k\}.$
Thus every unsuccessful iteration contracts the radius by $\gamma$, whereas
every successful iteration expands the next radius by at least $\gamma^{-1}$.
% \begin{algorithm}[h]
% 	\caption{DSE-DD}
% 	\label{alg:dse2}
% 	\begin{algorithmic}
% 		\par\vspace*{0.1cm}
% 		\item$\,\,\,$\hspace*{0.3truecm} \textbf{Input:} $x_0\in \mathbb{R}^n$, $\Delta_0 > 0$,  $\theta > 0, \gamma \in (0, 1)$
% 		\item$\,\,\,$\hspace*{0.30cm}  \textbf{For}  $k=0, 1,...$
% 			\item$\,\,\,$\hspace*{0.60cm} Select a descent direction $d_k$ in the unit sphere. 
%             \item Build a Positive Spanning Set $\mathbb{D}_k$ from $d_k$.
%             \textbf{For} $g \in \mathbb{D}_k$
% 		\item$\,\,\,$\hspace*{0.60cm} Compute $\delta_k, \Delta_{k + 1}$ with \textbf{StochasticLinesearch}($\Delta_k, x_k, d_k, \theta, \gamma$)  
% 		\item$\,\,\,$\hspace*{0.60cm} Set $x_{k+1}= x_k + \delta_k d_k$	\item$\,\,\,$\hspace*{0.30truecm} \textbf{End For} 
% 	\end{algorithmic}
% \end{algorithm}
% \begin{algorithm}[h]
%     \caption{DSE}
%     \label{alg:dse2}
%     \begin{algorithmic}[1]
%         \State \textbf{Input:} $x_0 \in \mathbb{R}^n$, $\delta_0 > 0$, $\theta > 0$, $\gamma \in (0, 1)$
%         \For{$k = 0, 1, 2, \dots$}
%             \State Sample a direction $d_k$ on the unit sphere \sara{just a single direction?}
%             % \State Build a Positive Spanning Set $\mathbb{D}_k$ from $d_k$
%             % \For{$g \in \mathbb{D}_k$}
%                 \State Compute $\delta_k$, $\delta_{k+1}$ with \textbf{StochasticLinesearch}{$(x_k, \delta_k, d_k, \theta, \gamma)$}
%                 \State Set $x_{k+1} = x_k + \delta_k d_k$
%             % \EndFor
%         \EndFor
%     \end{algorithmic}
% \end{algorithm}
\begin{algorithm}[H]
\caption{DSE}
\label{alg:dse2-cartis}
\begin{algorithmic}[1]
\State \textbf{Inputs:} $x_0\in\R^n$, $\delta_0>0$, $\theta>0$, $\gamma\in(0,1)$, $\bar{m}\in\N$, $\Imax \in \N$
\For{$k=0,1,2,\dots$}
  \State Sample $d_{k,1}, \ \dots, \ d_{k,\bar{m}}\in \mathbb S^{n-1}$ independently
  \State Observe the baseline estimate $\tilde f_k^{(-1)}\gets \tilde f(x_k)$
  \For{$m=1$ \textbf{to} $\bar{m}$}
    \State $(\delta^{\mathrm{acc}},\,\delta^{\mathrm{next}})
      \gets \textbf{StochasticLinesearch}\left(x_k,\delta_k,d_{k,m},\theta,\gamma, \f_k^{(-1)}, \Imax\right)$
    \If{$\delta^{\mathrm{acc}}>0$}
   
      \State \textbf{break}
    \EndIf
  \EndFor
     \State $d_k=d_{k,m}$
    \State $ \delta_{k+1}\gets \delta^{\mathrm{next}}$
  \State $x_{k+1}\gets x_k+\delta^{\mathrm{acc}}\,d_{k}$
\EndFor
\end{algorithmic}
\end{algorithm}

% \begin{algorithm}[h]
%     \caption{StochasticLinesearch}
%     \label{alg:lsep}
%     \begin{algorithmic}[1]
%         \State \textbf{Input:} $x \in \mathbb{R}^n$, $\delta > 0$, $d \in \mathbb{R}^n$, $\theta > 0$, $\gamma \in (0, 1), \f_k^{(-1)}$
%         \If{$\tilde{f}(x + \delta d) > \f_k^{(-1)} - \theta \delta^{\Dexp}$}
%             \State \textbf{Return} $(0, \gamma \delta)$
%         \EndIf
%          \While{$\tilde{f}(x + \delta d) \leq \f_k^{(-1)} - \theta \delta^{\Dexp}$}
%             \State Set $\delta = \delta / \gamma$
            
%         \EndWhile
%         \State \textbf{Return} $(\gamma \delta, \gamma\delta)$
%     \end{algorithmic}
% \end{algorithm}

% \begin{algorithm}[h]
%     \caption{StochasticLinesearch}
%     \label{alg:lsep}
%     \begin{algorithmic}[1]
%         \State \textbf{Input:} $x \in \mathbb{R}^n$, $\delta > 0$, $d \in \mathbb{R}^n$, $\theta > 0$, $\gamma \in (0,1)$, $\f_k^{(-1)}, \Imax$
%         \State Set $\delta_0 = \delta$

%         \If{$\tilde{f}(x+\delta_0 d) > \f_k^{(-1)} - \theta \delta_0^{\Dexp}$}
%             \State \textbf{Return} $(0,\gamma \delta_0)$
%         \EndIf
%         \State Set $h=0$
%         \State Set $\delta^{\mathrm{acc}} = \delta_0$
%         \State Set $\delta^{\mathrm{trial}} = \delta_0/\gamma$

%         \While{$h<\Imax$ \textbf{and} $\tilde{f}(x+\delta^{\mathrm{trial}} d) \leq \f_k^{(-1)} - \theta \delta^{\mathrm{trial}}^{\Dexp}$}
%             \State Set $h = h+1$
%             \State Set $\delta^{\mathrm{acc}} = \delta^{\mathrm{trial}}$
%             \State Set $\delta^{\mathrm{trial}} = \delta^{\mathrm{trial}}/\gamma$
%         \EndWhile

%         % \If{$\delta^{\mathrm{acc}} = \delta_0$}
%         \If{$h=0$}
%             \State \textbf{Return} $(\delta^{\mathrm{acc}},\delta_0/\gamma)$
%         \Else
%             \State \textbf{Return} $(\delta^{\mathrm{acc}},\delta^{\mathrm{acc}})$
%         \EndIf
%     \end{algorithmic}
% \end{algorithm}
% \newpage
% {\color{red}what about writing the algorithm this way?}

\begin{algorithm}[H]
\caption{StochasticLinesearch}
\label{alg:lsep}
\begin{algorithmic}[1]
\State \textbf{Input:} $x \in \mathbb{R}^n$, $\delta > 0$, $d \in \mathbb{R}^n$, $\theta > 0$, $\gamma \in (0,1)$, $\f_k^{(-1)}, \Imax$
        \State Set $\delta_0 = \delta$, $\delta_i=\gamma^{-i}\delta$, \quad $i=0,\ldots,\Imax$
\If{$\tilde f(x+\delta_0 d)>\tilde f_k^{(-1)}-\theta\delta_0^p$}
    \State \Return $(0,\gamma\delta_0)$
\EndIf
\State $h=\max\left\{j\in[0:\Imax]:
\tilde f(x+\delta_i d)\le \tilde f_k^{(-1)}-\theta\delta_i^p, \forall\ i\in[0:j] \right\}$
\If{$h=0$}
    \State \Return $(\delta_0,\delta_0/\gamma)$
\Else
    \State \Return $(\delta_h,\delta_h)$
\EndIf
\end{algorithmic}
\end{algorithm}

\noindent
We  would like to highlight that the maximization problem defined in Step 6 of the stochastic linesearch procedure is only a compact way of describing a sequential run-until-failure procedure.
In practice, after the level \(i=0\) test succeeds, the line search evaluates the
levels \(i=1,2,\ldots,\Imax\) one at a time and stops as soon as the
sufficient-decrease test fails. The returned value of \(h\) is the last
consecutive level for which all tests from \(0\) up to \(h\) have succeeded.
Thus, the algorithm does not evaluate extrapolation levels beyond the first
failed test.  The following remark finally clarifies the role of the maximum extrapolation depth
\(\Imax\) in Algorithm~\ref{alg:lsep}.


\begin{remark}[Role of the extrapolation cap]
The extrapolation cap \(\Imax\) is mainly an analytical device and is not
restrictive in practice. Indeed, any implementation is run under a finite
oracle-call budget, and \(\Imax\) can be chosen as large as this
budget allows. Its role in the analysis is to keep the extrapolation search
tree uniformly finite, which yields a simple maximal-error bound and a direct
tested-point complexity estimate (See Sections \ref{sec:assumptions}-\ref{sec:clarke complexity}).
\end{remark}




% \begin{algorithm}[h]
%     \caption{\textsc{StochasticLinesearchCapped}}
%     \label{alg:lsep-cartis}
%     \begin{algorithmic}[1]
%         \State \textbf{Input:} $x \in \mathbb{R}^n$, $\delta > 0$, $d \in \mathbb{R}^n$, $\theta > 0$, $\gamma \in (0, 1)$, $\delta_{\max}$
%         \If{$\tilde{f}(x + \delta d) > \tilde{f}(x) - \theta \delta^{\Dexp}$}
%             \State \textbf{Return} $(0, \gamma \delta)$
%         \EndIf
%         \While{$\tilde{f}(x + \delta d) < \tilde{f}(x) - \theta \delta^{\Dexp}$}
%         \If{$\delta = \delta_{\max}$}
%             \State $\delta = \gamma^{-1}\delta_{\max}$
%             \State \textbf{break}
%         \EndIf
%             \State Set $\delta = \min\{ \ \delta / \gamma, \, \delta_{\max} \ \}$ \sara{can do $\delta_{\max} / \gamma$ directly here and remove step 6-8}
%         \EndWhile
%         \State \textbf{Return} $(\gamma \delta, \gamma \delta)$
%     \end{algorithmic}
% \end{algorithm}

% initial set of assumptions
\section{Assumptions and Preliminaries}\label{sec:assumptions}
% We work on a filtered probability space $(\Omega,\mathcal F,\{\mathcal F_k\}_{k\ge 0},\mathbb P)$.
% At outer iteration $k$, let $X_k\in\mathbb R^n$ and $\Delta_k>0$ be the current point and stepsize.
% For each inner index $i\in\mathbb N_0\cup\{-1\}$ we write
% \[
% X_k^{(-1)}:=X_k,\qquad X_k^{(i)}:=X_k+\gamma^{-i}\Delta_k\,D_k,
% \qquad \Delta_{k}^{(i)}:=\gamma^{-i}\Delta_k, \quad 0<\gamma<1,
% \]
% and denote by $f_k^{(i)}:=f(X_k^{(i)})$ the true value and by $\tilde f_k^{(i)}$ its noisy estimate.
We now state the main assumptions on the objective and the stochastic oracle, deriving the key error bounds needed for the analysis. Throughout, we condition on $\F_{k-1}$, which contains $X_k, \D_k$, and the sampled directions $D_{k,1}, \ \dots, \ D_{k,\bar{m}}$, but not the stochastic estimates generated at iteration $k$. We assume noisy function estimates with the following basic structure.
%--------------------------------------------------------------------------
% Assumption: Compactness of the iterate (and trial) set
%--------------------------------------------------------------------------
\begin{assumption}[Compactness of the iterates]\label{ass:compact}
There exists a nonempty compact set \( {\mathcal C}\subset\mathbb R^n\) such that $X_k\in {\mathcal C} \ \text{a.s. for all $k\ge0$}.$
Moreover, the initial stepsize $\Delta_0>0$ is deterministic and finite.
% Moreover, for every sampled direction \(D_{k,m}\), \(m=1,\ldots,\bar{m}\), and every extrapolation
% level \(i\) actually tested by the line search, the corresponding trial point
% \[
% X_{k,m}^{(i)}
% :=
% X_k+\gamma^{-i}\Delta_kD_{k,m}
% \]
% also belongs to \(\mathcal C\).
\end{assumption}

\begin{remark}[Compactness of the tested points]
\label{rem:compact-tested-points}
Under the fixed maximum extrapolation depth
\(\Imax\in\mathbb N_0\),
% \sara{we are not assuming this, we are enforcing it}. Thus, at every iteration \(k\), every tested trial
% point has the form
% \[
%     X_{k,m}^{(i)}
%     =
%     X_k+\gamma^{-i}\Delta_kD_{k,m},
%     \qquad
%     1\le m\le \bar{m},
%     \qquad
%     0\le i\le \Imax.
% \]
and Assumption~\ref{ass:compact}, all tested points lie in a compact
enlargement of \({\mathcal C}\). Indeed, let
\[
    \Delta_{\mathcal C}:=\operatorname{diam}({\mathcal C}),
    \qquad
    \bar\Delta:=\max\{\Delta_0,\gamma^{-1}\Delta_{\mathcal C}\}.
\]
% \sara{shouldn't we also assume that $\Delta_0\le \gamma^{-1}\Delta_{\mathcal C}$ to stay in $\mathcal C$?} 
We first prove that $\Delta_k\le \bar\Delta,
    \ \forall k\ge0.$
The claim is true for \(k=0\) by the definition of \(\bar\Delta\). Suppose that
\(\Delta_k\le\bar\Delta\). If iteration \(k\) is unsuccessful, then
\[
    \Delta_{k+1}=\gamma\Delta_k\le\Delta_k\le\bar\Delta.
\]
If iteration \(k\) is successful with \(H_k=0\), then 
$  X_{k+1}=X_k+\Delta_kD_k$ is the accepted point.
Since \(X_k,X_{k+1}\in {\mathcal C}\), we have
\[
    \Delta_k=\|X_{k+1}-X_k\|\le \Delta_{\mathcal C}.
\]
Therefore,
\[
    \Delta_{k+1}=\gamma^{-1}\Delta_k\le\gamma^{-1}\Delta_{\mathcal C}\le\bar\Delta.
\]
Finally, if iteration \(k\) is successful with \(H_k\ge1\), then $X_{k+1}=X_k+\gamma^{-H_k}\Delta_kD_k,$ and hence
\[
    \Delta_{k+1}
    =
    \gamma^{-H_k}\Delta_k
    =
    \|X_{k+1}-X_k\|
    \le
    \Delta_{\mathcal C}
    \le
    \bar\Delta.
\]  
% \Ant{TO REMOVE: We cannot treat all successful cases through
% \[
%     \Delta_{k+1}=\|X_{k+1}-X_k\|,
% \]
% because when \(H_k=0\) the accepted displacement satisfies
% \[
%     \|X_{k+1}-X_k\|=\Delta_k,
% \]
% whereas the algorithm deliberately sets
% \[
%     \Delta_{k+1}=\gamma^{-1}\Delta_k>\Delta_k.
% \]}
Thus, by induction, $\Delta_k\le\bar\Delta, \ \forall k\ge0.$
Consequently, for every tested point,
\[
    \|X_{k,m}^{(i)}-X_k\|
    =
    \gamma^{-i}\Delta_k
    \le
    \gamma^{-\Imax}\bar\Delta.
\]
Therefore
\[
    X_{k,m}^{(i)}
    \in
    {{\mathcal C}_{enl}}
    :=
    {\mathcal C}+\gamma^{-\Imax}\bar\Delta\,\mathbb B,
    \qquad
    1\le m\le \bar{m},\quad 0\le i\le \Imax.
\]
Here,
\(
    \mathbb B:=\{z\in\mathbb R^n:\|z\|\le 1\}
\)
denotes the closed unit ball in \(\mathbb R^n\), and the symbol \(+\) denotes
the Minkowski sum of sets, namely
\(
    A+B:=\{a+b:a\in A,\ b\in B\}.
\)
Since \({\mathcal C}\) is compact and \(\gamma^{-\Imax}\bar\Delta\,\mathbb B\) is
compact, the set  ${\mathcal C}_{enl}$ is compact. Hence, all iterates and all tested trial
points generated by the algorithm lie in the compact set ${\mathcal C}_{enl}$.
\end{remark}

% \AnthonyNew{\\For now I'm using $|f|<f_{\max}$ since this is what I used for the proofs in the appendix. If we decide to stick with this boundedness I will move to $f_{\min} < f < f_{\max}$ and change the proofs accordingly}
% \begin{assumption}[Bounded function]\label{ass:compact}
% Assume that for any iterate $X_k$ the true function value is bounded \[|f(X_k)|\leq f_{\max}\] for every iteration $k$. 
% \end{assumption}
%--------------------------------------------------------------------------
% Assumption: Conditional Independence
%--------------------------------------------------------------------------
\begin{assumption}[Conditional independence]\label{ass:indep}
At iteration \(k\), conditionally on \(\mathcal F_{k-1}\), the stochastic
function estimates associated with the base point and with the search
tree,
\[
    \tilde f_k^{(-1)}
    \quad\text{and}\quad
    \left\{
        \tilde f_{k,m}^{(i)}
        :
        1\le m\le \bar{m},\ 0\le i\le \Imax
    \right\},
\]
are mutually independent.
\end{assumption}
%--------------------------------------------------------------------------
% Assumption: Quadratic tail bound
%--------------------------------------------------------------------------
\begin{assumption}[$\Dexp$-tail for single-point errors]\label{ass:quadratic tail}
There exists $\varepsilon_h>0$ such that, for every $\alpha\ge \varepsilon_h$:
\begin{equation}\label{eq:A4}
\begin{aligned}
&\mathbb P\!\left(
|\tilde f_{k}^{(-1)}-f_{k}^{(-1)}|
\ge
\alpha\,\Delta_k^{\Dexp}
\ \middle|\ 
\F_{k-1}
\right)
\le
\frac{\varepsilon_h}{\alpha^{\aexp}},
\\[1mm]
&\mathbb P\!\left(
|\tilde f_{k,m}^{(i)}-f_{k,m}^{(i)}|
\ge
\alpha\,\Delta_k^{\Dexp}
\ \middle|\ 
\F_{k-1}
\right)
\le
\frac{\varepsilon_h}{\alpha^{\aexp}},
\quad
i\in[0:\Imax],\ m\in[1:\bar{m}].
\end{aligned}
\end{equation}


\end{assumption}
%--------------------------------------------------------------------------
% Assumption: Post-selection L1 at the selected level
%--------------------------------------------------------------------------
% \begin{assumption}[Post-selection $L^1$ control for the difference estimator]\label{ass:postL1} %\label{lem:PS-L1}
% Define the difference estimator \[E_{k,i}:=(\tilde f_k^{(-1)}-\tilde f_k^{(i)})-(f_k^{(-1)}-f_k^{(i)}) = \bar E_{k,-1} - \bar E_{k,i}.\]
% Let $H_k$ be the (random) index selected by the inner routine at iteration $k$ (set $E_{k,-1}:=0$ if failure is encoded by $H_k=-1$). There exists a constant $c>0$ such that %the following bound
% \begin{equation}\tag{PS-$L^1$}\label{eq:postsel-diff}
% \mathbb{E}\!\left[\,|E_{k,H_k}|\,\middle|\,\mathcal F_{k-1}\right]\ \le\ c\,\Delta_k^{\Dexp}
% \qquad\text{for all $k$.}
% \end{equation}
% % with $c$ constant.
% \end{assumption}

% \begin{assumption}[Bounded maximal iteration error]
% \label{ass:postL1}
% At iteration \(k\), after sampling the trial directions
% \(D_{k,1},\ldots,D_{k,\bar{m}}\), define the trial points
% \[
% X_{k,m}^{(i)}
% :=
% X_k+\gamma^{-i}\Delta_kD_{k,m},
% \qquad
% 1\le m\le \bar{m},\qquad 0\le i\le \Imax.
% \]
% Using the notation
% \(
% \tilde f_k:=\tilde f_k(X_k),
% \;
% \tilde f_{k,m}^{(i)}
% :=
% \tilde f_k(X_{k,m}^{(i)}),
% \)
% define the maximal difference-estimation error over the potential capped search tree by
% \[
% E_k^{\max}
% :=
% \max_{1\le m\le \bar{m}}
% \max_{0\le i\le \Imax}
% \left|
% \left(
% \tilde f_k-\tilde f_{k,m}^{(i)}
% \right)
% -
% \left(
% f(X_k)-f(X_{k,m}^{(i)})
% \right)
% \right|.
% \]
% {\color {red} FR: No need to assume that. Thanks to $i\leq \Imax$ you have that c has a bound, right? Combine Ass. 3.3 and $i\leq \Imax$ you get following inequality satisfied and a bound for $c$} We assume that there exists a constant \(c>0\), independent of \(k\), such that
% \[
% \mathbb E\left[
% E_k^{\max}
% \mid
% \mathcal F_{k-1}
% \right]
% \le
% c\Delta_k^p
% \qquad
% \text{a.s. for all } k\ge 0.
% \]
% \end{assumption}

\noindent
To simplify the notation, from now on we set $\bar E_{k,m}^{(i)}:=\tilde f^{(i)}_{k,m} - f^{(i)}_{k,m}$, with
\(
    \bar E_k^{(-1)}
    :=
    \tilde f_k^{(-1)}-f_k^{(-1)}
\),
 and define the error difference estimator \[E_{k,m}^{(i)}:=(\tilde f_k^{(-1)}-\tilde f_k^{(i)})-(f_k^{(-1)}-f_k^{(i)}) = \bar E_{k}^{(-1)} - \bar E_{k,m}^{(i)},\]
%--------------------------------------------------------------------------
% Assumption: Post-selection L1 at the selected level
%--------------------------------------------------------------------------
Assumptions~\ref{ass:compact} and~\ref{ass:indep} are standard in stochastic derivative-free optimization (see, e.g., \cite{MR4238147,MR4758401}). 
In particular, compactness of the iterate set allows us to invoke tools from nonsmooth analysis
(such as the Clarke subdifferential) and ensures, under mild regularity (e.g., continuity on ${\mathcal C}_{enl}$), that
the objective $f$ attains its minimum. %\sara{why Assumption~\ref{ass:indep} needed for this?}. 
For notational simplicity, since our analysis will be restricted to a compact set, we will often write $f$ is $L$-Lipschitz.  Assumptions~\ref{ass:quadratic tail}, instead, is a technical condition that provides
quantitative control of the stochastic estimation error $\tilde f$ and
is needed to establish the merit-function drift bounds used throughout the analysis. From this basic set of hypotheses we derive several consequences that will be used in the convergence
and complexity proofs. For readability, we list the statements here and defer the proofs to the Appendix \ref{appendix:A}.
% From compactness we can immediately get several useful results regarding the behvaiour of the stepsize $\D_k$ and extrapolation index $H_k$. In particular, it is easy to show the followings.
% \begin{remark}[Compactness implies bounded $f$]\label{rem:C-bounded-and-Hk}
% Under Assumption \ref{ass:compact} and continuity, $f$ is bounded on $\mathcal{C}$:
% \[
% f_{\min}:=\inf_{x\in \mathcal{C}} f(x)>-\infty,\qquad f_{\max}:=\sup_{x\in \mathcal{C}} f(x)<\infty,
% \]
% so $f(X_k),f(X_k^{(i)})\in[f_{\min},f_{\max}]$ for all $k,i$.
% \end{remark}
% \begin{remark}[$H_k$ is finite a.s.]\label{rem: H_k is finite}
%     if \ref{ass:compact} holds, then the extrapolation index $H_k$ is finite almost surely, for every $k$.
% \end{remark}
Similarly, from Assumption \ref{ass:quadratic tail} we can get several other useful results by simply integrating the tails.
\begin{lemma}[Conditional Mean Bound]
\label{lem: conditional mean bound}
Let Assumption \ref{ass:quadratic tail} hold. For each $k$, $i\in[-1: \Imax]$ and $m\in[1:\bar{m}]$ recall the estimation error $\bar E_{k,m}^{(i)}:=\tilde f^{(i)}_{k,m} - f^{(i)}_{k,m}$. Then, there exists a constant $\mu_h>0$ such that the following bound on the conditional mean holds, uniformly in $k$ and $i$
\[
\mathbb{E}\!\left[\,|\bar E_{k,m}^{(i)}|\ \middle|\ \mathcal F_{k-1}\right]
\ \le\ \mu_h\,\Delta_k^{\Dexp},
\]
with $\mu_h$ a suitably chosen positive parameter.
% \begin{equation} \label{eq:mu_h}
% \mu_h\ :=\ \min\!\left\{\Dexp \varepsilon_h^{(\Dexp-1)/\Dexp}\,,\ \varepsilon_h \left[ 1+ (\Dexp-1)\varepsilon_h^{-1/(\Dexp-1)}\right]\,\right\}.
% \end{equation}
% In particular, for $i=-1$ we get a bound for the baseline estimate. \sara{remove since obvious?}
\end{lemma}
% Now, for $i\ge 0$ define the \emph{difference-estimator error}
% \sara{already defined in Assumption~\ref{ass:postL1}}\[
% E_{k,i}\ :=\ (\tilde f_k^{\saraedit{(-1)}}-\tilde f_k^{(i)})\;-\;(f_k^{(-1)}-f_k^{(i)}).
% \]
From the previous lemma we can also bound the conditional expectation of the difference estimator $E_{k,m}^{(i)}$.
\begin{lemma}[Conditional $L^1$ bound for the error difference estimator]
\label{lem:diff-L1-from-A44}
Under Assumption \ref{ass:quadratic tail}, 
% define
% \[
% E_{k,i}\ :=\ (\tilde f_k^{\saraedit{(-1)}}-\tilde f^{(i)}_k)\;-\;(f(X_k^{\saraedit{(-1)}})-f(X_k^{(i)}))
% \ =\ e_{k,-1}-e_{k,i}.
% \]
% Then, 
for every outer iteration $k$ and every inner index $i\in\{-1,0,1,\dots, \Imax\}$,
\begin{equation}\label{eq:diff-L1-core}
\mathbb E\!\left[\,|E_{k,m}^{(i)}|\,\middle|\,\mathcal F_{k-1}\right]
\ \le\ 2\,\mu_h\,\Delta_k^{\Dexp} \qquad \text{a.s.},
\end{equation}
with $\mu_h$ as in Lemma \ref{lem: conditional mean bound}. Consequently, the signed conditional mean is also controlled:
\begin{equation}\label{eq:diff-signed}
\big|\mathbb E\!\left[\,E_{k,m}^{(i)}\,\middle|\,\mathcal F_{k-1}\right]\big|
\ \le\ 2\,\mu_h\,\Delta_k^{\Dexp} \qquad \text{a.s.}.
\end{equation}
\end{lemma}
%%%%%%%%%%%%%%%%%%%%%%%%%%%%%%%%%%%%%%%%%%%%%%%%%%%%%%%%%%%%%%%%%%%%%%%%%%
\begin{lemma}[Bounded maximal iteration error]
\label{ass:postL1}
At iteration \(k\), after sampling the trial directions
\(D_{k,1},\ldots,D_{k,\bar{m}}\), 
% define the trial points \sara{already defined}
% \[
% X_{k,m}^{(i)}
% :=
% X_k+\gamma^{-i}\Delta_k D_{k,m},
% \qquad
% 1\le m\le \bar{m},\qquad 0\le i\le \Imax.
% \]
% Using the notation
% \(
% \tilde f_k:=\tilde f_k(X_k),
% \;
% \tilde f_{k,m}^{(i)}
% :=
% \tilde f_k(X_{k,m}^{(i)}),
% \)
define the maximal difference-estimation error over the potential search tree by
% \[
% E_k^{\max}
% :=
% \max_{1\le m\le \bar{m}}
% \max_{0\le i\le \Imax}
% \left|
% \left(
% \tilde f_k-\tilde f_{k,m}^{(i)}
% \right)
% -
% \left(
% f(X_k)-f(X_{k,m}^{(i)})
% \right)
% \right|.
% \]
\[
E_k^{\max}
:=
\max_{1\le m\le \bar{m}}
\max_{0\le i\le \Imax}
\left|
E_{k,m}^{(i)}
\right|.
\]
If Assumption \ref{ass:quadratic tail} holds, then via Lemma \ref{lem:diff-L1-from-A44}, one has
\[
\mathbb E\left[
E_k^{\max}
\mid
\mathcal F_{k-1}
\right]
\le
c_{\max}\Delta_k^p
\qquad
\text{a.s.},
\]
where \(c_{\max} = 2\mu_h\bar{m}(\Imax+1)\).
\end{lemma}
% \begin{remark}[Pre-selection control is insufficient at the selected level]
% \label{rem:pre_vs_selected}
% Lemma \ref{lem:diff-L1-from-A44} does \emph{not} imply a bound of the same order for the selected level:
% \[
% \mathbb E\!\left[\,|E_{k,H_k}|\,\middle|\,\mathcal F_{k-1}\right]
% \ \le\ c\,\Delta_k^{\Dexp} \ .
% \]
% In fact, since the selection event $\{H_k=i\}$ is data-dependent, one only has the selection-aware decomposition
% \begin{equation}
% \label{eq:selected_decomp}
% \mathbb E\!\left[\,|E_{k,H_k}|\,\middle|\,\mathcal F_{k-1}\right]
% =\sum_{i\ge 0}\mathbb P(H_k=i\,|\,\mathcal F_{k-1})\;
% \mathbb E\!\left[\,|E_{k,i}|\,\middle|\,\sigma(\mathcal F_{k-1},\ \cup\ \{H_k=i\})\right],
% \end{equation} where the $\sigma$-algebra generated by $\mathcal F_{k-1}$ and the event $\{H_k=i\}$ contains the ``post-selection'' information. Therefore, it 
% %and the conditional expectations on the right-hand side  
% cannot be controlled by \eqref{eq:diff-L1-core}. Assumption~\ref{ass:postL1} allows us to overcome such selection bias. To ease the notation, from now on we set $c=2\mu_h$ as well.
% \end{remark}
\noindent




\noindent
%%%%%%%%%%%%%%%%%%%%%%%%%%%%%%%%%%%%%%%%%%%%%%%%%%%%%%%%%%%%%%%%%%%%%%%%%
% Finally, from assumption \ref{ass:postL1} we have the following bound.
% \begin{lemma}[Post-selection $L^1$ control for the difference estimator]\label{lem:PS-L1}
% Let $E_{k,i}:=(\tilde f_k-\tilde f_k^{(i)})-(f_k^{(-1)}-f_k^{(i)})$ and let $H_k$ be the (random)
% index selected by the inner routine at iteration $k$ (set $E_{k,-1}:=0$ if failure is encoded by $H_k=-1$).
% Then, under assumption \ref{ass:postL1}, choosing $c=\mu_h$, we get the following bound
% \begin{equation}\tag{PS-$L^1$}\label{eq:postsel-diff}
% \mathbb{E}\!\left[\,|E_{k,H_k}|\,\middle|\,\mathcal F_{k-1}\right]\ \le\ 2\,\mu_h\,\Delta_k^{\Dexp}
% \qquad\text{for all $k$.}
% \end{equation}
% \end{lemma}
% \begin{proof}
% By the tower property and independence,
% \begin{align}
% \mathbb{E}\!\left[\,|E_{k,H_k}|\ \middle|\ \mathcal{F}_{k-1}\right]
% &= \mathbb{E}\!\left[\ \mathbb{E}\!\left[\,|E_{k,H_k}|\ \middle|\ \mathcal{F}_{k-1},\, H_k\right]\ \middle|\ \mathcal{F}_{k-1}\right] \\
% &= \mathbb{E}\!\left[\ \sum_{i\ge 0}\mathbf{1}_{\{H_k=i\}}\,
%       \mathbb{E}\!\left[\,|E_{k,H_k}|\ \middle|\ \mathcal{F}_{k-1},\, H_k=i\right]\ \middle|\ \mathcal{F}_{k-1}\right] \\
% &= \mathbb{E}\!\left[\ \sum_{i\ge 0}\mathbf{1}_{\{H_k=i\}}\,
%       \mathbb{E}\!\left[\,|E_{k,i}|\ \middle|\ \mathcal G_{k,i}\right]\ \middle|\ \mathcal{F}_{k-1}\right] \\
% &\le 2\mu_h\,\Delta_k^{\Dexp} \cdot \sum_{i\ge 0}\mathbb{P}(H_k=i\mid \mathcal{F}_{k-1}) \leq 2\mu_h\,\Delta_k^{\Dexp}.
% \end{align}
% \end{proof}
Finally, we note that the conditional \(\Dexp\)-tail bound in
Assumption~\ref{ass:quadratic tail} directly implies the commonly used notion of
\(\beta\)-probabilistically \(\varepsilon_f\)-accurate function estimates, as
made precise in the following remark.
\begin{remark}[From $\Dexp$-tails to $\beta$-probabilistic $\varepsilon_f$-accuracy]
\label{rem:A4-to-BPE}
Assumption~\ref{ass:quadratic tail} states that, for some $\varepsilon_h>0$ and all
$i\in\{-1,0,1,\dots, \Imax\}$,
\[
\mathbb P\!\left(|\tilde f_k^{(i)}-f_k^{(i)}|\ \ge\ \alpha\,\Delta_k^{\Dexp}\ \middle|\ \mathcal F_{k-1}\right)
\ \le\ \frac{\varepsilon_h}{\alpha^{\aexp}}\qquad\text{for all }\alpha\ge \varepsilon_h.
\]
Fix a target confidence $\beta\in(0,1)$. Choose
\[
\varepsilon_f\ \ge\ \max\!\left\{\varepsilon_h,\ \left(\frac{\varepsilon_h}{\,1-\beta\,}\right)^{\frac{\Dexp-1}{\Dexp}}\right\}.
\]
Then, setting $\alpha=\varepsilon_f$ in the tail bound gives
\[
\mathbb P\!\left(|\tilde f_k^{(i)}-f_k^{(i)}|\ \le\ \varepsilon_f\,\Delta_k^{\Dexp}\ \middle|\ \mathcal F_{k-1}\right)
\ \ge\ 1-\frac{\varepsilon_h}{\varepsilon_f^{\aexp}}\ \ge\ \beta,
\]
i.e., the usual $\beta$-probabilistically $\varepsilon_f$-accurate condition holds uniformly in $i$. If the tail in Assumption~\ref{ass:quadratic tail} holds for all $\alpha>0$, the simpler choice
$\varepsilon_f=(\varepsilon_h/(1-\beta))^{\frac{\Dexp-1}{\Dexp}}$ suffices.
\end{remark}
We conclude by observing that Assumption~\ref{ass:quadratic tail} is not restrictive in practice. Indeed, standard sampling strategies in stochastic derivative-free optimization often rely on averaging multiple independent realizations of the stochastic oracle; see, for example,~\cite{MR4238147,MR3800867,MR3880261}. Under the classical assumptions commonly imposed on such sample-average estimators, one obtains the stronger  conditions typically adopted in the stochastic derivative-free optimization literature, which in turn imply Assumption~\ref{ass:quadratic tail}. For completeness, Appendix~\ref{appendix:random estimates} establishes this connection for a standard sample-average estimator and derives the corresponding batch-size requirements.
% We conclude by observing that standard sampling strategies in stochastic derivative-free optimization often rely on averaging multiple independent realizations of the stochastic oracle; see, for example,~\cite{MR4238147,MR3800867,MR3880261}. The classical assumptions underlying these sample-average estimators are stronger than Assumption~\ref{ass:quadratic tail} and yield the tail bound required throughout our analysis. For completeness, Appendix~\ref{appendix:random estimates} establishes this connection and derives the corresponding batch-size requirements.
%We conclude by observing that stochastic estimates satisfying
%Assumption~\ref{ass:quadratic tail} can be constructed through sample averaging under standard unbiased finite-moment conditions on the underlying oracle noise, as is customary in the stochastic derivative-free optimization literature. To keep the main development focused on the abstract oracle assumptions and
%their algorithmic consequences, we defer the precise sampling conditions and batch-size requirements to Appendix~\ref{appendix:random estimates}.

%, and Appendix~\ref{appendix:random estimates} shows how stronger tail conditions on the noise model (e.g., sub-Gaussian decay) yield improved constants in the $\mathcal{O}(\cdot)$ bounds}.


% \paragraph{Applicability of reduced-sample sizing (Rinaldi-Vicente-Zeffiro \cite{MR4758401}).}
% We here indicate our sufficient-decrease exponent with $\Dexp\in (1, 2]$. Under the \emph{weak tail bound on the \emph{reduction estimate}} of \cite{MR4758401}, the per-iteration batch size needed to calibrate the acceptance test can be reduced as follows.  \cite[Theorem~2.4 and Corollary~2.5]{MR4758401}{\color{red}, when considering  finite $r$-th moment with $r = r(\Dexp) = \frac{\Dexp}{\Dexp - 1}$, } give
% \[
% m_k=\mathcal O(\Delta_k^{-2\Dexp}), 
% \qquad\text{e.g., with }\Dexp=1+\tfrac{\varepsilon}{2}:\;\; m_k=\mathcal O(\Delta_k^{-2-\varepsilon}).
% \]
% If CRN (common number generator) framework can be applied, e.g., by imposing uniformly Lipschitz sample paths (\cite[Proposition~2.6]{MR4758401}) or by a Gaussian process with exponentiated kernel (\cite[Proposition~2.7]{MR4758401}), the required sample size is reduced by a factor \(\Delta_k^{2}\), i.e.,
% \[
% m_k=\mathcal O(\Delta_k^{\,2-2\Dexp}), 
% \qquad\text{so that }\Dexp=1+\tfrac{\varepsilon}{2}\Rightarrow m_k=\mathcal O(\Delta_k^{-\varepsilon}),
% \]
% similarly to the results reported in \cite[Theorem~2.8 and Corollary~2.9]{MR4758401}. Finally, the weak reduction-tail condition is implied by classical finite-variance bounds used in StoMADS (\cite[Proposition~2.10]{MR4758401}) and by an extended \(\beta\)-probabilistic accuracy assumption (\cite[Proposition~2.12]{MR4758401}), providing drop-in compatibility with standard stochastic DFO hypotheses.

% \begin{lemma}[From single-point to reduction tails]
% \label{lem:single-to-reduction}
% Fix $\Dexp>1$ and $s=\Dexp/(\Dexp-1)$. Suppose that for all $\alpha>0$,
% \[
% \mathbb{P}\!\left(\,|\tilde f(x)-f(x)|\ \ge\ \alpha\,\Delta^\Dexp\ \middle|\ \mathcal F_{k-1}\right)
% \ \le\ \frac{\varepsilon_h}{\alpha^{\,s}}.
% \]
% If the two function estimates at two given points $x$ and $y$ are formed from independent batches, then the reduction error
% $R=(\tilde f(x)-\tilde f(y))-(f(x)-f(y))$ satisfies
% \[
% \mathbb{P}\!\left(\,|R|\ \ge\ \alpha\,\Delta^\Dexp\ \middle|\ \mathcal F_{k-1}\right)
% \ \le\ \frac{2^{\,s+1}\,\varepsilon_h}{\alpha^{\,s}}
% \qquad(\forall \alpha>0).
% \]
% Hence the weak reduction-tail condition holds with the same power $s$ and a rescaled constant.
% \end{lemma}}
%%%%%%%%%%%%%%%%%%%%%%%%%%%%%%%%%%%%%%%%%%%%%%%%%%%%%%%%%%%%%%%%%%
%%%%%%%%%%%%%%%%%%%%%%%%%%%%%%%%%%%%%%%%%%%%%%%%%%%%%%%%%%%%%%%%%%



% convergence analysis
\section{Main Convergence Result} \label{sec:convergence}

% \begin{lemma}[No infinite extrapolation and tail control]\label{lem:no-infinite-extrap}
% Recall $0<\gamma<1$, $\Delta_{k,i}:=\gamma^{-i}\Delta_k$, and define 
% \[H_k:=\sup\{i\ge 0:\ \tilde f_k-\tilde f^{(i)}_k \ge \theta\,\Delta_{k,i}^{\Dexp}\},\]
% namely, the last index at which the sufficient decrease test holds, with the convention $\sup\emptyset=-\infty$ and
% $H_k=\infty$ if the test holds for all $i$. Then, Assumptions \ref{ass:compact} (compactness) imply
% \[
% \mathbb P\!\left(H_k<\infty \,\middle|\, \mathcal F_{k-1}\right)=1.
% \]
% Moreover, assumptions \ref{ass:indep} (independence) and \ref{ass:quadratic tail} ($\Dexp$-tail) yield explicit tail bounds for $H_k$, hence control of the conditional mean \(\mathbb{E}[H_k\,|\,\mathcal F_{k-1}]\). The precise inequalities and constants are reported in appendix \ref{appendix:A}.
% \end{lemma}
Following the trust-region/direct-search literature (see,e.g., \cite{MR3506227,MR4758401}), we analyze progress through the \emph{merit function}
\begin{equation}\label{eq:phi-def}
  \Phi_k \;=\; f(X_k)-f_{min} \;+\; \eta\,\Delta_k^{\Dexp},
\end{equation}
where \(f_{\min}\) is a lower bound for \(f\) on the compact region containing the
iterates, and $\eta$ is a positive parameter. The term \(f(X_k)-f^\star\) measures objective decrease, while
\(\eta\,\Delta_k^{\Dexp}\) penalizes overly aggressive step-size growth during successful extrapolations. This choice ensures that the \emph{merit-function
drift}
\(\EE[\Phi_k-\Phi_{k+1}\mid\F_{k-1}]\) trades off true decrease against step-size changes in a way
that yields a uniform lower bound proportional to \(\Delta_k^{\Dexp}\).
%------------------------------------------------------------
% Notation for the merit-function drift argument
%------------------------------------------------------------
% For the drift analysis, let
% \[
%     \Phi_k := f(X_k)-f_{\min}+\eta\Delta_k^\Dexp,
%     \qquad \eta>0,
% \]
% where \(f_{\min}\) is a lower bound for \(f\) on the compact region containing the
% iterates.
At iteration \(k\), define
$H_k\in\{-1,0,\ldots,\Imax\},$ 
where \(H_k=-1\) denotes failure of the whole iteration, while \(H_k=i\ge0\)
denotes success at extrapolation depth \(i\). For \(i=0,\ldots,\Imax\), set
\[
    \Delta_{k,i}:=\gamma^{-i}\Delta_k,
    \qquad
    \pi_i:=\mathbb P(H_k=i\mid\mathcal F_{k-1}),
\]
and also
\[
    \pi_{-1}:=\mathbb P(H_k=-1\mid\mathcal F_{k-1}).
\]
Thus
\[
    \pi_{-1}+\sum_{i=0}^{\Imax}\pi_i=1.
\]
On a successful iteration, let \(m_k^*\) be the accepted direction index, so that
\[
    X_{k+1}
    =
    X_{k,m_k^*}^{(H_k)}
    =
    X_k+\Delta_{k,H_k}D_{k,m_k^*}.
\]
On a failed iteration, \(X_{k+1}=X_k\). The stepsize update is
\[
\Delta_{k+1}
=
\begin{cases}
    \gamma\Delta_k, & H_k=-1,\\[1mm]
    \gamma^{-1}\Delta_k, & H_k=0,\\[1mm]
    \gamma^{-i}\Delta_k, & H_k=i\ge1.
\end{cases}
\]
Recall the maximal difference-estimation error
\[
E_k^{\max}
:=
\max_{1\le m\le \bar{m}}
\max_{0\le i\le \Imax}
\left|
\left(
\tilde f_k^{(-1)}-\tilde f_{k,m}^{(i)}
\right)
-
\left(
f_k^{(-1)}-f_{k,m}^{(i)}
\right)
\right|.
\]

%------------------------------------------------------------
% Merit-function drift lemma
%------------------------------------------------------------
\begin{lemma}[Merit-function drift]
\label{lem:phi_drift_main}
Let Assumption~\ref{ass:quadratic tail} hold. Then, by
Lemma~\ref{ass:postL1}, there exists \(c_{\max}>0\) such that
\[
    \mathbb E\!\left[
        E_k^{\max}
        \,\middle|\,
        \mathcal F_{k-1}
    \right]
    \le
    c_{\max}\Delta_k^\Dexp
    \qquad \text{a.s.}
\]
Assume that
\begin{equation}
\label{ass:eta}
    \eta>
    \frac{c_{\max}}{1-\gamma^\Dexp},
\end{equation}
and
\begin{equation}
\label{ass:theta_eta_compatibility}
    \theta
    \ge
    \eta\max\{1,\gamma^{-\Dexp}-\gamma^\Dexp\}.
\end{equation}
Then
\begin{equation}
\label{eq:drift_uniform}
\mathbb E\!\left[
    \Phi_k-\Phi_{k+1}
    \,\middle|\,
    \mathcal F_{k-1}
\right]
\ge
\left(
    \eta(1-\gamma^\Dexp)-c_{\max}
\right)\Delta_k^\Dexp
>0
\qquad \textup{a.s.}
\end{equation}
\end{lemma}






% \begin{lemma}[Merit-function drift]\label{lem:phi_drift_main}
% Fix $0<\gamma<1$ and $\eta>0$, and define the merit function
% \[
%     \Phi_k := f(X_k)-f_{\min}+\eta\Delta_k^\Dexp .
% \]
% At iteration $k$, the algorithm sequentially tests up to $\bar{m}$ directions. Let
% $H_k=-1$ denote failure of the entire iteration, and let $H_k=i\ge0$ denote
% success at extrapolation depth $i$. Let
% \[
%     \pi_i:=\mathbb P(H_k=i\mid\mathcal F_{k-1}),
%     \qquad
%     \Delta_{k,i}:=\gamma^{-i}\Delta_k .
% \]
% Assume that, on $\{H_k=i\}$, the accepted point is
% \[
%     X_{k+1}=X_{k,m_k^*}^{(i)}
%     =
%     X_k+\Delta_{k,i}D_{k,m_k^*},
% \]
% and the estimate-based acceptance identity holds:
% \[
%     \tilde f_k-\tilde f_{k,m_k^*}^{(i)}
%     =
%     \theta\Delta_{k,i}^{\Dexp}
%     -
%     \rho_{k,i,m_k^*}\Delta_{k,i}^{\Dexp},
%     \qquad
%     \rho_{k,i,m_k^*}\le0 .
% \]
% Assume also that the stepsize update is
% \[
% \Delta_{k+1}
% =
% \begin{cases}
%     \gamma\Delta_k, & H_k=-1,\\[1mm]
%     \gamma^{-1}\Delta_k, & H_k=0,\\[1mm]
%     \gamma^{-i}\Delta_k, & H_k=i\ge1.
% \end{cases}
% \]
% Let Lemma~\ref{ass:postL1} hold, namely
% \[
%     \mathbb E[E_k^{\max}\mid\mathcal F_{k-1}]
%     \le
%     c_{\max}\Delta_k^\Dexp .
% \]
% Suppose that
% \begin{equation}\label{ass:eta}
%     \eta>\frac{c_{\max}}{1-\gamma^\Dexp}
%     =:\eta_{\rm req},
% \end{equation}
% and
% \begin{equation}\label{ass:theta_eta_compatibility}
%     \theta
%     \ge
%     \eta\max\{1,\gamma^{-\Dexp}-\gamma^\Dexp\} =: \theta_{\rm req}.
% \end{equation}
% Then
% \begin{equation}\label{eq:drift_uniform}
% \mathbb E\!\left[
%     \Phi_k-\Phi_{k+1}
%     \,\middle|\,
%     \mathcal F_{k-1}
% \right]
% \ge
% \left(
%     \eta(1-\gamma^\Dexp)-c_{\max}
% \right)\Delta_k^\Dexp
% >0
% \qquad \textit{a.s.}
% \end{equation}
% \end{lemma}

\begin{proof}
All expectations are taken conditionally on $\mathcal F_{k-1}$. We decompose
\[
    \Phi_k-\Phi_{k+1}
    =
    f(X_k)-f(X_{k+1})
    +
    \eta(\Delta_k^\Dexp-\Delta_{k+1}^\Dexp).
\]

\medskip
\noindent
\textbf{Step 1: True-reduction term.}
Fix $i\in[0:\Imax]$. Define
\[
E_{k,m_k^*}^{(i)}
:=
\bigl(\tilde f_k-\tilde f_{k,m_k^*}^{(i)}\bigr)
-
\bigl(f_k-f_{k,m_k^*}^{(i)}\bigr).
\]
Then 
\[
f_k-f_{k,m_k^*}^{(i)}
=
\bigl(\tilde f_k-\tilde f_{k,m_k^*}^{(i)}\bigr)
-
E_{k,m_k^*}^{(i)}.
\]
On $\{H_k=i\}$, the sufficient decrease holds, so
\[
\tilde f_k-\tilde f_{k,m_k^*}^{(i)}
\geq
\theta\Delta_{k,i}^{\Dexp}
.
\]
Therefore
\[
\begin{aligned}
&\mathbf 1_{\{H_k=i\}}
\bigl(f_k-f_{k,m_k^*}^{(i)}\bigr)
\geq
\mathbf 1_{\{H_k=i\}}
\Bigl(
    \theta\Delta_{k,i}^{\Dexp}
    -
    E_{k,m_k^*}^{(i)}
\Bigr).
\end{aligned}
\]
Taking conditional expectations and using
\[
    \Delta_{k,i}^{\Dexp}
    =
    \gamma^{-i\Dexp }\Delta_k^\Dexp,
    \qquad
    \mathbb E[\mathbf 1_{\{H_k=i\}}\mid\mathcal F_{k-1}]
    =
    \pi_i,
\]
we obtain
\[
\begin{aligned}
\mathbb E\!\left[
    \mathbf 1_{\{H_k=i\}}
    \bigl(f_k-f_{k,m_k^*}^{(i)}\bigr)
    \,\middle|\,
    \mathcal F_{k-1}
\right]
\geq
\pi_i\theta\Delta_{k,i}^{\Dexp}
% -
% \mathbb E\!\left[
%     \mathbf 1_{\{H_k=i\}}
%     \rho_{k,i,m_k^*}\Delta_{k,i}^{\Dexp}
%     \,\middle|\,
%     \mathcal F_{k-1}
% \right]
-
\mathbb E\!\left[
    \mathbf 1_{\{H_k=i\}}E_{k,m_k^*}^{(i)}
    \,\middle|\,
    \mathcal F_{k-1}
\right].
\end{aligned}
\]
Since $|E_{k,m_k^*}^{(i)}|\le E_k^{\max}$, it follows that 
\[
\begin{aligned}
&\mathbb E\!\left[
    \mathbf 1_{\{H_k=i\}}
    \bigl(f_k-f_{k,m_k^*}^{(i)}\bigr)
    \,\middle|\,
    \mathcal F_{k-1}
\right]
\ge
\pi_i\theta\Delta_{k,i}^{\Dexp}
-
\mathbb E\!\left[
    \mathbf 1_{\{H_k=i\}}E_k^{\max}
    \,\middle|\,
    \mathcal F_{k-1}
\right].
\end{aligned}
\]
Summing over $i\in[0:\Imax]$ gives
\[
\begin{aligned}
&\sum_{i=0}^{\Imax}
\mathbb E\!\left[
    \mathbf 1_{\{H_k=i\}}
    \bigl(f_k-f_{k,m_k^*}^{(i)}\bigr)
    \,\middle|\,
    \mathcal F_{k-1}
\right]
\ge
\sum_{i=0}^{\Imax}\pi_i\theta\Delta_{k,i}^{\Dexp}
-
\mathbb E\!\left[
    \mathbf 1_{\{H_k\ge0\}}E_k^{\max}
    \,\middle|\,
    \mathcal F_{k-1}
\right].
\end{aligned}
\]
By Lemma~\ref{ass:postL1},
\[
\mathbb E\!\left[
    \mathbf 1_{\{H_k\ge0\}}E_k^{\max}
    \,\middle|\,
    \mathcal F_{k-1}
\right]
\le
\mathbb E[E_k^{\max}\mid\mathcal F_{k-1}]
\le
c_{\max}\Delta_k^\Dexp.
\]
Hence
\begin{equation}\label{eq:true_red_explicit_2}
\sum_{i=0}^{\Imax}
\mathbb E\!\left[
    \mathbf 1_{\{H_k=i\}}
    \bigl(f_k-f_{k,m_k^*}^{(i)}\bigr)
    \,\middle|\,
    \mathcal F_{k-1}
\right]
\ge
\left(
    \sum_{i=0}^{\Imax}\pi_i\theta\gamma^{-\Dexp i}
    -
    c_{\max}
\right)
\Delta_k^\Dexp .
\end{equation}

\medskip
\noindent
\textbf{Step 2: Stepsize term.}
By the modified stepsize update,
\[
\Delta_{k+1}
=
\mathbf 1_{\{H_k=-1\}}(\gamma\Delta_k)
+
\mathbf 1_{\{H_k=0\}}(\gamma^{-1}\Delta_k)
+
\sum_{i=1}^{\Imax}
\mathbf 1_{\{H_k=i\}}(\gamma^{-i}\Delta_k).
\]
Therefore
% \[
% \begin{aligned}
% \eta(\Delta_k^\Dexp-\Delta_{k+1}^\Dexp)
% =
% \eta\Delta_k^\Dexp
% \Bigg[
%     &\mathbf 1_{\{H_k=-1\}}(1-\gamma^\Dexp)
%     +
%     \mathbf 1_{\{H_k=0\}}(1-\gamma^{-\Dexp})
% +
%     \sum_{i=1}^{\Imax}
%     \mathbf 1_{\{H_k=i\}}
%     \bigl(1-\gamma^{-\Dexp i}\bigr)
% \Bigg].
% \end{aligned}
% \]
\[
\begin{aligned}
\eta(\Delta_k^\Dexp-\Delta_{k+1}^\Dexp)
={}&
\eta\Delta_k^\Dexp
\Bigg[
    \mathbf 1_{\{H_k=-1\}}(1-\gamma^\Dexp)
    +
    \mathbf 1_{\{H_k=0\}}(1-\gamma^{-\Dexp})+
\\
&\qquad\qquad
    +
    \sum_{i=1}^{\Imax}
    \mathbf 1_{\{H_k=i\}}
    \bigl(1-\gamma^{-\Dexp i}\bigr)
\Bigg].
\end{aligned}
\]
Taking conditional expectations yields
% \begin{equation}\label{eq:stepsize}
% \begin{aligned}
% &\mathbb E\!\left[
%     \eta(\Delta_k^\Dexp-\Delta_{k+1}^\Dexp)
%     \,\middle|\,
%     \mathcal F_{k-1}
% \right]
% =
% \eta\Delta_k^\Dexp
% \left[
%     \pi_{-1}(1-\gamma^\Dexp)
%     +
%     \pi_0(1-\gamma^{-\Dexp})
%     +
%     \sum_{i=1}^{\Imax}\pi_i(1-\gamma^{-\Dexp i})
% \right].
% \end{aligned}
% \end{equation}
\begin{equation}\label{eq:stepsize}
\begin{aligned}
\mathbb E\!\left[
    \eta(\Delta_k^\Dexp-\Delta_{k+1}^\Dexp)
    \,\middle|\,
    \mathcal F_{k-1}
\right]
={}&
\eta\Delta_k^\Dexp
\Bigg[
    \pi_{-1}(1-\gamma^\Dexp)
    +
    \pi_0(1-\gamma^{-\Dexp})+
\\
&\qquad\qquad
    +
    \sum_{i=1}^{\Imax}
    \pi_i(1-\gamma^{-\Dexp i})
\Bigg].
\end{aligned}
\end{equation}

\medskip
\noindent
\textbf{Step 3: Combining the two estimates.}
Combining~\eqref{eq:true_red_explicit_2} and~\eqref{eq:stepsize}, and recalling that the
true-reduction term is zero on $\{H_k=-1\}$, we obtain
% \begin{equation}\label{eq:drift_full}
% \begin{aligned}
% \mathbb E\!\left[
%     \Phi_k-\Phi_{k+1}
%     \,\middle|\,
%     \mathcal F_{k-1}
% \right]
% \ge
% \Bigg[
%     \pi_{-1}\eta(1-\gamma^\Dexp)
%     +
%     \pi_0\Bigl(
%         \theta+\eta(1-\gamma^{-\Dexp})
%     \Bigr)
%     +
%     \sum_{i=1}^{\Imax}
%     \pi_i
%     \Bigl(
%         \theta\gamma^{-\Dexp i}
%         +
%         \eta(1-\gamma^{-\Dexp i})
%     \Bigr)
%     -
%     c_{\max}
% \Bigg]\Delta_k^\Dexp .
% \end{aligned}
% \end{equation}

\begin{equation}\label{eq:drift_full}
\begin{aligned}
\mathbb E\!\left[\Phi_k-\Phi_{k+1}\,\middle|\,\mathcal F_{k-1}\right]
&\ge \Bigg[\pi_{-1}\eta(1-\gamma^\Dexp)
+\pi_0\Bigl(\theta+\eta(1-\gamma^{-\Dexp})\Bigr)+ \\
&\qquad
+\sum_{i=1}^{\Imax}\pi_i
\Bigl(\theta\gamma^{-\Dexp i}
+\eta(1-\gamma^{-\Dexp i})\Bigr)
-c_{\max}\Bigg]\Delta_k^\Dexp .
\end{aligned}
\end{equation}
Define
\[
    A:=\eta(1-\gamma^\Dexp), \quad 
    B:=\theta+\eta(1-\gamma^{-\Dexp})
      =\theta-\eta(\gamma^{-\Dexp}-1),
\]
and, for $i\ge1$,
\[
    C_i
    :=
    \theta\gamma^{-\Dexp i}
    +
    \eta(1-\gamma^{-\Dexp i})
    =
    \eta+(\theta-\eta)\gamma^{-\Dexp i}.
\]
Then, under~\eqref{ass:theta_eta_compatibility},
\[
    B\ge A,
    \qquad
    C_i\ge A\quad \forall i\ge1.
\]
\noindent
Indeed, $B\ge A$ is equivalent to
\[
    \theta-\eta(\gamma^{-\Dexp}-1)
    \ge
    \eta(1-\gamma^\Dexp),
\]
that is,
\[
    \theta\ge \eta(\gamma^{-\Dexp}-\gamma^\Dexp),
\]
which follows from~\eqref{ass:theta_eta_compatibility}. Moreover,
\eqref{ass:theta_eta_compatibility} also implies $\theta\ge\eta$. Hence, for $i\ge1$,
\[
    C_i
    =
    \eta+(\theta-\eta)\gamma^{-\Dexp i}
    \ge
    \eta
    \ge
    \eta(1-\gamma^\Dexp)
    =
    A.
\]
Therefore
\[
    \pi_{-1}A+\pi_0B+\sum_{i=1}^{\Imax}\pi_iC_i
    \ge
    \left(\pi_{-1}+\pi_0+\sum_{i=1}^{\Imax}\pi_i\right)A
    =
    A.
\]
Using this in~\eqref{eq:drift_full}, we obtain
\[
\mathbb E\!\left[
    \Phi_k-\Phi_{k+1}
    \,\middle|\,
    \mathcal F_{k-1}
\right]
\ge
\left(
    A-c_{\max}
\right)\Delta_k^\Dexp.
\]
Substituting $A=\eta(1-\gamma^\Dexp)$ gives
\[
\mathbb E\!\left[
    \Phi_k-\Phi_{k+1}
    \,\middle|\,
    \mathcal F_{k-1}
\right]
\ge
\left(
    \eta(1-\gamma^\Dexp)-c_{\max}
\right)\Delta_k^\Dexp.
\]
By~\eqref{ass:eta}, the coefficient is strictly positive. This concludes the proof.
\end{proof}






\begin{lemma}\label{l:radius0}
Under the hypotheses of Lemma~\ref{lem:phi_drift_main},
% hold. In particular, assume that there exists a constant
% \[
%     \Theta := \eta(1-\gamma^{\Dexp})-c_{\max}>0
% \]
% such that, for all $k\in\mathbb N_0$,
% \[
%     \mathbb E\!\left[
%         \Phi_k-\Phi_{k+1}
%         \,\middle|\,
%         \mathcal F_{k-1}
%     \right]
%     \ge
%     \Theta \Delta_k^{\Dexp}
%     \qquad \text{a.s.}
% \] \sara{why assume? didn't we just prove this above?}
it holds
\[
    \sum_{k=0}^{\infty}\Delta_k^{\Dexp}<\infty
    \qquad \text{a.s.}
\]
In particular,
\[
    \Delta_k\to0
    \qquad \text{a.s. as } k\to\infty .
\]
\end{lemma}






\begin{proof}
Define \(\Theta := \eta(1-\gamma^{\Dexp})-c_{\max}, \;\Theta>0\). Taking total expectations in the drift inequality gives
\[
    \mathbb E[\Phi_k-\Phi_{k+1}]
    \ge
    \Theta\,\mathbb E[\Delta_k^{\Dexp}]
    \qquad \forall k\in\mathbb N_0 .
\]
Summing from $k=0$ to $N$ yields
\[
    \Theta\sum_{k=0}^{N}\mathbb E[\Delta_k^{\Dexp}]
    \le
    \sum_{k=0}^{N}\mathbb E[\Phi_k-\Phi_{k+1}]
    =
    \mathbb E[\Phi_0-\Phi_{N+1}] .
\]
Since
\[
    \Phi_{N+1}
    =
    f_{N+1}-f_{\min}
    +
    \eta\Delta_{N+1}^{\Dexp}
    \ge 0,
\]
we obtain
\[
    \Theta\sum_{k=0}^{N}\mathbb E[\Delta_k^{\Dexp}]
    \le
    \mathbb E[\Phi_0].
\]
Therefore,
\[
    \sum_{k=0}^{N}\mathbb E[\Delta_k^{\Dexp}]
    \le
    \frac{\mathbb E[\Phi_0]}{\Theta}
    \qquad \forall N\in\mathbb N_0 .
\]
Letting $N\to\infty$ and using monotone convergence,
\[
    \mathbb E\!\left[
        \sum_{k=0}^{\infty}\Delta_k^{\Dexp}
    \right]
    =
    \sum_{k=0}^{\infty}\mathbb E[\Delta_k^{\Dexp}]
    \le
    \frac{\mathbb E[\Phi_0]}{\Theta}
    <
    \infty .
\]
Since the random variable
\(
    \sum_{k=0}^{\infty}\Delta_k^{\Dexp}
\)
is nonnegative and has finite expectation, it is finite almost surely. Hence
\(
    \sum_{k=0}^{\infty}\Delta_k^{\Dexp}<\infty
\)
almost surely.
% \sara{no, if $E[X]<\infty$ it doesn't imply $X<\infty$ a.s.}. \Ant{TO REMOVE: It does. Let
% \[
% S:=\sum_{k=0}^{\infty}\Delta_k^{\Dexp}\ge0.
% \]
% For a nonnegative random variable,
% \[
% \mathbb E[S]
% =
% \int_0^\infty \mathbb P(S>t)\,dt.
% \]
% If \(\mathbb P(S=\infty)>0\), then
% \[
% \mathbb P(S>t)\ge \mathbb P(S=\infty)>0
% \qquad\text{for every }t\ge0,
% \]
% and hence \(\mathbb E[S]=\infty\), a contradiction.\\
% In general, let \(S:\Omega\to[0,+\infty]\) be a nonnegative random variable such that
% \[
% \mathbb{E}[S]<\infty.
% \]
% For every \(n\in\mathbb{N}\), Markov's inequality gives
% \[
% \mathbb{P}(S>n)
% \leq
% \frac{\mathbb{E}[S]}{n}.
% \]
% Moreover,
% \[
% \{S=+\infty\}
% =
% \bigcap_{n=1}^{\infty}\{S>n\}.
% \]
% Since the events \(\{S>n\}\) form a decreasing sequence, continuity of
% probability from above yields
% \[
% \mathbb{P}(S=+\infty)
% =
% \lim_{n\to\infty}\mathbb{P}(S>n).
% \]
% Using Markov's inequality,
% \[
% 0
% \leq
% \mathbb{P}(S=+\infty)
% \leq
% \lim_{n\to\infty}\frac{\mathbb{E}[S]}{n}
% =
% 0.
% \]
% Therefore,
% \[
% \mathbb{P}(S=+\infty)=0,
% \]
% and hence
% \[
% \mathbb{P}(S<+\infty)=1.
% \]
% Thus, \(S\) is finite almost surely.\\
% Now let \(X \sim \mathrm{Exp}(\lambda)\), with \(\lambda > 0\). For every \(t \geq 0\),
% \[
% \mathbb{P}(X > t) = e^{-\lambda t}.
% \]
% As before,
% \[
% \{X = +\infty\}
% =
% \bigcap_{n=1}^{\infty} \{X > n\}.
% \]
% Since the sequence of events \(\{X > n\}\) is decreasing, continuity of probability from above yields
% \[
% \mathbb{P}(X = +\infty)
% =
% \lim_{n \to \infty} \mathbb{P}(X > n)
% =
% \lim_{n \to \infty} e^{-\lambda n}
% =
% 0.
% \]
% Therefore,
% \[
% \mathbb{P}(X < +\infty) = 1,
% \]
% and hence \(X\) is finite almost surely.
% }\\
Finally, since $\Delta_k^{\Dexp}\ge0$ and $p>1$, summability implies
\(
    \Delta_k^{\Dexp}\to0
    \ \ \text{a.s.}
\)
and therefore
\(
    \Delta_k\to0
    \ \ \text{a.s.}
\)
because $\Dexp>0$.
\end{proof}
\noindent
% OLD PROOF. Replaced with a cleaner one
% \begin{lemma} \label{l:radius0}
% Under Assumptions \ref{ass:indep}, \ref{ass:quadratic tail}, \ref{ass:postL1} and the bound \eqref{ass:eta}, then almost surely
% 	\begin{equation}
% 		\sum_{k \in \mathbb{N}_0} \Delta_k^{\Dexp} < + \infty \, .
% 	\end{equation}
% In particular, \(\Delta_k \to 0\) almost surely as \(k \to \infty\).
% \end{lemma}
% \begin{proof}
% % From Lemma \ref{lem:phi_drift_main}, by telescoping the identity and recalling the definition of expected value, it follows in a straightforward way (see, e.g., \cite[Theorem 3.1]{MR4758401})
% % \begin{equation}
% % 	\sum_{k \in \mathbb{N}_0} \tdelta_k^{\Dexp} < \infty \, ,
% % \end{equation}
% % Alternatively, for any $\bar{m}>0$, Markov's inequality gives
% % \[
% % \mathbb P( \sum_{k \in \mathbb{N}_0} \tdelta_k^{\Dexp}\ge \bar{m})\ \le\ \frac{\mathbb E[\sum_{k \in \mathbb{N}_0} \tdelta_k^{\Dexp}]}{\bar{m}}.
% % \]
% % Letting $\bar{m}\to\infty$ yields $\mathbb P(\sum_{k \in \mathbb{N}_0} \tdelta_k^{\Dexp}=\infty)\le \lim_{\bar{m}\to\infty}\mathbb P(\sum_{k \in \mathbb{N}_0} \tdelta_k^{\Dexp}\ge \bar{m})=0$, hence $\sum_{k \in \mathbb{N}_0} \tdelta_k^{\Dexp}<\infty$ a.s.
% Fix $\varepsilon>0$. By Markov inequality and monotone convergence theorem,
% \[
% \sum_{k=0}^{\infty} \mathbb{P}(\Delta_k^{\Dexp} \ge \varepsilon)
% \;\le\; \frac{1}{\varepsilon}\,\mathbb{E}\!\left[\sum_{k=0}^{\infty} \Delta_k^{\Dexp}\right]
% \;<\; \infty.
% \]
% Borel--Cantelli gives $\mathbb{P}(\Delta_k^{\Dexp} \ge \varepsilon \ \text{i.o.})=0$. Letting $\varepsilon_m\downarrow 0$
% (e.g., $\varepsilon_m=1/m$) yields $\Delta_k^{\Dexp}\to 0$ a.s., hence $\Delta_k\to 0$ a.s. Moreover, for every $k \in \mathbb{N}_0$ we have
% 	\begin{equation}
% 		\sum_{i \in I_k} (\Delta_k^{(i)})^{\Dexp} = \sum_{i \in I_k} \gamma^{\Dexp(H_k - i)} \Delta_{k + 1}^{\Dexp} \leq \sum_{j =0}^{+\infty}\gamma^{\Dexp j} \Delta_{k + 1}^{\Dexp} \leq \frac{1}{1 - \gamma^{\Dexp}}  \Delta_{k + 1}^{\Dexp}  \, .
% 	\end{equation}
% Thus it also holds
% 	\begin{equation}
% 	\sum_{k \in \mathbb{N}_0} \sum_{i \in I_k} (\Delta_k^{(i)})^{\Dexp} \leq  \frac{1}{1 - \gamma^{\Dexp}}\sum_{k \in \mathbb{N}} \tdelta_k^{\Dexp} < +\infty \,.
% \end{equation}
% \end{proof}
Thus, from the merit-function drift we have
\(\sum_{k=0}^\infty \Delta_k^{\Dexp}<\infty\) a.s., hence \(\Delta_k\to0\) along the sequence almost surely. This in turn guarantees the existence of a subsequence of unsuccessful iterations.
\begin{lemma}[Existence of infinitely many unsuccessful iterations] \label{lem:inf-unsuccessful}
Let
\[
    \mathcal U
    :=
    \{k\in\mathbb N_0 : \text{iteration } k \text{ is unsuccessful}\}
\]
be the random set of indices of unsuccessful iterations. Assume that
\[
    \sum_{k=0}^{\infty}\Delta_k^{\Dexp}<\infty
    \qquad \text{a.s.}
\]
Then, with probability one, the set $\mathcal U$ is infinite.
\end{lemma}

\begin{proof}
Argue by contradiction. Suppose that, with strictly positive probability,
there are only finitely many unsuccessful iterations. On this event, there
exists a finite random index $k^u$ such that every iteration $k\ge k^u$ is
successful. By the stepsize update rule, on a successful iteration we have
\[
\Delta_{k+1}
=
\begin{cases}
    \gamma^{-1}\Delta_k, & H_k=0,\\[1mm]
    \gamma^{-H_k}\Delta_k, & H_k\ge1.
\end{cases}
\]
Since $0<\gamma<1$, it follows in both cases that
\[
    \Delta_{k+1}\ge \gamma^{-1}\Delta_k>\Delta_k .
\]
In particular, on the event under consideration, the tail sequence
$\{\Delta_k\}_{k\ge k^u}$ is strictly increasing. Since $\Delta_{k^u}>0$, we get
\[
    \Delta_k\ge \Delta_{k^u}>0
    \qquad \forall k\ge k^u .
\]
Therefore
\[
    \sum_{k=k^u}^{\infty}\Delta_k^{\Dexp}
    \ge
    \sum_{k=k^u}^{\infty}\Delta_{k^u}^{\Dexp}
    =
    \infty,
\]
which contradicts the assumption
\[
    \sum_{k=0}^{\infty}\Delta_k^{\Dexp}<\infty
    \qquad \text{a.s.}
\]
Hence the probability that $\mathcal U$ is finite must be zero. Therefore
$\mathcal U$ is infinite almost surely.
\end{proof}
\noindent
% OLD PROOF
% \begin{lemma}[Existence of infinitely many unsuccessful iterations]
% \label{lem:inf-unsuccessful}
% Let $\mathcal{U} := \{k \in \mathbb{N}_0 : \text{iteration } k \text{ is unsuccessful}\}$ be the random set of indices of unsuccessful iterations. Assume that $\sum_{k=0}^\infty \Delta_k^{\Dexp} < \infty$ almost surely. Then, with probability one, the set $\mathcal{U}$ is infinite.
% \end{lemma}
% \begin{proof}
% By contradiction, suppose that with strictly positive probability there are only finitely many unsuccessful iterations. On this event, there exists a finite index $k^u$ such that for all $k \ge k^u$, iteration $k$ is successful.
% By the algorithmic update rule, on successful iterations we have $\Delta_{k+1} = \gamma^{-H_k}\Delta_k$. Since $\gamma \in (0,1)$ and the extrapolation depth $H_k \ge 0$, this implies $\Delta_{k+1} \ge \Delta_k$. 
% Thus, the tail sequence $\{\Delta_k\}_{k \ge k^u}$ is monotonically nondecreasing. Because $\Delta_{k^u} > 0$ strictly, it follows that $\Delta_k \ge \Delta_{k^u} > 0$ for all $k \ge k^u$.
% This implies that the infinite tail of the series diverges:
% $$
% \sum_{k=k^u}^\infty \Delta_k^{\Dexp} = \infty,
% $$
% which directly contradicts the hypothesis that $\sum_{k=0}^\infty \Delta_k^{\Dexp} < \infty$ almost surely. Therefore, the probability of having only finitely many unsuccessful iterations must be zero, concluding that $\mathcal{U}$ is infinite almost surely.
% \end{proof}
It is on these unsuccessful iterations that our convergence analysis will concentrate.\\
We say that a subsequence $\{x_k\}_{k \in L}$ is refining if it converges to $x^*$ and if $\{d_k\}_{k \in L}$ is dense in the unit sphere. We will show that under \eqref{ass:eta}, \eqref{ass:theta_eta_compatibility}, if Assumptions  \ref{ass:compact}, \ref{ass:indep}, \ref{ass:quadratic tail} hold, then with probability 1 all refining subsequences converge to a Clarke stationary point. We start by recalling the definition of Clarke generalized directional derivative  and Clarke subdifferential.
\begin{definition}[Clarke generalized directional derivative]
Let $f:\mathbb{R}^n\to\mathbb{R}$ be Lipschitz near $x$.  Its \emph{Clarke generalized directional derivative} at $x$ in the direction $h$ is
\[
  f^\circ(x;h)
  :=
  \limsup_{\substack{y\to x\\t\downarrow0}}
    \frac{f(y + t\,h) - f(y)}{t}.
\]
\end{definition}
\begin{definition}[Clarke subdifferential]
Under the same hypotheses, the \emph{Clarke subdifferential} of $f$ at $x$ is
\[
  \partial_C f(x)
  := \bigl\{v\in\mathbb{R}^n : f^\circ(x;h)\;\ge\;\langle v,\,h\rangle
    \quad\forall\,h\in\mathbb{R}^n\bigr\}.
\]
Equivalently, one shows
\[
  \partial_C f(x)
  = \mathrm{conv}\Bigl\{
      \lim_{i\to\infty}\nabla f(x_i):
      x_i\to x,\ f\text{ is differentiable at }x_i
    \Bigr\}.
\]
\end{definition}
\begin{definition}[Clarke-stationary point]\label{def:clarke-stationary}
Let $f:\mathbb{R}^n\to\mathbb{R}$ be locally Lipschitz and let $\partial_{C} f(x)$ denote its
Clarke subdifferential. A point $x^\star\in\mathbb{R}^n$ is \emph{Clarke-stationary} if and only if $0 \in \partial_{C} f(x^\star).$
Equivalently, the Clarke directional derivative satisfies
\[
f^{\circ}(x^\star; d)\ \ge\ 0 \qquad \text{for all } d\in\mathbb{R}^n.
\]
\end{definition}
Now some preliminary results. The next lemma extends Assumption \ref{ass:quadratic tail} to the case where $\alpha$ is a random variable.
%  OLD VERSION: too strong, since by assumption 3.3 we have the bound only for alpha >= eps_h
% \begin{lemma} \label{l:rv}
% Let $A$ be a nonnegative $\F_{k - 1}$ measurable random variable. If  Assumption \ref{ass:quadratic tail} holds, then
% \begin{equation}
%     \Pb\left(|\f_{k}^{(i)} - f_k^{(i)}| \geq A \Delta_{k}^{\Dexp} \ | \ \F_{k - 1} \right) \leq \eh(A)
% \end{equation}
% with $\eh(A) \defeq +\infty \ind_{\{A=0\}} + \frac{\eh}{A^{\aexp}}\ind_{\{A>0\}}.$
% \end{lemma}
\begin{lemma}\label{l:rv}
Let \(A\) be an \(\mathcal F_{k-1}\)-measurable random variable such that
\(A\ge\varepsilon_h\) almost surely. If
Assumption~\ref{ass:quadratic tail} holds, then, for every
\(i\in[0:\Imax]\) and \(m\in[1:\bar m]\), almost surely
\begin{equation*}
\begin{aligned}
    \mathbb P\left(
        |\f_k^{(-1)}-f_k^{(-1)}|
        \ge
        A\Delta_k^{\Dexp}
        \,\middle|\,
        \mathcal F_{k-1}
    \right)
    &\le
    \frac{\varepsilon_h}{A^{\aexp}},
    \\
    \mathbb P\!\left(
|\tilde f_{k,m}^{(i)}-f_{k,m}^{(i)}|
\ge
A\,\Delta_k^{\Dexp}
\ \middle|\ 
\F_{k-1}
\right)
&\le
\frac{\varepsilon_h}{A^{\aexp}}.
\end{aligned}
% \qquad\text{a.s.}
\end{equation*}
\end{lemma}
\begin{proof}
Similar to \cite[Lemma~2.3]{MR4758401}. It is enough to test the desired conditional inequality against arbitrary
events $G\in\mathcal F_{k-1}$. For simple thresholds
$A=\sum_j a_j\mathbf 1_{G_j}$, with $G_j\in\mathcal F_{k-1}$ and
$a_j\ge\varepsilon_h$, the claim follows by applying
Assumption~\ref{ass:quadratic tail} on each $G\cap G_j$. The general case
$A\ge\varepsilon_h$ follows by taking simple
$\mathcal F_{k-1}$-measurable thresholds $A_n$ such that
$\varepsilon_h\le A_n\le A$ and $A_n\uparrow A$, and then passing to the limit
by dominated convergence.
\end{proof}
To prove $0\in\partial_C f(\bar x)$ at a limit point $\bar x$ of $\{X_k\}$, we proceed in two steps. First, we show that along the (infinite) set $\mathcal{U}$ of \emph{unsuccessful} iterations the forward
difference in the sampled direction is asymptotically nonnegative,
\[
\liminf_{k\in \mathcal{U},\ k\to\infty}\ \frac{f(X_k+\Delta_k D_k)-f(X_k)}{\Delta_k}\ \ge\ 0\quad\text{a.s.}
\]
which is Lemma~\ref{lem:unsuccessful-dir-nonneg} below. 
%(Lemma~\ref{lem:no-infinite-extrap} and the oracle bounds ensure the required estimates as $\Delta_k\to 0$). 
Second, we exploit that the distribution of $D_k$ has \emph{dense support} on the unit sphere: for any fixed unit vector $u$, there exist indices $k_j\in \mathcal{U}$ with $D_{k_j}\to u$ and $\Delta_{k_j}\to 0$. Passing to the limit and using the outer semicontinuity of the Clarke directional derivative then yields $f^\circ(\bar x;u)\ge 0$ for every $u$, which is equivalent to $0\in\partial_C f(\bar x)$. The next lemma establishes the first step.
\begin{lemma}\label{lem:unsuccessful-dir-nonneg}
Let $\mathcal{U}$ be the (random) set of indices of unsuccessful iterations. Under
Assumptions \ref{ass:compact}, \ref{ass:indep}, \ref{ass:quadratic tail}, and the stepsize/merit
conditions ensuring $\sum_k \EE[\Delta_k^{\Dexp}]<\infty$ (hence $\Delta_k\to 0$ a.s.), we have a.s.
\[
\liminf_{k\in \mathcal{U},\,k\to\infty}\ \frac{f(X_k+\Delta_k D_k)-f(X_k)}{\Delta_k}\ \ge\ 0.
\]
\end{lemma}
\begin{proof}
Fix $l\in\mathbb N$. On an unsuccessful outer iteration $k\in \mathcal{U}$, the inner extrapolation does not
start and for every tested direction $D_{k,m}$, with $m\in[1:\bar{m}]$, the level $i=0$ sufficient-decrease
test fails. Since, on unsuccessful iterations, no direction is accepted, we define the
direction used in the refining subsequence as the last tested direction,
namely $D_k:=D_{k,\bar{m}}$. Hence
$$
\tilde f_k^{(-1)}=\tilde f(X_k),
\qquad
\tilde f_k^{(0)}=\tilde f(X_k+\Delta_kD_k).
$$
Because the iteration is unsuccessful, the sufficient-decrease test fails at
level $i=0$ for this direction. Therefore,
\begin{equation}\label{eq:fail}
\tilde f_k^{(0)}-\tilde f_k^{(-1)}
>
-\theta\Delta_k^{\Dexp}.
\end{equation}
Apply Assumption~\ref{ass:quadratic tail} with the $\F_{k-1}$-measurable threshold
$\Upsilon_k:=\frac{1}{2l}\Delta_k^{1-\Dexp}$ (note $\Upsilon_k\ge \varepsilon_h$ for all $k$ large enough since
$\Delta_k\to 0$ a.s.). Conditionally on $\F_{k-1}$,
\[
\Pb\!\left(\,|\tilde f_k^{(i)}-f_k^{(i)}|\ \ge\ \frac{\Delta_k}{2l}\ \Big|\ \F_{k-1}\right)
\ \le\ \frac{\varepsilon_h}{\Upsilon_k^{\aexp}}
\ =\ \varepsilon_h\,(2l)^{\aexp}\,\Delta_k^{\Dexp},
\  i\in\{-1,0\}.
\]
Taking expectations and summing (Tonelli) over $k$ and $i\in\{-1,0\}$ gives
\[
\sum_{k\ge 0}\ \sum_{i\in\{-1,0\}}\Pb\!\left(|\tilde f_k^{(i)}-f_k^{(i)}|\ \ge\ \frac{\Delta_k}{2l}\right)
\ \le\ \varepsilon_h (2l)^{\aexp}\ \EE\!\left[\sum_{k\ge 0}\Delta_k^{\Dexp}\right]\ <\ \infty.
\]
By Borel--Cantelli (first lemma), there is an a.s.-finite random index $\bar{k}$ such that for all
$k\ge \bar{k}$ (in particular, for all large $k\in \mathcal{U}$),
\begin{equation}\label{eq:err-small}
|\tilde f_k^{(i)}-f_k^{(i)}|\ \le\ \frac{\Delta_k}{2l},\qquad i\in\{-1,0\}.
\end{equation}
For $k\in \mathcal{U}$ and $k\ge \bar{k}$, decompose
\begin{align*}
f(X_k+\Delta_k D_k)-f(X_k)
&=\big(f(X_k+\Delta_k D_k)-\tilde f_k^{(0)}\big)
 +\big(\tilde f_k^{(0)}-\tilde f_k^{(-1)}\big)
 +\big(\tilde f_k^{(-1)}-f(X_k)\big) \\
&\ge -\frac{\Delta_k}{2l}\ +\ \big(\tilde f_k^{(0)}-\tilde f_k^{(-1)}\big)\ -\ \frac{\Delta_k}{2l}
\qquad\text{by \eqref{eq:err-small}} \\
&\ge -\,\theta\,\Delta_k^{\Dexp}\ -\ \frac{\Delta_k}{l}
\qquad\text{by \eqref{eq:fail}}.
\end{align*}


Dividing by $\Delta_k>0$ yields, for all large $k\in \mathcal{U}$,
\[
\frac{f(X_k+\Delta_k D_k)-f(X_k)}{\Delta_k}\ \ge\ -\,\theta\,\Delta_k^{\Dexp-1}\ -\ \frac{1}{l}.
\]
Since $\Delta_k\to 0$ a.s., we obtain a.s.
\[
\liminf_{k\in \mathcal{U},\,k\to\infty}\ \frac{f(X_k+\Delta_k D_k)-f(X_k)}{\Delta_k}
\ \ge\ -\,\frac{1}{l}.
\]
As $l\in\mathbb N$ was arbitrary, letting $l\to\infty$ gives the claim.
\end{proof}
We now report the main convergence result for our stochastic direct-search scheme. The result requires the existence of accumulation points for the sequence $\left\{x_{k}\right\}$, which can be obtained assuming that the iterates generated by the algorithm lie in a compact set (see Assumption \ref{ass:compact}). We omit the proof, as it follows verbatim from \cite[Theorem 3.3]{MR4758401}.
\begin{theorem}
    % Assume that $f$ is Lipschitz continuous with constant $L$ around any limit point of the sequence of iterates $\left\{X_{k}\right\}$. Let $\mathcal{U}$ be the random set of indices of unsuccessful iterations. Let Assumptions \ref{ass:compact}, \ref{ass:indep}, \ref{ass:quadratic tail}, and \eqref{ass:eta}, \eqref{ass:theta_eta_compatibility} hold. Then, the following property holds a.s. in $\Omega$: if $V \subset \mathcal{U}$ is a random set such that the sequence $\left\{D_{k}\right\}_{k \in V}$ is dense in the unit sphere and $\lim _{k \in V, k \rightarrow \infty} X_{k}=X^{*}$, then the point $X^{*}$ is Clarke stationary, i.e., $f^{\circ}\left(X^{*}, d\right) \geq 0$ for every $d \in \mathbb{R}^{n}$.
Assume that \(f\) is Lipschitz continuous in a
neighborhood of every limit point of the sequence of iterates
\(\{X_k\}\). Let \(\mathcal U\) denote the random set of indices of
unsuccessful iterations. Suppose that
Assumptions~\ref{ass:compact}, \ref{ass:indep},
\ref{ass:quadratic tail}, and conditions~\eqref{ass:eta} and
\eqref{ass:theta_eta_compatibility} hold.
Then, with probability one, the following implication holds: if
\(V\subseteq\mathcal U\) is a random index set such that
\(\{D_k\}_{k\in V}\) is dense in \(\mathbb S^{n-1}\) and
\(
    \lim_{\substack{k\to\infty\\ k\in V}} X_k = X^*,
\)
then \(X^*\) is Clarke stationary, namely,
\(
    f^\circ(X^*;d)\ge0
\) for every $d\in\mathbb R^n$.
\end{theorem} 
% \begin{proof}
%     Refer to \cite[Theorem 3.3]{MR4758401}.
% \end{proof}

% % quick complexity analysis via merit function drift
% \input{sections/complexity delta}

% more sophisticated complexity analysis via Clarke subdifferential
% \input{sections/complexity clarke}
\section{Expected Complexity with Goldstein-stationarity stopping}\label{sec:clarke complexity}

To establish the expected iteration complexity of the proposed algorithm, we adapt the supermartingale stopping-time framework introduced in \cite{MR4151319}. This framework allows us to bound the expected number of iterations required to reach an approximate stationary condition. In contrast to the almost-sure convergence analysis of the previous section, we now characterize stationarity via the Goldstein $\goldelta$-subdifferential, $\partial_\goldelta f$, rather than the Clarke subdifferential, $\partial_C f$. The reason for this shift is detailed in Remark~\ref{rem:rescale-Lipschitz}.
% \begin{algorithm}[h]
% \caption{DSE (sequential directions, capped by $\bar{m}$)}
% \label{alg:dse2-cartis}
% \begin{algorithmic}[1]
% \State \textbf{Inputs:} $x_0\in\R^n$, $\delta_0>0$, $\theta>0$, $\gamma\in(0,1)$, $\delta_{\max}>0$, cap $\bar{m}\in\N$
% \For{$k=0,1,2,\dots$}
%   \State Draw (decision stream) the baseline estimate $\tilde f_k^{(-1)}\gets \tilde f(x_k)$
%   \For{$t=1$ \textbf{to} $\bar{m}$}
%     \State Sample $d_{k,t}\in \mathbb S^{n-1}$ independently \sara{do we prove this is dense?}% of $\F_{k-1}$
%     \State $(\delta_{\mathrm{acc}},\,\delta_{\mathrm{next}})
%       \gets \textsc{StochasticLinesearchCapped}(x_k,\delta_k,d_{k,t},\theta,\gamma,\delta_{\max})$
%     \If{$\delta_{\mathrm{acc}}>0$}
%       \State \textbf{break}
%     \EndIf
%   \EndFor
%   \State $x_{k+1}\gets x_k+\delta_{\mathrm{acc}}\,d_{k,t}$;\quad $\delta_{k+1}\gets \delta_{\mathrm{next}}$
% \EndFor
% \end{algorithmic}
% \end{algorithm}
% \begin{algorithm}[h]
%     \caption{\textsc{StochasticLinesearchCapped}}
%     \label{alg:lsep-cartis}
%     \begin{algorithmic}[1]
%         \State \textbf{Input:} $x \in \mathbb{R}^n$, $\delta > 0$, $d \in \mathbb{R}^n$, $\theta > 0$, $\gamma \in (0, 1)$, $\delta_{\max}$
%         \If{$\tilde{f}(x + \delta d) > \tilde{f}(x) - \theta \delta^{\Dexp}$}
%             \State \textbf{Return} $(0, \gamma \delta)$
%         \EndIf
%         \While{$\tilde{f}(x + \delta d) < \tilde{f}(x) - \theta \delta^{\Dexp}$}
%         \If{$\delta = \delta_{\max}$}
%             \State $\delta = \gamma^{-1}\delta_{\max}$
%             \State \textbf{break}
%         \EndIf
%             \State Set $\delta = \min\{ \ \delta / \gamma, \, \delta_{\max} \ \}$ \sara{can do $\delta_{\max} / \gamma$ directly here and remove step 6-8}
%         \EndWhile
%         \State \textbf{Return} $(\gamma \delta, \gamma \delta)$
%     \end{algorithmic}
% \end{algorithm}
We start by recalling the definition of stopping times.
\begin{definition}[Stopping time]
Let $(\mathcal A_k)_{k\ge0}$ be a filtration on
$(\Omega,\mathcal A,\mathbb P)$. A random variable
\(
    T:\Omega\to\mathbb N_0\cup\{\infty\}
\)
is called a stopping time with respect to $(\mathcal A_k)_{k\ge0}$ if
\[
    \{T\le k\}\in\mathcal A_k
    \qquad
    \forall k\in\mathbb N_0 .
\]
Equivalently, \(T\) is a stopping time if
\[
    \{T=k\}\in\mathcal A_k
    \qquad
    \forall k\in\mathbb N_0 .
\]
\end{definition}
% OLD DEFINITION. Newer one is more standard
% \begin{definition}
% Given a stochastic process $\{\Xk\} = \{\Xk : k \geq 0\}$, we say that $T$ is a stopping time with respect to $\{\Xk\}$ if for each $m \geq 0$ the occurrence of the event $\{T = m\}$ is determined by observing $X_1, \dots , X_m$. That is, $\{T = m\} \in \sigma(X_0, \dots, X_m)$, the $\sigma$-field generated by $X_1, \dots , X_m$, for each $m \geq 0$.
% \end{definition}
For the complexity analysis, we use the pre-direction algorithmic-history filtration
\[
    \mathcal A_k:=\mathcal G_{k-1},
\]
where \(\mathcal G_{k-1}\) is the information available at the beginning of
iteration \(k\), before sampling the directions and before drawing the oracle
estimates of iteration \(k\). Thus \(X_k\), \(\Delta_k\), and \(\Phi_k\) are
\(\mathcal A_k\)-measurable. Moreover, with the notation of Section \ref{sec:algo},
\[
    \mathcal A_k=\mathcal G_{k-1}
    \subseteq
    \mathcal F_{k-1}
    =
    \mathcal G_{k-1}
    \vee
    \sigma(D_{k,1},\ldots,D_{k,\bar{m}})
    \subseteq
    \mathcal A_{k+1}=\mathcal G_k .
\]
The filtration \((\mathcal F_{k-1})\) is used for oracle-conditioning
arguments after the directions have been sampled, whereas
\((\mathcal A_k)\) is used for the stopping-time and radius-dynamics
arguments. Since \(f\) is possibly nonsmooth, the stopping time is defined in terms of the
\(\goldelta\)-Goldstein subdifferential.
\begin{definition}[Goldstein subdifferential]
For \(r>0\), let
$B_r(x):=\{y\in\mathbb{R}^n:\|y-x\|\le r\}
$
denote the closed ball of radius \(r\) centered at \(x\). The
\emph{\(r\)-Goldstein subdifferential} of \(f\) at \(x\) is defined as the
convex hull of the Clarke subdifferentials at all points \(y\in B_r(x)\), namely
\[
\partial_r f(x)
:=
\operatorname{conv}\left(
\bigcup_{y\in B_r(x)} \partial_C f(y)
\right).
\]
\end{definition}
For fixed tolerances $\goldelta>0$ and $\varepsilon>0$, we define the stopping time as
\begin{equation}\label{Teps}
    T_{\goldelta, \varepsilon} = \inf\left\{ k \in \mathbb{N}_0 : \min_{g \in \partial_\goldelta f(\Xk)} \|g\| \leq \varepsilon \right\}.
\end{equation}
% \begin{equation}\label{Teps}
%     T_{\goldelta, \varepsilon} = \inf\left\{ k \in \mathbb{N} : \|g_k\| \leq \varepsilon,\  g_k \in \argmin \|g\| \ | \ g \in \partial_\goldelta f(\Xk) \right\}.
% \end{equation}
Equivalently, \(T_{\goldelta,\varepsilon}\) is the first iteration at which
\(X_k\) is a \((\goldelta,\varepsilon)\)-Goldstein stationary point.\\
Since the definition only depends on \((\Xk)\) which are $\mathcal A_{k}$ measurable, it follows that $T_{\goldelta, \varepsilon}$ is indeed a ($\mathcal A_k$)-stopping time. Adapting the framework of \cite{MR4151319}, we can provide a bound on $\me[T_{\goldelta, \varepsilon}]$. At a given iteration \(k\), define the base-estimate event
\[
\mathcal E_{k,0}
:=
\left\{
    |\tilde f(X_k)-f(X_k)|
    \le
    \varepsilon_f\Delta_k^{\Dexp}
\right\},
\]
and, for each \(m=1,\ldots,\bar{m}\), the trial-estimate event
\[
\mathcal E_{k,m}
:=
\left\{
    |\tilde f(X_k+\Delta_kD_{k,m})
      -f(X_k+\Delta_kD_{k,m})|
    \le
    \varepsilon_f\Delta_k^{\Dexp}
\right\}.
\]
Then, fix a stationarity tolerance
\(\varepsilon>0\) and a Goldstein neighborhood radius \(\goldelta>0\).
Let
\[
    G_{\goldelta,k}^*
    \in
    \operatorname*{argmin}_{g\in\partial_{\goldelta}f(X_k)}\|g\|,
\]
and, on the event \(\{T_{\goldelta,\varepsilon}>k\}\), define
\[
    D_k^*
    :=
    -\frac{G_{\goldelta,k}^*}{\|G_{\goldelta,k}^*\|}.
\]
Indeed, on this event we have \(\|G_{\goldelta,k}^*\|>\varepsilon\), so
\(D_k^*\) is well-defined. The direction \(D_k^*\) is the Goldstein descent direction at iteration $k$. Since the algorithm
samples random directions, we only need one sampled direction to be sufficiently
aligned with \(D_k^*\). To quantify the probability that a randomly sampled direction is sufficiently
aligned with a target descent direction, we use the standard formula for the
surface measure of a spherical cap. 

\begin{definition}[Spherical cap]
\label{def:spherical-cap-probability}
Fix \(0<\rho<1\) and a unit vector \(u\in\mathbb S^{n-1}\). The
\emph{spherical cap with center \(u\) and alignment parameter \(\rho\)} is
\[
    \mathcal C_\rho(u)
    :=
    \{d\in\mathbb S^{n-1}:\langle u,d\rangle\ge\rho\}.
\]
Equivalently, \(\mathcal C_\rho(u)\) is the cap with half-angle
\(\arccos(\rho)\). 
\end{definition}
Thus, \(D_{k,m}\in\mathcal C_\rho(D_k^*)\) means that the sampled direction
\(D_{k,m}\) is sufficiently aligned with the Goldstein descent direction at iteration $k$. Now, define the good event that both estimates are accurate and  $\exists\ m\in\{1,\ldots,\bar{m}\}:
D_{k,m}\in \mathcal C_\rho(D_k^\ast)$, with \(D_k^*\) being the Goldstein descent direction at iteration $k$, and \(C_\rho(D_k^\ast)\) being the spherical cap therein centered. That is, define
\[
\mathcal I_k
:=
\mathcal E_{k,0}
\cap
\left[
\bigcup_{m=1}^{\bar{m}}
\left(
\{D_{k,m}\in\mathcal C_\rho(D_k^*)\}
\cap
\mathcal E_{k,m}
\right)
\right],
\]
and let $J_k:=\mathbf 1_{\mathcal I_k}. $
We then set
\begin{equation} \label{eq:Wk}
    W_{k+1}:=2J_k-1\in\{-1,+1\}.
\end{equation}
Hence \(W_{k+1}=+1\) if the good event \(\mathcal I_k\) occurs, and
\(W_{k+1}=-1\) otherwise. In particular,
we will show that,
% \(\mathcal I_k\) will be chosen so that,
whenever \(T_{\goldelta,\varepsilon}>k\) and
\(\Delta_k\le \bar\Delta_{\goldelta,\varepsilon}\), the occurrence of \(\mathcal I_k\)
implies that iteration \(k\) is successful. We set \(W_0=1\) only for
initialization. Notice that \(W_k\) is \(\mathcal A_k\)-measurable, whereas
\(W_{k+1}\) is not \(\mathcal A_k\)-measurable. We now state the abstract conditions needed to apply the stopping-time
framework (see also \cite[Assumption 1]{MR4151319}).

\begin{assumption}
\label{ass:5}
Let \(T_{\goldelta,\varepsilon}\) be the stopping time defined in
\eqref{Teps}. Let the stochastic process \(\{(\Phi_k,\Delta_k,W_k)\}_{k\ge0}\) be adapted to
\((\mathcal A_k)_{k\ge0}\). The following conditions hold.

\begin{enumerate}
    \item[(i)] There exist constants \(\lambda>0\) and
    \(\Delta_{\max}= \Delta_0 e^{\lambda j_{max}}>0\), \(j_{max}\in\mathbb Z\) such that
    \[
        \lambda=-\log\gamma
        \qquad\text{and}\qquad
        \Delta_k\le \Delta_{\max}
        \quad\text{for all }k\ge0.
    \]

    \item[(ii)] There exist a threshold
    \[
        \bar\Delta_{\goldelta,\varepsilon}
        =
        \Delta_0 e^{\lambda j_{\goldelta,\varepsilon}},
        \qquad
        j_{\goldelta,\varepsilon}\in\mathbb Z,
        \quad j_{\goldelta,\varepsilon}\le 0,
    \]
    and a constant \(q>1/2\) such that
    \[
        \mathbb P(W_{k+1}=+1\mid \mathcal A_k)= q,
        \qquad
        \mathbb P(W_{k+1}=-1\mid \mathcal A_k)= 1-q,
    \]
    and, for every \(k\ge0\),
    \begin{equation}
    \label{eq:dynamics}
        \mathbf 1_{\{T_{\goldelta,\varepsilon}>k\}}
        \Delta_{k+1}
        \ge
        \mathbf 1_{\{T_{\goldelta,\varepsilon}>k\}}
        \min\left\{
            \Delta_k e^{\lambda W_{k+1}},
            \bar\Delta_{\goldelta,\varepsilon}
        \right\}.
    \end{equation}

    \item[(iii)] There exist a nondecreasing function
    \(h:[0,\infty)\to [0,\infty)\) and a constant \(\Theta>0\) such that,
    for every \(k\ge0\),
    \begin{equation}
    \label{eq:abstract_drift}
        \mathbb E[
            \Phi_k-\Phi_{k+1}
            \mid
            \mathcal A_k
        ]
        \mathbf 1_{\{T_{\goldelta,\varepsilon}>k\}}
        \ge
        \Theta h(\Delta_k)
        \mathbf 1_{\{T_{\goldelta,\varepsilon}>k\}}
        \qquad\text{a.s.}
    \end{equation}
    In particular, we take
    $h(\Delta)=\Delta^{\Dexp}. $
\end{enumerate}
\end{assumption}
Assumption~\ref{ass:5} states that, before the stopping time, the
nonnegative stochastic process \(\Phi_k\) has an expected decrease of at least
\(\Theta h(\Delta_k)\) at each iteration. Moreover, before the stopping time
and whenever \(\Delta_k\le \bar\Delta_{\goldelta,\varepsilon}\), the radius
has an upward tendency: the good event \(\mathcal I_k\), which implies
\(W_{k+1}=+1\), occurs with conditional probability larger that  \(1/2\), namely
\[\mathbb P(W_{k+1}=+1\mid\mathcal A_k)
   = q >1/2.
\]
Our goal is to bound \(\mathbb E[T_{\goldelta,\varepsilon}]\) in terms of
\(h(\bar\Delta_{\goldelta,\varepsilon})\). What happens is that, on average, $\D_k \geq \bar \D_{\goldelta,\varepsilon}$ frequently, and hence, $\me[\Phi_{k+1} - \Phi_k]$ can be bounded by a negative fixed value (dependent on $\goldelta, \varepsilon$), sufficiently frequently, which will allow us
to apply Wald’s identity and derive the bound on $\me[T_{\goldelta,\varepsilon}]$ (for more details refer to \cite[Theorem 2]{MR4151319}). 
The next result is a direct specialization of the supermartingale stopping-time bound of Blanchet, Cartis, Menickelly and Scheinberg~\cite[Theorem~2]{MR4151319} to our framework.
\begin{theorem}\label{t:expected complexity}
    Let Assumption~\ref{ass:5} hold. Then
    \[
        \mathbb E[T_{\goldelta,\varepsilon}]
        \le
        \frac{q}{2q-1}
        \cdot
        \frac{\Phi_0}{\Theta h(\bar\Delta_{\goldelta,\varepsilon})}
        +1 .
    \]
\end{theorem}
Assumption~\ref{ass:5}(i) follows from the compactness assumption (see Remark \ref{rem:compact-tested-points}). For Assumption~\ref{ass:5}(iii), the merit-function drift established in Lemma~\ref{lem:phi_drift_main} gives
\[
    \mathbb E[
        \Phi_k-\Phi_{k+1}
        \mid
        \mathcal F_{k-1}
    ]
    \ge
    \Theta\Delta_k^{\Dexp},
\]
where
    $\Theta:=\eta(1-\gamma^{\Dexp})-c_{\max}>0.$
Since
\(
    \mathcal A_k\subseteq\mathcal F_{k-1},
\)
the tower property gives
\[
\begin{aligned}
    \mathbb E[
        \Phi_k-\Phi_{k+1}
        \mid
        \mathcal A_k
    ]
    =
    \mathbb E\!\left[
        \mathbb E[
            \Phi_k-\Phi_{k+1}
            \mid
            \mathcal F_{k-1}
        ]
        \,\middle|\,
        \mathcal A_k
    \right]                                                   
    \ge
    \Theta\Delta_k^{\Dexp}.
\end{aligned}
\]
Multiplying by the \(\mathcal A_k\)-measurable indicator
\(\mathbf 1_{\{T_{\goldelta,\varepsilon}>k\}}\), we obtain
\[
    \mathbb E[
        \Phi_k-\Phi_{k+1}
        \mid
        \mathcal A_k
    ]
    \mathbf 1_{\{T_{\goldelta,\varepsilon}>k\}}
    \ge
    \Theta\Delta_k^{\Dexp}
    \mathbf 1_{\{T_{\goldelta,\varepsilon}>k\}},
\]
which is Assumption~\ref{ass:5}(iii).\\
It remains to verify Assumption~\ref{ass:5}(ii), namely the radius dynamics and
the lower bound
\[
    \mathbb P(W_{k+1}=+1\mid\mathcal A_k)\ge q>1/2
\]
before the stopping time and for sufficiently small \(\Delta_k\). The next two
lemmas provide exactly this. The first lemma proves a uniform lower bound
on the probability that at least one of the \(\bar{m}\) sampled directions yields a
sufficient-decrease step. The second lemma shows that, with the given stepsize update rule, the radius process satisfies the required comparison
inequality. Before stating the lemmas, we report the following definition that collects the
notation for the regularized incomplete Beta function and the cap
probability under the uniform distribution of directions on the sphere.

\begin{definition}[Regularized incomplete Beta function and Spherical-cap probability]
\label{def:reg-inc-beta-funct}
 For \(a,b>0\) and \(z\in[0,1]\), the
\emph{regularized incomplete Beta function} is
\[
    I_z(a,b)
    :=
    \frac{B(z;a,b)}{B(a,b)},
    \qquad
    B(z;a,b)
    :=
    \int_0^z t^{a-1}(1-t)^{b-1}\,dt,
\]
where
\[
    B(a,b)
    :=
    \frac{\Gamma(a)\Gamma(b)}{\Gamma(a+b)}.
\]
Equivalently, if \(U\sim\mathrm{Beta}(a,b)\), then
\[
    I_z(a,b)=\mathbb P(U\le z).
\]
Now, assume \(n\ge2\). By standard properties of a uniformly distributed point on the sphere (see, e.g., \cite[Chapter~9]{MardiaJupp2000}),
the surface measure of the spherical cap \(\mathcal C_\rho(u)\) is independent
of its center \(u\). More precisely, if \(D\sim\mathrm{Unif}(\mathbb S^{n-1})\), then
\begin{equation}\label{capprob}
    \mathbb P(D\in\mathcal C_\rho(u))
    =
    p_\rho
    :=
    \frac12
    I_{1-\rho^2}
    \left(
        \frac{n-1}{2},
        \frac12
    \right).
\end{equation}
 Moreover,
$p_\rho>0
    \ \text{for every fixed }\rho\in(0,1),
    \ \text{and}\ 
    p_\rho\downarrow0
     \text{ as }\rho\uparrow1.$

\end{definition}
If
\(D_{k,m}\) is uniformly distributed on \(\mathbb S^{n-1}\), then $\mathbb P(D_{k,m}\in \mathcal C_\rho(D_k^*)\mid\mathcal A_k)
    =
    p_\rho.$
We first show that, before the stopping time and for sufficiently small
stepsizes, the occurrence of the good event \(\mathcal I_k\) forces the algorithmic
sufficient-decrease test to succeed. This is the analogue, in the present
Goldstein-stationarity setting, of the key implication used in Dzahini's
analysis \cite{MR4359469}: if the stationarity measure is still larger than the prescribed
tolerance, the stepsize is small enough, and the estimates are accurate, then
the iteration must be successful. Here, the role of the descent direction is
played by any sampled direction lying in the spherical cap~
\(\mathcal C_\rho(D_k^\ast)\).

\begin{lemma}[Sufficient decrease on the good event]
\label{lem:good-event-sufficient-decrease}
Assume
\[
   c_1
   :=
   \rho\varepsilon-\sqrt{1-\rho^{2}}\,L
   >
   0 .
\]
Let \(\theta>0\) be the sufficient-decrease parameter used by the algorithm, and
choose
\[
  \Delta_{\goldelta,\varepsilon}
  <
  \min\left\{
      \goldelta,
      \left(
          \frac{c_1}{\theta+2\varepsilon_f}
      \right)^{1/(\Dexp-1)}
  \right\}.
\]
Let
\[
    B_k
    :=
    \{T_{\goldelta,\varepsilon}>k\}
    \cap
    \{\Delta_k\le\Delta_{\goldelta,\varepsilon}\}.
\]
% Then, on \(B_k\), if there exists \(m\in\{1,\ldots,\bar{m}\}\) such that
% \[
%     D_{k,m}\in\mathcal C_\rho(D_k^*)
% \]
% and both estimates \(\tilde f(X_k)\) and
% \(\tilde f(X_k+\Delta_kD_{k,m})\) are \(\varepsilon_f\)-accurate \sara{why not use the notation $\mathcal I_k$ defined earlier?}, we have
% \[
%     \tilde f(X_k+\Delta_kD_{k,m})-\tilde f(X_k)
%     \le
%     -\theta\Delta_k^{\Dexp}.
% \]
% In particular, iteration \(k\) is successful.
Then, on the event \(B_k\cap\mathcal I_k\), there exists
\(m\in\{1,\ldots,\bar m\}\) such that
\[
    D_{k,m}\in\mathcal C_\rho(D_k^*)
\]
and both the base estimate and the corresponding trial estimate are
\(\varepsilon_f\)-accurate. For such an index \(m\),
\[
    \tilde f(X_k+\Delta_kD_{k,m})
    -
    \tilde f(X_k)
    \le
    -\theta\Delta_k^{\Dexp}.
\]
Consequently, iteration \(k\) is successful.
\end{lemma}

\begin{proof}
We argue on the event \(B_k\). By definition of
\(T_{\goldelta,\varepsilon}\),
\[
    \min_{g\in\partial_{\goldelta} f(X_k)}\|g\|>\varepsilon .
\]
Choose
\[
    G_{\goldelta,k}^*
    \in
    \operatorname*{argmin}_{g\in\partial_{\goldelta} f(X_k)}
    \|g\|,
    \qquad
    D_k^*
    :=
    -\frac{G_{\goldelta,k}^*}{\|G_{\goldelta,k}^*\|}.
\]
Then \(\|G_{\goldelta,k}^*\|>\varepsilon\).

Fix \(m\in\{1,\ldots,\bar{m}\}\) such that
\(D_{k,m}\in\mathcal C_\rho(D_k^*)\), and set $ S_{k,m}:=\Delta_kD_{k,m}.$
By Lebourg's mean-value theorem, there exist
\[
    \bar X_{k,m}=X_k+\Xi_{k,m}S_{k,m},
    \qquad
    \Xi_{k,m}\in(0,1),
\]
and
\[
    \bar G_{k,m}\in\partial_C f(\bar X_{k,m})
\]
such that
\[
    f(X_k+S_{k,m})-f(X_k)
    =
    \langle \bar G_{k,m},S_{k,m}\rangle
    =
    \Delta_k\langle \bar G_{k,m},D_{k,m}\rangle .
\]
Since \(\Delta_k\le\Delta_{\goldelta,\varepsilon}\le\goldelta\), we have
\[
    \|\bar X_{k,m}-X_k\|\le\Delta_k\le\goldelta .
\]
Therefore \(\bar X_{k,m}\in B_{\goldelta}(X_k)\), and hence
\[
    \bar G_{k,m}\in\partial_{\goldelta}f(X_k).
\]
By the projection theorem applied to the closed convex set
\(\partial_{\goldelta} f(X_k)\),
\[
    \langle \bar G_{k,m},G_{\goldelta,k}^*\rangle
    \ge
    \|G_{\goldelta,k}^*\|^2 .
\]
Consequently,
\[
    \langle \bar G_{k,m},D_k^*\rangle
    =
    -\frac{
        \langle \bar G_{k,m},G_{\goldelta,k}^*\rangle
    }{
        \|G_{\goldelta,k}^*\|
    }
    \le
    -\|G_{\goldelta,k}^*\|
    <
    -\varepsilon .
\]
Write
\[
    D_{k,m}
    =
    R_{k,m}D_k^*
    +
    \sqrt{1-R_{k,m}^2}\,V_{k,m},
\]
where
\[
    R_{k,m}:=\langle D_k^*,D_{k,m}\rangle,
    \qquad
    V_{k,m}\perp D_k^*,
    \qquad
    \|V_{k,m}\|=1.
\]
Since \(D_{k,m}\in\mathcal C_\rho(D_k^*)\), we have
\(R_{k,m}\ge\rho\). Moreover, since \(f\) is \(L\)-Lipschitz on the relevant
neighborhood, \(\|\bar G_{k,m}\|\le L\). Thus
\[
\begin{aligned}
    \langle \bar G_{k,m},D_{k,m}\rangle
    &=
    R_{k,m}\langle \bar G_{k,m},D_k^*\rangle
    +
    \sqrt{1-R_{k,m}^2}
    \langle \bar G_{k,m},V_{k,m}\rangle                                      \\
    &\le
    -R_{k,m}\varepsilon
    +
    \sqrt{1-R_{k,m}^2}\,L                                                   \\
    &\le
    -\rho\varepsilon+\sqrt{1-\rho^2}\,L                                      \\
    &=
    -c_1 .
\end{aligned}
\]
Hence
\[
    f(X_k+\Delta_kD_{k,m})-f(X_k)
    \le
    -c_1\Delta_k.
\]
If the two estimates are \(\varepsilon_f\)-accurate, then
\[
\begin{aligned}
    \tilde f(X_k+\Delta_kD_{k,m})-\tilde f(X_k)
    &\le
    f(X_k+\Delta_kD_{k,m})-f(X_k)
    +
    2\varepsilon_f\Delta_k^{\Dexp}                                           \\
    &\le
    -c_1\Delta_k
    +
    2\varepsilon_f\Delta_k^{\Dexp}.
\end{aligned}
\]
By the definition of \(\Delta_{\goldelta,\varepsilon}\),
\[
    \Delta_k^{\Dexp-1}
    \le
    \Delta_{\goldelta,\varepsilon}^{\Dexp-1}
    <
    \frac{c_1}{\theta+2\varepsilon_f}.
\]
Therefore
\[
    c_1\Delta_k
    >
    (\theta+2\varepsilon_f)\Delta_k^{\Dexp}.
\]
It follows that
\[
    \tilde f(X_k+\Delta_kD_{k,m})-\tilde f(X_k)
    \le
    -\theta\Delta_k^{\Dexp}.
\]
Thus the sufficient-decrease test succeeds for direction \(m\), and iteration
\(k\) is successful.
\end{proof}



% \begin{lemma}[Probabilistic sufficient decrease with Goldstein subdifferential]
% \label{lem:random-cap}
% Assume
% \[
%    c_1
%    :=
%    \rho\varepsilon-\sqrt{1-\rho^{2}}\,L
%    >
%    0
%    \qquad
%    \text{(cf.~Remark~\ref{rem:rescale-Lipschitz})}.
% \]
% Let \(\theta>0\) be the sufficient-decrease parameter used by the algorithm, and
% choose
% \[
%   \Delta_{\goldelta,\varepsilon}
%   <
%   \min\left\{
%       \goldelta,
%       \left(
%           \frac{c_1}{\theta+2\varepsilon_f}
%       \right)^{1/(\Dexp-1)}
%   \right\}.
% \]
% Assume that, conditionally on \(\mathcal A_k\), the directions
% \(D_{k,1},\ldots,D_{k,\bar{m}}\) are independent and uniformly distributed on
% \(\mathbb S^{n-1}\). Assume also that the estimates are
% \(\beta\)-probabilistically \(\varepsilon_f\)-accurate, namely
% \[
%     \mathbb P(\mathcal E_{k,0}\mid\mathcal A_k)\ge\beta,
%     \qquad
%     \mathbb P(\mathcal E_{k,m}\mid\mathcal A_k,D_{k,m})\ge\beta,
%     \quad m=1,\ldots,\bar{m}.
% \]
% Moreover, assume that, conditionally on \(\mathcal A_k\), the base-estimate
% event \(\mathcal E_{k,0}\) is independent of the trial events, and that the
% pairs
% \[
%     \{(D_{k,m},\mathcal E_{k,m})\}_{m=1}^{\bar{m}}
% \]
% are conditionally independent.
% Let us consider the event: 
% \[
%     B_k
%     :=
%     \{T_{\goldelta,\varepsilon}>k\}
%     \cap
%     \{\Delta_k\le\Delta_{\goldelta,\varepsilon}\}.
% \]
% Then
% \[
%     \mathbf 1_{B_k}\,
%     \mathbb P(J_k=1\mid\mathcal A_k)
%     =
%     \mathbf 1_{B_k}\,
%     \mathbb P(W_{k+1}=+1\mid\mathcal A_k)
%     \ge
%     q_M\,\mathbf 1_{B_k},
% \]
% with 
% \[
%     q_M
%     :=
%     \beta\left[1-(1-p_\rho\beta)^{\bar{m}}\right].
% \]

% \end{lemma}


% \begin{lemma}[Probabilistic sufficient decrease with Goldstein subdifferential]
% \label{lem:random-cap}
% Fix a stationarity tolerance \(\varepsilon>0\), a Goldstein neighborhood radius
% \(\goldelta>0\), and choose a spherical cap parameter
% \[
%    \rho\in(0,1),
%    \qquad
%    p_\rho
%    :=
%    \frac12\,
%    I_{1-\rho^{2}}
%    \left(
%         \frac{n-1}{2},
%         \frac12
%    \right)
%    \in(0,1),
% \]
% where \(I\) is the regularized incomplete Beta function
% (cf.~Remark~\ref{rem:reg-inc-beta}). Assume
% \[
%    c_1
%    :=
%    \rho\varepsilon-\sqrt{1-\rho^{2}}\,L
%    >
%    0
%    \qquad
%    \text{(cf.~Remark~\ref{rem:rescale-Lipschitz})}.
% \]
% Let \(\theta>0\) be the sufficient-decrease parameter used by the algorithm, and
% choose
% \[
%   \Delta_{\goldelta,\varepsilon}
%   <
%   \min\left\{
%       \goldelta,
%       \left(
%           \frac{c_1}{\theta+2\varepsilon_f}
%       \right)^{1/(\Dexp-1)}
%   \right\}.
% \]
% Define
% \[
%     q_M
%     :=
%     \beta\left[1-(1-p_\rho\beta)^{\bar{m}}\right].
% \]
% Assume that, conditionally on \(\mathcal A_k\), the directions
% \(D_{k,1},\ldots,D_{k,\bar{m}}\) are independent and uniformly distributed on
% \(\mathbb S^{n-1}\). Define the base-estimate event
% \[
% \mathcal E_{k,0}
% :=
% \left\{
%     |\tilde f(X_k)-f(X_k)|
%     \le
%     \varepsilon_f\Delta_k^{\Dexp}
% \right\},
% \]
% and, for each \(m=1,\ldots,\bar{m}\), the trial-estimate event
% \[
% \mathcal E_{k,m}
% :=
% \left\{
%     |\tilde f(X_k+\Delta_kD_{k,m})
%       -f(X_k+\Delta_kD_{k,m})|
%     \le
%     \varepsilon_f\Delta_k^{\Dexp}
% \right\}.
% \]
% Assume that the estimates are \(\beta\)-probabilistically
% \(\varepsilon_f\)-accurate, in the sense that
% \[
%     \mathbb P(\mathcal E_{k,0}\mid\mathcal A_k)\ge\beta, \quad \text{and} \quad \mathbb P(\mathcal E_{k,m}\mid\mathcal A_k,D_{k,m})\ge\beta,
%     \quad m=1,\ldots,\bar{m}.
% \]
% Assume also that, conditionally on \(\mathcal A_k\), the base-estimate event
% \(\mathcal E_{k,0}\) is independent of the trial events, and that the pairs
% \(
%     \{(D_{k,m},\mathcal E_{k,m})\}_{m=1}^{\bar{m}}
% \)
% are conditionally independent. Then, on the event
% \(
%     \{T_{\goldelta,\varepsilon}>k\}
%     \cap
%     \{\Delta_k\le\Delta_{\goldelta,\varepsilon}\},
% \)
% we have
% \[
%    \mathbb P(J_k=1\mid\mathcal A_k)
%    =
%    \mathbb P(W_{k+1}=+1\mid\mathcal A_k)
%    \ge
%    q_M .
% \]
% \end{lemma}

% \begin{proof}
% We prove the desired lower bound on the event \(B_k\). On \(B_k\), by definition of \(T_{\goldelta,\varepsilon}\), the minimum-norm
% element of the Goldstein \(\goldelta\)-subdifferential satisfies
% \[
%     \min_{g\in\partial_{\goldelta} f(X_k)}\|g\|>\varepsilon .
% \]
% Choose
% \[
%     G_{\goldelta,k}^*
%     \in
%     \operatorname*{argmin}_{g\in\partial_{\goldelta} f(X_k)}
%     \|g\|.
% \]
% Then
% \[
%     \|G_{\goldelta,k}^*\|>\varepsilon.
% \]
% We defined the normalized Goldstein descent direction as follows: 
% \[
%     D_k^*
%     :=
%     -\frac{G_{\goldelta,k}^*}{\|G_{\goldelta,k}^*\|}.
% \]
% By Definition~\ref{def:spherical-cap-probability}, the spherical cap centered
% at \(D_k^*\) is
% \[
%     \mathcal C_\rho(D_k^*)
%     :=
%     \left\{
%         d\in\mathbb S^{n-1}:
%         \langle D_k^*,d\rangle\ge\rho
%     \right\}.
% \]
% Since \(D_{k,m}\) is uniformly distributed on \(\mathbb S^{n-1}\) conditionally
% on \(\mathcal A_k\),  from equation \eqref{capprob} we have
% \[
%     \mathbb P(D_{k,m}\in\mathcal C_\rho(D_k^*)\mid\mathcal A_k)=p_\rho .
% \]

% Fix \(m\in\{1,\ldots,\bar{m}\}\) and suppose that $D_{k,m}\in\mathcal C_\rho(D_k^*) $.
% Set $S_{k,m}:=\Delta_kD_{k,m}.$
% By Lebourg's mean-value theorem, there exist
% \[
%     \bar X_{k,m}
%     =
%     X_k+\Xi_{k,m}S_{k,m},
%     \qquad
%     \Xi_{k,m}\in(0,1),
% \]
% and
% \[
%     \bar G_{k,m}\in\partial_C f(\bar X_{k,m})
% \]
% such that
% \[
%     f(X_k+S_{k,m})-f(X_k)
%     =
%     \langle \bar G_{k,m},S_{k,m}\rangle
%     =
%     \Delta_k\langle \bar G_{k,m},D_{k,m}\rangle .
% \]

% Since, on \(B_k\), $\Delta_k\le\Delta_{\goldelta,\varepsilon}\le\goldelta,$ 
% we have $\|\bar X_{k,m}-X_k\|
%     \le
%     \Delta_k
%     \le
%     \goldelta .$
% Therefore $\bar X_{k,m}\in B_{\goldelta}(X_k),$
% and hence $ \bar G_{k,m}\in\partial_{\goldelta}f(X_k).$
% By the projection theorem applied to the closed convex set
% \(\partial_{\goldelta} f(X_k)\), we have $\langle \bar G_{k,m},G_{\goldelta,k}^*\rangle
%     \ge
%     \|G_{\goldelta,k}^*\|^2 .
% $
% Consequently,
% \[
%     \langle \bar G_{k,m},D_k^*\rangle
%     =
%     -\frac{
%         \langle \bar G_{k,m},G_{\goldelta,k}^*\rangle
%     }{
%         \|G_{\goldelta,k}^*\|
%     }
%     \le
%     -\|G_{\goldelta,k}^*\|
%     <
%     -\varepsilon .
% \]

% Now decompose \(D_{k,m}\) as
% \[
%     D_{k,m}
%     =
%     R_{k,m}D_k^*
%     +
%     \sqrt{1-R_{k,m}^2}\,V_{k,m},
% \]
% where
% \[
%     R_{k,m}:=\langle D_k^*,D_{k,m}\rangle,
%     \qquad
%     V_{k,m}\perp D_k^*,
%     \qquad
%     \|V_{k,m}\|=1.
% \]
% On the event \(D_{k,m}\in\mathcal C_\rho(D_k^*)\), we have
% \[
%     R_{k,m}\ge\rho.
% \]
% Moreover, since \(f\) is \(L\)-Lipschitz on the relevant neighborhood, every
% Clarke subgradient appearing above has norm at most \(L\). Thus
% \[
%     \langle \bar G_{k,m},V_{k,m}\rangle
%     \le
%     \|\bar G_{k,m}\|
%     \le
%     L.
% \]
% Therefore
% \[
% \begin{aligned}
%     \langle \bar G_{k,m},D_{k,m}\rangle
%     &=
%     R_{k,m}\langle \bar G_{k,m},D_k^*\rangle
%     +
%     \sqrt{1-R_{k,m}^2}
%     \langle \bar G_{k,m},V_{k,m}\rangle                                      \\
%     &\le
%     -R_{k,m}\varepsilon
%     +
%     \sqrt{1-R_{k,m}^2}\,L                                                     \\
%     &\le
%     -\rho\varepsilon+\sqrt{1-\rho^2}\,L                                      \\
%     &=
%     -c_1 .
% \end{aligned}
% \]
% Hence
% \[
%     f(X_k+\Delta_kD_{k,m})-f(X_k)
%     \le
%     -c_1\Delta_k.
% \]

% On the event \(\mathcal E_{k,0}\cap\mathcal E_{k,m}\), we then have
% \[
% \begin{aligned}
%     \tilde f(X_k+\Delta_kD_{k,m})-\tilde f(X_k)
%     &\le
%     f(X_k+\Delta_kD_{k,m})-f(X_k)
%     +
%     2\varepsilon_f\Delta_k^{\Dexp}                                           \\
%     &\le
%     -c_1\Delta_k
%     +
%     2\varepsilon_f\Delta_k^{\Dexp}.
% \end{aligned}
% \]
% By the definition of \(\Delta_{\goldelta,\varepsilon}\),
% \[
%     \Delta_k^{\Dexp-1}
%     \le
%     \Delta_{\goldelta,\varepsilon}^{\Dexp-1}
%     <
%     \frac{c_1}{\theta+2\varepsilon_f}.
% \]
% Therefore
% \[
%     c_1\Delta_k
%     >
%     (\theta+2\varepsilon_f)\Delta_k^{\Dexp}.
% \]
% It follows that
% \[
%     \tilde f(X_k+\Delta_kD_{k,m})-\tilde f(X_k)
%     \le
%     -\theta\Delta_k^{\Dexp}.
% \]
% Thus the sufficient-decrease test succeeds for direction \(m\), and iteration
% \(k\) is successful.

% Now define
% \[
%     \mathcal S_{k,m}
%     :=
%     \{D_{k,m}\in\mathcal C_\rho(D_k^*)\}
%     \cap
%     \mathcal E_{k,m}.
% \]
% From the conditional uniformity of \(D_{k,m}\) and the conditional
% \(\beta\)-accuracy of the trial estimate,
% \[
% \begin{aligned}
%     \mathbb P(\mathcal S_{k,m}\mid\mathcal A_k)
%     &=
%     \mathbb E\!\left[
%         \mathbf 1_{\{D_{k,m}\in\mathcal C_\rho(D_k^*)\}}
%         \mathbb P(\mathcal E_{k,m}\mid\mathcal A_k,D_{k,m})
%         \,\middle|\,
%         \mathcal A_k
%     \right]                                      \\
%     &\ge
%     \beta\,
%     \mathbb P(D_{k,m}\in\mathcal C_\rho(D_k^*)\mid\mathcal A_k)              \\
%     &=
%     p_\rho\beta .
% \end{aligned}
% \]
% By the conditional independence assumption, the events
% \(\mathcal S_{k,1},\ldots,\mathcal S_{k,\bar{m}}\) are conditionally independent
% given \(\mathcal A_k\). Hence
% \[
%     \mathbb P\left(
%         \bigcup_{m=1}^{\bar{m}} \mathcal S_{k,m}
%         \,\middle|\,
%         \mathcal A_k
%     \right)
%     \ge
%     1-(1-p_\rho\beta)^{\bar{m}}.
% \]
% Furthermore, \(\mathcal E_{k,0}\) is conditionally independent of these trial
% events and satisfies
% \[
%     \mathbb P(\mathcal E_{k,0}\mid\mathcal A_k)\ge\beta.
% \]
% Therefore, on \(B_k\),
% \[
% \begin{aligned}
%     \mathbb P(J_k=1\mid\mathcal A_k)
%     &\ge
%     \mathbb P\left(
%         \mathcal E_{k,0}
%         \cap
%         \bigcup_{m=1}^{\bar{m}} \mathcal S_{k,m}
%         \,\middle|\,
%         \mathcal A_k
%     \right)                                                               \\
%     &\ge
%     \beta\left[1-(1-p_\rho\beta)^{\bar{m}}\right]                                  \\
%     &=
%     q_M.
% \end{aligned}
% \]
% Since \(W_{k+1}=+1\) if and only if \(J_k=1\), this gives
% \[
%     \mathbb P(W_{k+1}=+1\mid\mathcal A_k)\ge q_M
% \]
% on \(B_k\). Equivalently,
% \[
%     \mathbf 1_{B_k}\,
%     \mathbb P(W_{k+1}=+1\mid\mathcal A_k)
%     \ge
%     q_M\,\mathbf 1_{B_k}.
% \]
% This proves the claim.
% \end{proof}


% \begin{proof}
% Work on the event
% \((
%     \{T_{\goldelta,\varepsilon}>k\}
%     \cap
%     \{\Delta_k\le\Delta_{\goldelta,\varepsilon}\}.
% \)
% By definition of \(T_{\goldelta,\varepsilon}\), the minimum-norm element of the Goldstein $\goldelta$-subdifferential $g_{\goldelta, k}^* := \operatorname{argmin}_{g \in \partial_\goldelta^G f(x_k)} \|g\|$ satisfies $\|g_{\goldelta, k}^*\| > \varepsilon$.
% Define the normalized Goldstein descent direction
% \[
%     d_k^*
%     :=
%     -\frac{g_{\goldelta,k}^*}{\|g_{\goldelta,k}^*\|}.
% \]
% For each \(m=1,\ldots,\bar{m}\), define the spherical cap
% \[
%    \mathcal C_k
%    :=
%    \left\{
%         d\in\mathbb S^{n-1}:
%         \langle d_k^*,d\rangle\ge\rho
%    \right\}.
% \]
% Since \(D_{k,m}\) is uniformly distributed on \(\mathbb S^{n-1}\) conditionally
% on \(\mathcal A_k\),
% \[
%     \mathbb P(D_{k,m}\in\mathcal C_k\mid\mathcal A_k)=p_\rho.
% \]
% Fix \(m\) and suppose that \(D_{k,m}\in\mathcal C_k\). Set
% \(
%     s_{k,m}:=\Delta_kD_{k,m}.
% \)
% By Lebourg's mean-value theorem, there exist
% \[
%     \bar x_{k,m}
%     =
%     X_k+\xi_{k,m}s_{k,m},
%     \qquad
%     \xi_{k,m}\in(0,1),
% \]
% and
% \[
%     \bar g_{k,m}\in\partial_C f(\bar x_{k,m})
% \]
% such that
% \[
%     f(X_k+s_{k,m})-f(X_k)
%     =
%     \langle \bar g_{k,m},s_{k,m}\rangle
%     =
%     \Delta_k\langle \bar g_{k,m},D_{k,m}\rangle .
% \]
% Since
% \(
%     \Delta_k\le\Delta_{\goldelta,\varepsilon}\le\goldelta,
% \)
% we have
% \(
%     \|\bar x_{k,m}-X_k\|
%     \le
%     \Delta_k
%     \le
%     \goldelta.
% \)
% Therefore
% \(
%     \bar x_{k,m}\in B_\goldelta(X_k),
% \)
% and hence
% \(
%     \bar g_{k,m}
%     \in
%     \partial_\goldelta f(X_k).
% \)
% By the projection theorem onto the closed convex set \(\partial_\goldelta f(X_k)\), we have
% \(
%     \langle \bar g_{k,m},g_{\goldelta,k}^*\rangle
%     \ge
%     \|g_{\goldelta,k}^*\|^2.
% \)
% Consequently,
% \[
%     \langle \bar g_{k,m},d_k^*\rangle
%     =
%     -\frac{
%         \langle \bar g_{k,m},g_{\goldelta,k}^*\rangle
%     }{
%         \|g_{\goldelta,k}^*\|
%     }
%     \le
%     -\|g_{\goldelta,k}^*\|
%     <
%     -\varepsilon.
% \]
% Now decompose \(D_{k,m}\) as
% \[
%     D_{k,m}
%     =
%     \rho_{k,m}d_k^*
%     +
%     \sqrt{1-\rho_{k,m}^2}\,v_{k,m},
% \]
% where
% \[
%     \rho_{k,m}:=\langle d_k^*,D_{k,m}\rangle,
%     \qquad
%     v_{k,m}\perp d_k^*,
%     \qquad
%     \|v_{k,m}\|=1.
% \]
% On the event \(D_{k,m}\in\mathcal C_k\), we have
% \[
%     \rho_{k,m}\ge\rho.
% \]
% Moreover, since \(f\) is \(L\)-Lipschitz on the relevant neighborhood, every
% Clarke subgradient appearing above has norm at most \(L\). Thus
% \[
%     \langle \bar g_{k,m},v_{k,m}\rangle
%     \le
%     \|\bar g_{k,m}\|
%     \le
%     L.
% \]
% Therefore
% \[
% \begin{aligned}
%     \langle \bar g_{k,m},D_{k,m}\rangle
%     &=
%     \rho_{k,m}\langle \bar g_{k,m},d_k^*\rangle
%     +
%     \sqrt{1-\rho_{k,m}^2}
%     \langle \bar g_{k,m},v_{k,m}\rangle                                      \\
%     &\le
%     -\rho_{k,m}\varepsilon
%     +
%     \sqrt{1-\rho_{k,m}^2}\,L                                                   \\
%     &\le
%     -\rho\varepsilon+\sqrt{1-\rho^2}\,L                                        \\
%     &=
%     -c_1 .
% \end{aligned}
% \]
% Hence
% \[
%     f(X_k+\Delta_kD_{k,m})-f(X_k)
%     \le
%     -c_1\Delta_k.
% \]
% On the event \(\mathcal E_{k,0}\cap\mathcal E_{k,m}\), we then have
% \[
% \begin{aligned}
%     \tilde f(X_k+\Delta_kD_{k,m})-\tilde f(X_k)
%     &\le
%     f(X_k+\Delta_kD_{k,m})-f(X_k)
%     +
%     2\varepsilon_f\Delta_k^{\Dexp}                                             \\
%     &\le
%     -c_1\Delta_k
%     +
%     2\varepsilon_f\Delta_k^{\Dexp}.
% \end{aligned}
% \]
% By the definition of \(\Delta_{\goldelta,\varepsilon}\),
% \[
%     \Delta_k^{\Dexp-1}
%     \le
%     \Delta_{\goldelta,\varepsilon}^{\Dexp-1}
%     <
%     \frac{c_1}{\theta+2\varepsilon_f}.
% \]
% Therefore
% \[
%     c_1\Delta_k
%     >
%     (\theta+2\varepsilon_f)\Delta_k^{\Dexp}.
% \]
% It follows that
% \[
%     \tilde f(X_k+\Delta_kD_{k,m})-\tilde f(X_k)
%     \le
%     -\theta\Delta_k^{\Dexp}.
% \]
% Thus the sufficient-decrease test succeeds for direction \(m\), and iteration
% \(k\) is successful.\\
% Now define
% \[
%     S_{k,m}
%     :=
%     \{D_{k,m}\in\mathcal C_k\}
%     \cap
%     \mathcal E_{k,m}.
% \]
% From the conditional uniformity of \(D_{k,m}\) and the conditional accuracy of
% the trial estimate,
% \[
%     \mathbb P(S_{k,m}\mid\mathcal A_k)\ge p_\rho\beta.
% \]
% By the conditional independence assumption, the events
% \(S_{k,1},\ldots,S_{k,\bar{m}}\) are conditionally independent given
% \(\mathcal A_k\). Hence
% \[
%     \mathbb P\left(
%         \bigcup_{m=1}^{\bar{m}} S_{k,m}
%         \,\middle|\,
%         \mathcal A_k
%     \right)
%     \ge
%     1-(1-p_\rho\beta)^{\bar{m}}.
% \]
% Furthermore, \(\mathcal E_{k,0}\) is conditionally independent of these trial
% events and satisfies
% \[
%     \mathbb P(\mathcal E_{k,0}\mid\mathcal A_k)\ge\beta.
% \]
% Therefore
% \[
% \begin{aligned}
%     \mathbb P(J_k=1\mid\mathcal A_k)
%     &\ge
%     \mathbb P\left(
%         \mathcal E_{k,0}
%         \cap
%         \bigcup_{m=1}^{\bar{m}} S_{k,m}
%         \,\middle|\,
%         \mathcal A_k
%     \right)                                                               \\
%     &\ge
%     \beta\left[1-(1-p_\rho\beta)^{\bar{m}}\right]                                  \\
%     &=
%     q_M.
% \end{aligned}
% \]
% Since \(W_{k+1}=+1\) if and only if \(J_k=1\), this also gives
% \[
%     \mathbb P(W_{k+1}=+1\mid\mathcal A_k)\ge q_M.
% \]
% This proves the claim.
% \end{proof}

%%%%%%%%%%%%%%%%%%%%%%%%%%%%%%%%%%%%%%%%%%%%%%%%%%%%%%%%%%%%%%%%%%%%%%
% I think there was an error in the following proof since I assumed conditionally independence of some events that may in fact be dependent on each other. That's why i replaced the proof with the one above

% \begin{lemma}[Probabilistic sufficient decrease with dense directions]
% \label{lem:random-cap}
% Fix a stationarity tolerance $\varepsilon>0$ and choose a cap parameter
% \[
%    \rho\in(0,1), 
%    \quad
%    p_\rho:=\tfrac12\,
%           I_{1-\rho^{2}}\!\bigl(\tfrac{n-1}{2},\tfrac12\bigr)\in(0,1),
% \]
% where $I_{\!\cdot}(a,b)$ is the regularised incomplete Beta function
% (cf.\ Remark~\ref{rem:reg-inc-beta}).
% Assume the \emph{large‑tolerance} inequality (cf.\ Remark~\ref{rem:rescale-Lipschitz})
% \[
%    c_1:=\rho\varepsilon-2L>0.
% \tag{A1}
% \]
% Set
% \[
%     \bar{\Delta}_{\varepsilon}
%    :=\min\Bigl\{
%         1,\;
%         \frac{c_1}{4\varepsilon_f},\;
%         \frac{c_1}{4\theta_{\rm req}}
%   \Bigr\}, 
%   \qquad
%   \theta:=\frac{c_1}{4\,\bar{\Delta}_{\varepsilon}},
%   \qquad
%   q_M:=\bigl[1-(1-p_\rho)^{\bar{m}}\bigr]\,\beta.
% \]
% Draw \emph{independently} $\bar{m}$ directions
% $d_{k,1},\dots,d_{k,\bar{m}}\sim\mathrm{Unif}(\mathbb S^{n-1})$ (cf.\ Remark~\ref{rem:\bar{m}-lower-bound}).
% For each $m=1,\dots,\bar{m}$ define
% \[
% G_{k,m} =
%   \Bigl\{|\tilde f(x_k)-f(x_k)| \le \varepsilon_f\Delta_k^{\Dexp} \Bigr\}
%   \cap
%   \Bigl\{|\tilde f(x_k+\Delta_k d_{k,m})-f(x_k+\Delta_k d_{k,m})|
%           \le \varepsilon_f\Delta_k^{\Dexp}\Bigr\},
% \]
% and the indicator \(J_{k,m}=\ind_{G_{k,m}}\).
% Assume the random estimates are $\beta$‑probabilistically
% $\varepsilon_f$‑accurate, i.e.
% \begin{equation}\label{eq:Gkm_rand}
%    \Pb\bigl(J_{k,m}=1 \;|\; \F_{k-1}\bigr)\;\ge\;\beta\qquad\forall m.
% \end{equation}
% Then, for every \(k<T_\varepsilon\) with \(\Delta_k\le\bar{\Delta}_{\varepsilon}\),
% \[
%    \Pb\bigl(\text{iteration $k$ is successful}\mid \F_{k-1}\bigr)
%    \;\ge\;
%    q_M
%    \;>\;0.
% \]
% \end{lemma}

% \begin{proof}
% Because $k<T_\varepsilon$, there exists $g_k\in\partial_C f(x_k)$ with
% $\|g_k\|>\varepsilon$.  For each trial \(m\) define the cap
% \[
%    \mathcal C_{k,m}
%    :=\Bigl\{ d\in\mathbb{S}^{n-1} :
%              -\langle g_k/\|g_k\|,d\rangle\ge\rho \Bigr\}.
% \]
% The conditional probability of sampling $d_{k,m}$ in
% \(\mathcal C_{k,m}\) is exactly the cap’s surface fraction,
% \[
%    \Pb(d_{k,m}\in\mathcal C_{k,m}\mid\F_{k-1})=p_\rho.
% \]
% Fix an $m$ and set $s_{k,m}:=\Delta_k d_{k,m}$.
% By the Lebourg mean‑value theorem there are
% $\bar x_{k,m}=x_k+\xi_{k,m}s_{k,m}$ and
% $\bar g_{k,m}\in\partial_C f(\bar x_{k,m})$ such that
% \[
%    f(x_k+s_{k,m})-f(x_k)=\langle\bar g_{k,m},s_{k,m}\rangle.
% \]
% Split $\bar g_{k,m}=g_k+(\bar g_{k,m}-g_k)$ and use:
% (i) the cap condition
% $-\langle g_k,d_{k,m}\rangle\ge\rho\|g_k\|\ge\rho\varepsilon$,
% (ii) $\|\bar g_{k,m}-g_k\|\le\|\bar g_{k,m}\|+\|g_k\|\le2L$,
% (iii) $\|s_{k,m}\|=\Delta_k$.  Then
% \[
%   f(x_k+s_{k,m})-f(x_k)
%   \le -\rho\varepsilon\,\Delta_k + 2L\,\Delta_k.
% \]
% On the good‑estimate event \(G_{k,m}\),
% \[
%    \tilde f(x_k+s_{k,m})-\tilde f(x_k)
%    \le -c_{1}\Delta_k + 2\varepsilon_f\Delta_k^{\Dexp}.
% \]
% Since $\Delta_k\le\bar{\Delta}_{\varepsilon} \le c_{1}/(4\varepsilon_f)$,
% \[
%    2\varepsilon_f\Delta_k^{\Dexp}\;\le\;\frac{c_{1}}{2}\,\Delta_k,
% \]
% hence
% \[
%    \tilde f(x_k+s_{k,m})-\tilde f(x_k)
%    \le -\frac{c_{1}}{2}\,\Delta_k
%    =   -\theta\,\Delta_k^{\Dexp},
% \]
% so the quadratic sufficient‑decrease test is satisfied and
% \emph{that} trial would render iteration~\(k\) successful.
% The events \(\mathcal C_{k,m}\) are i.i.d.\ with probability \(p_\rho\),
% and by~\eqref{eq:Gkm_rand} each \(G_{k,m}\) is conditionally
% independent of \(\mathcal C_{k,m}\) with probability at least
% \(\beta\).
% Thus, for a single trial
% \[
%    \Pb\bigl(\mathcal C_{k,m}\cap G_{k,m}\mid\F_{k-1}\bigr)
%    \;\ge\; p_\rho\,\beta.
% \]
% Independence of the \(\bar{m}\) trials implies that the probability
% \emph{no} trial is simultaneously in the cap and has good estimates
% is at most \(\bigl(1-p_\rho\beta\bigr)^{\bar{m}}\).
% Therefore
% \[
%    \Pb(\text{iteration $k$ successful}\mid\F_{k-1})
%    \;\ge\;
%    1-\bigl(1-p_\rho\beta\bigr)^{\bar{m}}
%    \ge\bigl[1-(1-p_\rho)^{\bar{m}}\bigr]\beta
%    =q_M.
% \]
% \end{proof}
%--------------------------------------------------------------------






\begin{remark}[Overcoming the coarse Clarke bound via Goldstein subdifferentials]%
\label{rem:rescale-Lipschitz}
In a pure Clarke framework, bounding the Lebourg subgradient relies on the coarse Lipschitz bound $\|\bar G_{k,m}-G_k\|\le 2L$. This requires a descent margin of $\rho\varepsilon-2L>0$, forcing $\rho > 2L/\varepsilon$. As $\varepsilon \to 0$, this constraint inevitably demands $\rho \ge 1$. The spherical cap becomes empty, collapsing the success probability $p_\rho$ to zero, and restricting the analysis to low-accuracy targets where $\varepsilon > 2L$. Shifting to the Goldstein $\goldelta$-subdifferential resolves this topological barrier. By targeting an $(\varepsilon, \goldelta)$-Goldstein stationary point and restricting $\Delta_k \le \goldelta$, the intermediate Lebourg point is trapped in $B_\goldelta(X_k)$, ensuring $\bar G_{k,m} \in \partial_\goldelta f(X_k)$. The projection theorem then eliminates the $2L$ residual, refining the descent condition to $c_1 := \rho\varepsilon - \sqrt{1-\rho^2}L > 0$. This yields $\rho > \frac{L}{\sqrt{\varepsilon^2 + L^2}}$, which is strictly less than $1$ for all $\varepsilon > 0$, guaranteeing a positive success probability as $\varepsilon \to 0$. In the smoother regime \(f\in C^{1,1}\), define 
\(
    G_k:=\nabla f(X_k),
    \;
    \bar G_{k,m}:=\nabla f(\bar X_{k,m}),
\)
where \(\bar X_{k,m}\) is the intermediate point on the segment
\([X_k,X_k+\Delta_kD_{k,m}]\). If \(\nabla f\) is
\(\bar L\)-Lipschitz, then
\[
    \left|
    \langle \bar G_{k,m}-G_k,D_{k,m}\rangle
    \right|
    \le
    \bar L\Delta_k.
\]
% In smoother regimes ($f\in C^{1,1}$ with an $\bar L$-Lipschitz gradient), $\partial_C f(X)=\{\nabla f(X)\}$ and
% \[
% \langle \bar G_{k,m}-G_k,\ D_{k,m}\rangle
% \le \bar L\,\Delta_k.
% \] \sara{why $-g_k$?}
The $2L$ penalty is replaced by $\mathcal{O}(\Delta_k)$, yielding $c_{1,k}=\rho\varepsilon-\mathcal{O}(\Delta_k)$. Since $\Delta_k\to 0$ almost surely, $c_{1,k}>0$ eventually holds regardless of $\varepsilon$, as established in~\cite{MR4359469}.
\end{remark}
% \begin{remark}[Overcoming the coarse Clarke bound via Goldstein subdifferentials]%
% \label{rem:rescale-Lipschitz}
% In a pure Clarke framework, bounding the Lebourg mean-value subgradient $\bar g_{k,m}$ against the minimum-norm element $g_k \in \partial_C f(x_k)$ relies on the coarse Lipschitz bound $\|\bar g_{k,m}-g_k\|\le \|\bar g_{k,m}\|+\|g_k\|\le 2L$. This dictates a required descent margin of $\rho\varepsilon-2L>0$, imposing the strict geometric constraint $\rho > 2L/\varepsilon$. Consequently, driving the stationarity tolerance $\varepsilon \to 0$ inevitably enters the regime where $\varepsilon \le 2L$, forcing $\rho$ to hard-threshold to $1$. This in turn leads to \(\mathbb P(U\leq 1-\rho^2) = \mathbb P(U \leq 0) = 0\). The spherical cap becomes mathematically empty, causing the success probability $p_\rho$ to hard-collapse to zero. This fundamentally restricts the standard Clarke analysis to low-accuracy targets where $\varepsilon > 2L$.\\
% Shifting to the Goldstein $\delta$-subdifferential circumvents this topological barrier. By defining the target as an $(\varepsilon, \delta)$-Goldstein stationary point and restricting the step size to $\Delta_k \le \delta$, the intermediate Lebourg evaluation point is confined within $B_\delta(x_k)$. Thus, $\bar g_{k,m} \in \partial_\delta^G f(x_k)$ by definition. This allows us to apply the projection theorem across the entire neighborhood, effectively eliminating the $2L$ residual. The descent condition simplifies to $c_1 := \rho\varepsilon - \sqrt{1-\rho^2}L > 0$. Solving for $\rho$ yields the lower bound $\rho > \frac{L}{\sqrt{\varepsilon^2 + L^2}}$, which remains strictly less than $1$ for all $\varepsilon > 0$, preserving a strictly positive spherical cap probability as $\varepsilon \to 0$.\\
% In smoother regimes, this geometric restriction can be removed entirely. If $f\in C^{1,1}$ with an $\bar L$-Lipschitz gradient, then $\partial_C f(x)=\{\nabla f(x)\}$ and
% \[
% \langle \bar g_{k,m}-g_k,\ d_{k,m}\rangle
% =\langle \nabla f(x_k+\xi_{k,m}\Delta_k d_{k,m})-\nabla f(x_k),\ d_{k,m}\rangle
% \le \bar L\,\Delta_k.
% \]
% Thus, one obtains $c_{1,k}=\rho\varepsilon-\mathcal{O}(\Delta_k)$.
% Since $\Delta_k\to 0$ almost surely under the usual step-size updates, $c_{1,k}>0$ holds eventually regardless of $\varepsilon$, as in the expected complexity analysis given in~\cite{MR4359469} for stochastic direct-search schemes that handle smooth objective functions.
% \end{remark}
% \begin{remark}[Ensuring $c_{1}>0$ is not restrictive]%
% \label{rem:rescale-Lipschitz}
% The requirement
% \(
%    c_{1}:=\rho\varepsilon-2L>0
% \)
% can always be satisfied. Indeed, this is \emph{only a scaling issue}: we can always multiply the objective by a positive constant and make $L$ as small/large as we like without altering the optimisation landscape.
% \end{remark}
% \begin{remark}[Polling capacity \(\bar{m}\) and dimensional complexity]
% \label{rem:\bar{m}-lower-bound}
% To apply the expected complexity bound of Theorem \ref{t:expected complexity}, we must ensure the per-iteration success probability satisfies \(q_M>\tfrac12\). Assuming the estimator-accuracy probability satisfies \(\beta>\tfrac12\), we can explicitly solve \(q_M = \beta\,[\,1-(1-p_\rho\beta)^{\bar{m}}\,] > \tfrac12\) for \(\bar{m}\) to obtain the lower bound:
% $$\bar{m}\;>\; M_{\min}\;:=\; \left\lceil \frac{\displaystyle \log\!\bigl(1-\tfrac1{2\beta}\bigr)}{\displaystyle \log(1-p_\rho\beta)} \right\rceil.$$
% Thus, choosing \(\bar{m} > M_{\min}\) polling directions guarantees \(q_M>\tfrac12\), satisfying the supermartingale descent condition.\\
% However, targeting the Goldstein $\goldelta$-subdifferential requires the descent margin $c_1 > 0$, forcing $\rho > L / \sqrt{\varepsilon^2 + L^2}$. As the stationarity tolerance $\varepsilon \downarrow 0$, the required cap alignment $\rho \uparrow 1$.\\
% To quantify the asymptotic impact on $M_{\min}$, let $z := 1-\rho^2 = \mathcal O(\varepsilon^2)$. Using the integral representation of the regularized incomplete Beta function (cf.\ Definition~\ref{def:spherical-cap-probability}), the leading-order approximation of the cap probability is:
% $$p_\rho = \frac{1}{(n-1)B\bigl(\frac{n-1}{2}, \frac12\bigr)} z^{\frac{n-1}{2}} + o\bigl(z^{\frac{n-1}{2}}\bigr) = \mathcal O\left(\left(\varepsilon/L\right)^{n-1}\right).$$
% Because $\log(1-x) \approx -x$ for infinitesimally small $x$, the denominator in the $M_{\min}$ bound scales as $\log(1-p_\rho\beta) \approx -p_\rho\beta$. Therefore, the required number of polling directions scales asymptotically as:
% $$\bar{m}(\varepsilon) = \mathcal O\left(\frac{1}{p_\rho}\right) = \mathcal O\bigl(L^{n-1}\cdot\varepsilon^{1-n}\bigr).$$
% This exponential dependence on the dimension $n$ quantifies the fundamental topological difficulty of finding a descent direction in a purely nonsmooth landscape via dense random search.
% \end{remark}
We next quantify the probability of the good event introduced above. More specifically, the following lemma shows that,
under conditional independence and accuracy assumptions, this event occurs
with probability bounded from below uniformly in \(k\).



\begin{lemma}[Probability of the good event]
\label{lem:good-event-probability}
Assume that, conditionally on \(\mathcal A_k\), the directions
\(D_{k,1},\ldots,D_{k,\bar{m}}\) are independent and uniformly distributed on
\(\mathbb S^{n-1}\). Assume also that the estimates are
\(\beta\)-probabilistically \(\varepsilon_f\)-accurate, namely
\[
    \mathbb P(\mathcal E_{k,0}\mid\mathcal A_k)\ge\beta,
    \qquad
    \mathbb P(\mathcal E_{k,m}\mid\mathcal A_k,D_{k,m})\ge\beta,
    \quad m=1,\ldots,\bar{m}.
\]
Moreover, assume that, conditionally on \(\mathcal A_k\), the base-estimate
event \(\mathcal E_{k,0}\) is independent of the trial events, and that the
pairs
\[
    \{(D_{k,m},\mathcal E_{k,m})\}_{m=1}^{\bar{m}}
\]
are conditionally independent.\\
Then,
we have
\[
    \mathbb P(\mathcal I_k\mid\mathcal A_k)
    =
    \mathbb P(W_{k+1}=+1\mid\mathcal A_k)
    \ge
    q_{\bar{m}}\,
\]
where
\[
    q_{\bar{m}}
    :=
    \beta\left[1-(1-p_\rho\beta)^{\bar{m}}\right].
\]
\end{lemma}

\begin{proof}
Since \(D_{k,m}\) is uniformly distributed on
\(\mathbb S^{n-1}\) conditionally on \(\mathcal A_k\), and since the spherical
cap has probability \(p_\rho\), we have
\[
    \mathbb P(D_{k,m}\in\mathcal C_\rho(D_k^*)\mid\mathcal A_k)=p_\rho .
\]
Define  $\mathcal S_{k,m}
    :=
    \{D_{k,m}\in\mathcal C_\rho(D_k^*)\}
    \cap
    \mathcal E_{k,m}. $
Using the conditional uniformity of \(D_{k,m}\) and the conditional
\(\beta\)-accuracy of the trial estimate,
\[
\begin{aligned}
    \mathbb P(\mathcal S_{k,m}\mid\mathcal A_k)
    &=
    \mathbb E\!\left[
        \mathbf 1_{\{D_{k,m}\in\mathcal C_\rho(D_k^*)\}}
        \mathbb P(\mathcal E_{k,m}\mid\mathcal A_k,D_{k,m})
        \,\middle|\,
        \mathcal A_k
    \right]                                      \\
    &\ge
    \beta\,
    \mathbb P(D_{k,m}\in\mathcal C_\rho(D_k^*)\mid\mathcal A_k)              \\
    &=
    p_\rho\beta .
\end{aligned}
\]
By conditional independence, the events
\(\mathcal S_{k,1},\ldots,\mathcal S_{k,\bar{m}}\) are conditionally independent
given \(\mathcal A_k\). Hence
\[
    \mathbb P\left(
        \bigcup_{m=1}^{\bar{m}} \mathcal S_{k,m}
        \,\middle|\,
        \mathcal A_k
    \right)
    \ge
    1-(1-p_\rho\beta)^{\bar{m}}.
\]
Furthermore, \(\mathcal E_{k,0}\) is conditionally independent of these trial
events and satisfies
\[
    \mathbb P(\mathcal E_{k,0}\mid\mathcal A_k)\ge\beta.
\]
Therefore,
\[
\begin{aligned}
    \mathbb P(\mathcal I_k\mid\mathcal A_k)
    &=
    \mathbb P\left(
        \mathcal E_{k,0}
        \cap
        \bigcup_{m=1}^{\bar{m}} \mathcal S_{k,m}
        \,\middle|\,
        \mathcal A_k
    \right)                                                               \\
    &\ge
    \beta\left[1-(1-p_\rho\beta)^{\bar{m}}\right]                                  \\
    &=
    q_{\bar{m}}.
\end{aligned}
\]
Since \(W_{k+1}=+1\) if and only if \(\mathcal I_k\) occurs, we obtain
\[
    \mathbb P(W_{k+1}=+1\mid\mathcal A_k)
    =
    \mathbb P(\mathcal I_k\mid\mathcal A_k)
    \ge q_{\bar{m}},
\]
which gives the claim.
\end{proof}



\begin{remark}[Polling capacity \(\bar{m}\) and dimensional complexity]
\label{rem:M-lower-bound}
To apply the expected complexity bound of Theorem~\ref{t:expected complexity}
with constants bounded uniformly, it is convenient to fix a number
\(q_0\in(1/2,\beta)\) and choose \(\bar{m}\) so that
\[
    q_{\bar{m}}
    :=
    \beta\,[\,1-(1-p_\rho\beta)^{\bar{m}}\,]
    \ge q_0 .
\]
This is possible only if \(\beta>1/2\). Solving the inequality
\(q_{\bar{m}}\ge q_0\) gives
\[
    (1-p_\rho\beta)^{\bar{m}}
    \le
    1-\frac{q_0}{\beta},
\]
and hence
\[
    \bar{m}
    \ge
    \frac{\log(1-q_0/\beta)}
         {\log(1-p_\rho\beta)}.
\]
Thus, for example, it is sufficient to choose
\[
    \bar{m}
    >
    \bar{m}_{\min}(q_0)
    :=
    \left\lceil
    \frac{\log(1-q_0/\beta)}
         {\log(1-p_\rho\beta)}
    \right\rceil .
\]
With this choice, \(q_{\bar{m}}\ge q_0>1/2\), satisfying the supermartingale condition. However, targeting the Goldstein $\goldelta$-subdifferential requires the descent margin $c_1 > 0$, forcing $\rho > L / \sqrt{\varepsilon^2 + L^2}$.
As the stationarity tolerance \(\varepsilon\downarrow0\), the required cap
alignment satisfies \(\rho\uparrow1\), and the cap probability \(p_\rho\)
goes to zero. To quantify the resulting dependence of \(\bar{m}\) on \(\varepsilon\), choose
\(\rho\) so that
\[
    1-\rho^2\asymp \left(\frac{\varepsilon}{L}\right)^2 .
\]
Equivalently, the cap remains nonempty but becomes narrower as
\(\varepsilon\downarrow0\). Let \(z:=1-\rho^2\). Using the small-\(z\)
asymptotics of the regularized incomplete Beta function from Definition~\ref{def:reg-inc-beta-funct}, we have
\[
    p_\rho
    =
    \frac12
    I_z\left(\frac{n-1}{2},\frac12\right)
    \asymp
    z^{(n-1)/2}.
\]
Therefore
\[
    p_\rho
    \asymp
    \left(\frac{\varepsilon}{L}\right)^{n-1}.
\]
Since $\log(1-p_\rho\beta)\asymp -p_\rho\beta
    $ as $p_\rho\downarrow0, $
the sufficient polling budget satisfies
\[
    \bar{m}_{\min}(q_0)
    =
    \mathcal O\left(\frac{1}{p_\rho}\right)
    =
    \mathcal O\left(
        \left(\frac{L}{\varepsilon}\right)^{n-1}
    \right).
\]
Thus, for fixed \(L\), choosing
\[
    \bar{m}=\mathcal O(\varepsilon^{1-n})
\]
is sufficient to keep the success probability uniformly bounded away from
\(1/2\). This dimensional dependence reflects the geometric difficulty of
finding a descent direction in a purely nonsmooth landscape via a dense set of random directions.
The above argument, together with the use of the spherical cap, is somehow related to the probabilistic-descent framework of
Gratton et al.~\cite[Appendix~B]{GrattonRoyerVicenteZhang2015}. There, for
smooth objectives, random polling directions are used to obtain, with high
probability, a direction lying in a spherical cap around the negative gradient;
and a fixed positive cosine threshold  (for
instance of order \(1/\sqrt n\)) is sufficient to guarantee descent. In our nonsmooth
Goldstein setting, the spherical-cap idea is used, but nonsmoothness lets the cap shrink with the stationarity
tolerance, leading to the polling budget
\(\bar m=O(\varepsilon^{1-n})\). We finally note that the constant
\(c_{\max}=2\mu_h \bar{m}(\Imax+1)\) depends on \(\bar{m}\). We can however
guarantee that \(c_{\max}\) remains uniformly bounded by tightening of
the sample-average accuracy, that is by increasing the batch size in the stochastic estimates construction, if needed.
\end{remark}
% \begin{remark}[Picking \(\bar{m}\)]
% \label{rem:\bar{m}-lower-bound}
% To apply the renewal-reward complexity bound of Theorem \ref{t:expected complexity} we need \(q_M>\tfrac12\). Assume the estimator-accuracy probability satisfies \(\beta>\tfrac12\). Then we can derive the explicit
% lower bound
% \[
%    \bar{m}\;>\;
%    M_{\min}\;:=\;
%    \left\lceil
%       \frac{\displaystyle
%             \log\!\bigl(1-\tfrac1{2\beta}\bigr)}
%            {\displaystyle
%             \log(1-p_\rho\beta)}
%    \right\rceil .
% \]
% Thus choosing \(\bar{m}>M_{\min}\) polling directions per iteration guarantees
% \(q_M>\tfrac12\) and therefore the super‑martingale argument of
% Theorem \ref{t:expected complexity} applies.
% \end{remark}
% \begin{remark}[Geometric sampling, truncation, and choice of the cap $\bar{m}$]
% \label{rem:geom-trunc}
% Within an outer iteration, directions are drawn i.i.d.\ until the first that passes the
% sufficient-decrease test, with at most $\bar{m}$ trials (truncated geometric).
% Let $q:=p_\rho$ be the spherical-cap probability (Remark \ref{rem:reg-inc-beta}), and let
% $\beta\in(0,1)$ be the accuracy level for a single estimate. If both the baseline and the trial
% estimates are $\beta$-accurate, then the per-trial success probability satisfies
% \[
% p \ \ge\ q\,\beta \quad\text{(conditional on a $\beta$-accurate baseline)},\qquad
% \text{or}\qquad p \ \ge\ q\,\beta^2 \ \text{ unconditionally.}
% \]
% For the truncated geometric with parameter $p$:
% \[
% \PP(N_k>\bar{m} \mid \F_{k-1}) \;=\; (1-p)^{\bar{m}}, 
% \qquad 
% \EE[N_k \mid \F_{k-1}] \;\le\; \min\!\Big\{\frac{1}{p},\,\bar{m}\Big\}.
% \]
% Hence a uniform one-iteration success bound is
% \[
% \PP(\text{success at iteration }k \mid \F_{k-1}) \;\ge\; 1-(1-p)^{\bar{m}} 
% \ \ge\ 1-(1-q\beta)^{\bar{m}}
% \quad\big(\text{or } 1-(1-q\beta^2)^{\bar{m}} \text{ uncond.}\big).
% \]
% Given a target per-iteration failure level $\delta\in(0,1)$, choose
% \[
% \bar{m} \ \ge\ \left\lceil \frac{\log \delta}{\log(1-q\beta)} \right\rceil
% \quad \bigg(\text{or }\ \left\lceil \frac{\log \delta}{\log(1-q\beta^2)} \right\rceil \text{ uncond.}\bigg).
% \]
% With this choice of $\bar{m}$, we have
% \[
% \PP(\text{no success in }\bar{m}\text{ trials}\mid\F_{k-1})=(1-p)^{\bar{m}}
% \ \le\ (1-q\beta)^{\bar{m}}\ \le\ \delta,
% \]
% hence
% \[
% \PP(\text{success at iteration }k\mid\F_{k-1})\ \ge\ 1-\delta
% \qquad\text{for all }k.
% \]
% \end{remark}
% \begin{remark}[Reconciling the $\theta$ bounds and the Goldstein radius]
% \label{rem:theta-two-bounds} 
% From the spherical-cap lemma (with $\beta$-accurate estimates), converting the linear decrease into the $\Dexp$-test requires:
% \[
% \Delta_{\goldelta,\varepsilon}\;<\;\min\left\{\goldelta,\ \left(\frac{c_1}{\,\eta_{\rm req}+2\varepsilon_f\,}\right)^\frac{1}{\Dexp-1}\right\},
% \qquad
% \theta\ \in\ \Big(\,\eta_{\rm req},\ \frac{c_1}{\Delta_{\goldelta,\varepsilon}^{\Dexp-1}}-2\varepsilon_f\,\Big),
% \]
% where \(\eta_{\rm req}:=\tfrac{2\mu_h}{1-\gamma^{\Dexp}}\) is the lower bound on the merit parameter required by the drift. These inequalities are mathematically compatible and easy to enforce: pick any margin \(\mu>0\) and define the critical threshold as
% \[
% \Delta_{\goldelta,\varepsilon}\ :=\ \min\!\left\{\goldelta,\ \Delta_{\max},\ \left(\frac{c_1}{\,\eta_{\rm req}+2\varepsilon_f+\mu\,}\right)^\frac{1}{\Dexp-1}\right\}.
% \]
% Because \(\Delta_{\goldelta,\varepsilon} \le \left(\frac{c_1}{\eta_{\rm req}+2\varepsilon_f+\mu}\right)^{\frac{1}{\Dexp-1}}\), it directly follows that \(\Delta_{\goldelta,\varepsilon}^{\Dexp-1} \le \frac{c_1}{\eta_{\rm req}+2\varepsilon_f+\mu}\), which rearranges exactly to
% \[
% \frac{c_1}{\Delta_{\goldelta,\varepsilon}^{\Dexp-1}}-2\varepsilon_f\ \ge\ \eta_{\rm req}+\mu.
% \]
% Thus, the interval \(\big(\eta_{\rm req},\,\tfrac{c_1}{\Delta_{\goldelta,\varepsilon}^{\Dexp-1}}-2\varepsilon_f\big)\) is strictly nonempty, ensuring a valid \(\theta\) always exists. 
% \end{remark}
% Recall that an iteration is called \emph{successful} if at least one of
% the \(\bar{m}\) tested directions satisfies the sufficient-decrease test, and
% \emph{unsuccessful} otherwise. Equivalently,
% \[
%     J_k=1
%     \quad\Longleftrightarrow\quad
%     H_k\ge0,
% \]
% whereas \(J_k=0\) corresponds to \(H_k=-1\). By the stepsize update rule, an
% unsuccessful iteration contracts the polling radius by \(\gamma\), while a
% successful iteration expands the next polling radius by at least
% \(\gamma^{-1}\). 
Recall that an iteration is called \emph{successful} if at least one of
the \(\bar{m}\) tested directions satisfies the sufficient-decrease test, and
\emph{unsuccessful} otherwise. Equivalently, success corresponds to
\(H_k\ge 0\), whereas failure corresponds to \(H_k=-1\). By the stepsize
update rule, an unsuccessful iteration contracts the polling radius by
\(\gamma\), while a successful iteration expands the next polling radius by at
least \(\gamma^{-1}\). The good-event indicator used in the stopping-time argument is instead
\(J_k=\mathbf 1_{\mathcal I_k}\). Hence \(J_k=1\) does not define success in
general; rather, on the event
\(\{T_{r,\varepsilon}>k\}\cap\{\Delta_k\le\Delta_{r,\varepsilon}\}\),
Lemma~\ref{lem:good-event-sufficient-decrease} shows that
\(\mathcal I_k\) implies \(H_k\ge0\). The following lemma verifies the radius-dynamics condition in
Assumption~\ref{ass:5}.
% \begin{lemma}[Assumption (ii) holds with $q=q_M$ for DSE]
% \label{lem:assump-ii-DSE-final}
% Let \(\gamma\in(0,1)\) and set
% \[
%     \lambda:=\log(\gamma^{-1})=-\log\gamma>0.
% \]
% Let \(\bar\Delta_{\goldelta,\varepsilon}\in
% (0,\Delta_{\goldelta,\varepsilon}]\), where
% \(\Delta_{\goldelta,\varepsilon}\) is the threshold from
% Lemma~\ref{lem:good-event-sufficient-decrease}. Then, for every \(k\ge0\),
% \begin{equation}\label{eq:radius-recursion}
%     \mathbf 1_{\{T_{\goldelta,\varepsilon}>k\}}\Delta_{k+1}
%     \ge
%     \mathbf 1_{\{T_{\goldelta,\varepsilon}>k\}}
%     \min\left\{
%         \Delta_k e^{\lambda W_{k+1}},
%         \bar\Delta_{\goldelta,\varepsilon}
%     \right\}.
% \end{equation}
% Moreover, whenever
% \[
%     T_{\goldelta,\varepsilon}>k
%     \qquad\text{and}\qquad
%     \Delta_k\le\bar\Delta_{\goldelta,\varepsilon}
%     \le
%     \Delta_{\goldelta,\varepsilon},
% \]
% we have
% \begin{equation}\label{eq:success-prob}
%     \mathbb P(W_{k+1}=+1\mid\mathcal A_k)
%     \ge
%     q_M
%     :=
%     \beta\left[1-(1-p_\rho\beta)^{\bar{m}}\right].
% \end{equation}
% In particular, if \(\bar{m}\) is chosen so that \(q_M>1/2\), then
% Assumption~\ref{ass:5}(ii) holds with \(q=q_M\), threshold
% \(\bar\Delta_{\goldelta,\varepsilon}\), and
% \(\lambda=\log(\gamma^{-1})\).
% \end{lemma}
% \begin{proof}
% First, consider the radius recursion. If iteration \(k\) is unsuccessful, then
% \(J_k=0\), \(W_{k+1}=-1\), and the algorithm sets
% \[
%     \Delta_{k+1}
%     =
%     \gamma\Delta_k
%     =
%     \Delta_k e^{-\lambda}
%     =
%     \Delta_k e^{\lambda W_{k+1}}.
% \]
% If iteration \(k\) is successful, then \(J_k=1\), \(W_{k+1}=+1\), and by the stepsize update rule
% \[
%     \Delta_{k+1}
%     =
%     \begin{cases}
%         \gamma^{-1}\Delta_k, & H_k=0,\\[1mm]
%         \gamma^{-H_k}\Delta_k, & H_k\ge1.
%     \end{cases}
% \]
% Therefore, in both successful cases,
% \[
%     \Delta_{k+1}
%     \ge
%     \gamma^{-1}\Delta_k
%     =
%     \Delta_k e^\lambda
%     =
%     \Delta_k e^{\lambda W_{k+1}}.
% \]
% Hence, in all cases,
% \[
%     \Delta_{k+1}
%     \ge
%     \Delta_k e^{\lambda W_{k+1}}.
% \]
% Consequently,
% \[
%     \Delta_{k+1}
%     \ge
%     \min\left\{
%         \Delta_k e^{\lambda W_{k+1}},
%         \bar\Delta_{\goldelta,\varepsilon}
%     \right\}.
% \]
% Multiplying by
% \(\mathbf 1_{\{T_{\goldelta,\varepsilon}>k\}}\) gives
% \eqref{eq:radius-recursion}.
% Now suppose that
% \[
%     T_{\goldelta,\varepsilon}>k
%     \qquad\text{and}\qquad
%     \Delta_k\le\bar\Delta_{\goldelta,\varepsilon}
%     \le
%     \Delta_{\goldelta,\varepsilon}.
% \]
% By Lemma~\ref{lem:random-cap},
% \[
%     \mathbb P(J_k=1\mid\mathcal A_k)
%     \ge
%     \beta\left[1-(1-p_\rho\beta)^{\bar{m}}\right]
%     =
%     q_M.
% \]
% Since \(W_{k+1}=+1\) if and only if \(J_k=1\), this gives
% \[
%     \mathbb P(W_{k+1}=+1\mid\mathcal A_k)
%     \ge
%     q_M.
% \]
% If \(\bar{m}\) is chosen so that \(q_M>1/2\), then the success-probability condition
% in Assumption~\ref{ass:5}(ii) holds. This concludes the proof.
% \end{proof}

\begin{lemma}[Verification of Assumption~\ref{ass:5}(ii)]
\label{lem:assumption-ii-holds}
Let \(\gamma\in(0,1)\) and set
$ \lambda:=-\log\gamma>0.$
Let
\[
    \bar\Delta_{\goldelta,\varepsilon}
    =
    \Delta_0\gamma^{j_{\goldelta,\varepsilon}},
    \qquad
    j_{\goldelta,\varepsilon}\in\mathbb Z,\qquad j_{r,\varepsilon} \leq 0,
\]
be such that
\[
    \bar\Delta_{\goldelta,\varepsilon}
    \le
    \Delta_{\goldelta,\varepsilon},
\]
where \(\Delta_{\goldelta,\varepsilon}\) is defined in
Lemma~\ref{lem:good-event-sufficient-decrease}. Then Assumption~\ref{ass:5}(ii) holds with
\[
    q=q_{\bar{m}}
    :=
    \beta\left[1-(1-p_\rho\beta)^{\bar{m}}\right].
\]
\end{lemma}

\begin{proof}
The inequality in Assumption~\ref{ass:5}(ii) is trivial when
\(\mathbf 1_{\{T_{\goldelta,\varepsilon}>k\}}=0\). Hence, assume that
\(T_{\goldelta,\varepsilon}>k\). Since the stepsizes lie on the geometric mesh, we can write
\[
    \Delta_k=\Delta_0\gamma^{i_k}
    \qquad\text{for some } i_k\in\mathbb Z.
\]
If \(\Delta_k>\bar\Delta_{\goldelta,\varepsilon}\), then
\[
    \Delta_k\ge \gamma^{-1}\bar\Delta_{\goldelta,\varepsilon}.
\]
Moreover, from the stepsize update rule, we always have
\[
    \Delta_{k+1}\ge \gamma\Delta_k.
\]
Thus
\[
    \Delta_{k+1}
    \ge
    \gamma\Delta_k
    \ge
    \bar\Delta_{\goldelta,\varepsilon},
\]
and therefore
\[
    \Delta_{k+1}
    \ge
    \min\left\{
        \Delta_k e^{\lambda W_{k+1}},
        \bar\Delta_{\goldelta,\varepsilon}
    \right\}.
\]
Now assume that
\[
    \Delta_k\le\bar\Delta_{\goldelta,\varepsilon}
    \le
    \Delta_{\goldelta,\varepsilon}.
\]
If \(W_{k+1}=+1\), then \(\mathcal I_k\) occurs. By
Lemma~\ref{lem:good-event-sufficient-decrease}, iteration \(k\) is successful.
Hence
\[
    \Delta_{k+1}
    \ge
    \gamma^{-1}\Delta_k
    =
    \Delta_k e^\lambda
    =
    \Delta_k e^{\lambda W_{k+1}}.
\]
If \(W_{k+1}=-1\), then, independently of whether the iteration is successful
or unsuccessful, the update rule gives
\[
    \Delta_{k+1}
    \ge
    \gamma\Delta_k
    =
    \Delta_k e^{-\lambda}
    =
    \Delta_k e^{\lambda W_{k+1}}.
\]
Therefore, in both cases,
\[
    \Delta_{k+1}
    \ge
    \min\left\{
        \Delta_k e^{\lambda W_{k+1}},
        \bar\Delta_{\goldelta,\varepsilon}
    \right\}.
\]
Finally, by Lemma~\ref{lem:good-event-probability},
\[
    \mathbb P(W_{k+1}=+1\mid\mathcal A_k)
    =
    \mathbb P(\mathcal I_k\mid\mathcal A_k)
    \ge
    \beta\left[1-(1-p_\rho\beta)^{\bar{m}}\right]
    =
    q_{\bar{m}}.
\]
If \(\bar{m}\) is chosen so that \(q_{\bar{m}}>1/2\), then
Assumption~\ref{ass:5}(ii) holds with \(q=q_{\bar{m}}\).
\end{proof}



\begin{remark}[Choice of $\bar{\Delta}_{r,\varepsilon}$ and mesh index $j_{r,\varepsilon}$]\label{rem: choice of delta'}
Lemma~\ref{lem:good-event-sufficient-decrease} and Lemma~\ref{lem:good-event-probability} provide an analytic threshold
\(\Delta_{\goldelta,\varepsilon}\) such that, for
\(k<T_{\goldelta,\varepsilon}\), the spherical-cap argument yields a uniformly
positive probability of success whenever
\(
    \Delta_k\le\Delta_{\goldelta,\varepsilon}.
\)
For the stopping-time theorem, we need a threshold belonging to
the geometric mesh generated by the radius updates. Since the updates multiply
the radius by powers of \(\gamma\), the radii belong to the mesh
\[
    \Delta_0\gamma^j,
    \qquad j\in\mathbb Z.
\]
We therefore define
\[
    j_{\goldelta,\varepsilon}
    :=
    \min\left\{
        j\in\mathbb Z:
        \Delta_0\gamma^j\le\Delta_{\goldelta,\varepsilon}
    \right\},
\]
and set
\[
    \bar\Delta_{\goldelta,\varepsilon}
    :=
    \Delta_0\gamma^{j_{\goldelta,\varepsilon}}.
\]
Then
\[
    \bar\Delta_{\goldelta,\varepsilon}
    \le
    \Delta_{\goldelta,\varepsilon},
\]
so Lemma~\ref{lem:good-event-sufficient-decrease} and  Lemma~\ref{lem:good-event-probability} apply whenever
 \(\Delta_k\le\bar\Delta_{\goldelta,\varepsilon}\).
\end{remark}
We are now ready to combine the merit-function drift, the radius recursion, and
the lower bound on the good-event probability within the stopping-time framework.
This gives an expected bound on the number of outer iterations needed to reach
\((\goldelta,\varepsilon)\)-Goldstein stationarity. It is important to highlight that assuming $c_{max}$ is uniformly bounded (see end of Remark \ref{rem:M-lower-bound}) keeps the constants in the
outer-iteration bound of Corollary~\ref{cor:expected-stopping} uniform in \(\varepsilon\). 


\begin{corollary}[Expected stopping time]
\label{cor:expected-stopping}
Assume the merit-function drift condition given in  Lemma~\ref{lem:phi_drift_main}
holds with $\Theta:=\eta(1-\gamma^{\Dexp})-c_{\max}>0,$
and let the radius recursion and success-probability bound be as in
Lemma~\ref{lem:good-event-sufficient-decrease} and Lemma~\ref{lem:good-event-probability}. Suppose that \(\bar{m}\) is chosen so that
\[
    q_{\bar{m}}
    :=
    \beta\left[1-(1-p_\rho\beta)^{\bar{m}}\right]
    > q_0>
    \frac12 .
\]
Then we obtain the expected iteration complexity bound
\[
    \mathbb E[T_{\goldelta,\varepsilon}]
    =
    \mathcal O\left(
        \max\left\{
            \goldelta^{-\Dexp},
            \varepsilon^{-\Dexp/(\Dexp-1)}
        \right\}
    \right).
\]
% up to constants depending on \(q_M\), \(\Theta\), \(\gamma\), and the initial
% potential \(\Phi_0\).






\end{corollary}
\begin{proof}
By Assumption (i) (bounded step sizes) and Lemma~\ref{lem:phi_drift_main} (merit-function drift with $h(\Delta)=\Delta^{\Dexp}$), items (i) and (iii) of Assumptions \ref{ass:5} hold. Lemma~\ref{lem:assumption-ii-holds} verifies item (ii) with the Bernoulli driver $W_{k+1}\in\{-1,+1\}$, parameter $q=q_{\bar{m}}>\tfrac12$, and $\lambda=\log(\gamma^{-1})$, together with the cap at $\bar{\Delta}_{\goldelta, \varepsilon}$. Substituting these constants into the renewal-reward bound (Theorem~\ref{t:expected complexity}) gives the stated inequality. More specifically, Theorem~\ref{t:expected complexity} applies with
\[
    h(\Delta)=\Delta^{\Dexp},
    \qquad
    q=q_{\bar{m}},
    \qquad
    \Delta_{\goldelta,\varepsilon}
    =
    \bar\Delta_{\goldelta,\varepsilon},
\]
and gives
\begin{equation}\label{eq:Cartis}
    \mathbb E[T_{\goldelta,\varepsilon}]
    \le
    \frac{q_{\bar{m}}}{2q_{\bar{m}}-1}
    \frac{\Phi_0}{\Theta\bar\Delta_{\goldelta,\varepsilon}^{\Dexp}}
    +1 .
\end{equation}
By construction of the analytic threshold in Lemma~\ref{lem:good-event-sufficient-decrease}, we have that $\Delta_{\goldelta,\varepsilon}
    \asymp
    \min\left\{
        \goldelta,
        \varepsilon^{1/(\Dexp-1)}
    \right\}.$ 
Therefore, it is easy to see that also
$ \bar\Delta_{\goldelta,\varepsilon}
    \asymp
    \min\left\{
        \goldelta,
        \varepsilon^{1/(\Dexp-1)}
    \right\}.$
It thus follows that
\[
    \frac{1}{\bar\Delta_{\goldelta,\varepsilon}^{\Dexp}}
    =
    \mathcal O\left(
        \max\left\{
            \goldelta^{-\Dexp},
            \varepsilon^{-\Dexp/(\Dexp-1)}
        \right\}
    \right).
\]
Substituting this estimate into~\eqref{eq:Cartis} gives
\[
    \mathbb E[T_{\goldelta,\varepsilon}]
    =
    \mathcal O\left(
        \max\left\{
            \goldelta^{-\Dexp},
            \varepsilon^{-\Dexp/(\Dexp-1)}
        \right\}
    \right),
\]
which proves the claim.



\end{proof}
% Unlike the stepsize-based complexity (where $\Delta_k\ge\varepsilon$ before stopping and we can use a
% uniform bound $\overline H(\varepsilon)$ on $\EE[H_k\mid\F_{k-1}]$), the Clarke-stationarity stopping time
% does not provide a stepsize lower bound. Hence we cannot control $H_k$ iteration-wise. Instead, we control
% the \emph{sum} of extrapolation counts up to the Clarke stopping time by a telescoping identity on the
% log-radius. 
The previous corollary counts outer iterations. We now translate this bound into
a tested-point complexity estimate by using the fact that, at each iteration, the
algorithm evaluates at most \(\bar{m}\) directions and at most \(\Imax+1\)
extrapolation levels per direction. Also in this case, we note that tightening
the sample-average accuracy to make $c_{max}$ uniformly bounded may increase the raw stochastic sample complexity, but not the tested-point complexity considered in Corollary~\ref{cor:finalcompl}.

\begin{corollary}\label{cor:finalcompl}[Expected tested-point complexity to Goldstein stationarity]
\label{cor:tested-point-Goldstein}
Let \(T_{\goldelta,\varepsilon}\) be the stopping time defined in
\eqref{Teps}. Assume the hypotheses of Corollary~\ref{cor:expected-stopping}. 
Fix $q_0\in(1/2,\beta)$, and choose \(\bar{m}\) large enough so that
\(q_{\bar{m}}\ge q_0\). Suppose moreover that \(\Imax\) is fixed and that the
line search tests at most \(\bar{m}\) directions per outer iteration and at most
\(\Imax+1\) extrapolation levels per direction. Assume $\bar{m}$ is chosen
according to Remark~\ref{rem:M-lower-bound}. Then
\[
    \mathbb E\!\left[
        \sum_{k=0}^{T_{\goldelta,\varepsilon}-1} (1+P_k)
    \right]
    =
    \mathcal O\left(
        \varepsilon^{1-n}
        \max\left\{
            \goldelta^{-\Dexp},
            \varepsilon^{-\Dexp/(\Dexp-1)}
        \right\}
    \right),
\]
where \(P_k\) denotes the number of trial points evaluated during iteration
\(k\), excluding the baseline point \(X_k\).


% up to constants depending on \(\gamma,\beta,L,n,\Theta,\Phi_0\), and the cap
% geometry.
\end{corollary}
\begin{proof}
At every outer iteration, the algorithm tests at most \(\bar{m}\) directions. For
each direction, the line search evaluates at most the levels
$ i=0,1,\ldots,\Imax.$
Therefore the number of trial points evaluated during iteration \(k\) satisfies
\[
    P_k\le \bar{m}(\Imax+1)
    \qquad\text{a.s.}
\]
Summing up to the stopping time gives
\[
    \sum_{k=0}^{T_{\goldelta,\varepsilon}-1} P_k
    \le
    \bar{m}(\Imax+1)T_{\goldelta,\varepsilon}.
\]
Taking expectations yields
\[
    \mathbb E\!\left[
        \sum_{k=0}^{T_{\goldelta,\varepsilon}-1} P_k
    \right]
    \le
    \bar{m}(\Imax+1)\mathbb E[T_{\goldelta,\varepsilon}].
\]
If the baseline estimate is counted once per outer iteration, the same argument
gives
\[
    \mathbb E\!\left[
        \sum_{k=0}^{T_{\goldelta,\varepsilon}-1} (1+P_k)
    \right]
    \le
    \bigl(1+\bar{m}(\Imax+1)\bigr)
    \mathbb E[T_{\goldelta,\varepsilon}].
\]
The rest follows by substituting the bound from
Corollary~\ref{cor:expected-stopping}. Finally, if \(\bar{m}\) is chosen minimally to ensure
\(q_{\bar{m}}>1/2\), the scaling \(\bar{m}=\mathcal O(\varepsilon^{1-n})\) follows from
Remark~\ref{rem:M-lower-bound}. Thus the result is proved.

\end{proof}

\begin{remark}[Tested points versus raw stochastic oracle samples]
Corollary~\ref{cor:tested-point-Goldstein} counts calls to the estimator
\(\tilde f\), or equivalently the number of tested points. It does not count the
raw stochastic oracle samples used to construct each estimate. If \(\tilde f\) is computed by leveraging a batch size \(m_k\) of function evaluations at
iteration \(k\), then the raw sample cost is proportional to
\[
    \sum_{k<T_{\goldelta,\varepsilon}} m_k P_k,
\]
where \(P_k\) is the number of estimator calls at iteration \(k\). Thus Corollary~\ref{cor:tested-point-Goldstein} should be interpreted as a
tested-point complexity bound, not as a raw stochastic-oracle sample complexity
bound.
\end{remark}

% \begin{remark}[Complexity Comparison]
% \label{rem:complexity_comparison}
% To put the expected total tested-point complexity of Corollary \ref{cor:finalcompl} into perspective, consider the symmetric target regime where we seek an $(\varepsilon, \varepsilon)$-Goldstein stationary point (i.e., setting $\goldelta = \varepsilon$). The complexity bound algebraically factors as:
% \[
%    \mathcal O\Bigl( \varepsilon^{1-n} \max\bigl\{\varepsilon^{-p},\ \varepsilon^{-\frac{p}{p-1}}\bigr\} \Bigr) 
%    = \mathcal O\Bigl( \varepsilon^{1-n} \cdot \varepsilon^{-\frac{p}{\min(p-1, 1)}} \Bigr).
% \]
% This factorization perfectly isolates the theoretical cost of our barebones, purely nonsmooth approach. The second term, $\mathcal O(\varepsilon^{-p/\min(p-1, 1)})$, is exactly the expected complexity bound established by \cite{MR4359469} for stochastic directional direct search on \emph{smooth} ($C^{1,1}$) objective functions. The first term, $\mathcal O(\varepsilon^{1-n})$, represents the strict topological penalty for navigating a nonsmooth landscape: it is the geometric cost of sampling a uniform random direction and hoping it aligns with an arbitrarily narrow nonsmooth descent cone, without the guarantee provided by the cosine measure of positive spanning sets in smooth settings.\\
% It is also useful to compare this bound with randomized-smoothing
% zeroth-order methods, such as those of Nesterov and Spokoiny~\cite{MR3627456}.
% Those methods replace the original objective \(f\) by a smoothed objective $f_\mu(x)=\mathbb E[f(x+\mu u)]
% $
% and use random finite-difference or directional-derivative estimators to
% approximate the gradient of the smoothed function. This approach can lead to
% sharper worst-case bounds in regimes where the smoothed problem is an
% appropriate proxy for the original one. However, the resulting performance
% depends on the choice of the smoothing radius \(\mu\). A larger value of
% \(\mu\) gives a smoother and more stable surrogate but may introduce a larger
% bias between \(f_\mu\) and \(f\); a smaller value of \(\mu\) reduces this
% bias but makes finite-difference estimators more sensitive to stochastic
% noise. Indeed, for estimators of the form $\frac{\tilde f(x+\mu u)-\tilde f(x)}{\mu}\,u,$
% independent observational noise is divided by \(\mu\). Thus, under a simple
% additive noise model with variance \(\sigma^2\), the noise contribution to
% the estimator variance scales like \(\mathcal O(\sigma^2/\mu^2)\). As
% \(\mu\) is reduced to improve the approximation of \(f\) by \(f_\mu\), one
% may therefore need additional sampling to stabilize the finite-difference
% estimate.

% % By contrast, the DSE method analyzed here does not optimize a smoothed
% % surrogate and does not attempt to reconstruct a gradient. It evaluates
% % candidate points and accepts steps through a sufficient-decrease test on stochastic function estimates directly. The price
% % paid for this more direct nonsmooth treatment is the worse cap-probability
% % factor in the expected complexity bound. The advantage is that the stationarity
% % guarantee is stated directly in terms of the Goldstein subdifferential of
% % the original locally Lipschitz function, and no smoothing radius has to be
% % tuned. This helps explain why, despite the more conservative complexity, the method can be competitive in noisy nonsmooth
% % black-box problems.

% By contrast, the DSE method analyzed here does not optimize a smoothed
% surrogate and does not attempt to reconstruct a gradient. It evaluates
% candidate points and accepts steps through a sufficient-decrease test on
% stochastic function estimates directly. The price paid for this direct
% nonsmooth treatment is the worse cap-probability factor in the expected
% complexity bound. This factor should however be understood as a worst-case geometric
% penalty.
% Moreover, for locally Lipschitz functions, points of nondifferentiability are
% negligible by Rademacher's theorem  (see, e.g.,
% \cite[Section~2.5]{Clarke1990}). Hence random trial points typically lie in
% smooth pieces of the objective, where the set of successful directions may be
% larger than the single certified cap used in the proof. Thus the true chance of success
% can be substantially larger than what the conservative lower bound used
% in the analysis says. The advantage in our framework is that the stationarity guarantee is stated
% directly in terms of the Goldstein subdifferential of the original locally
% Lipschitz function, and no smoothing radius has to be tuned. This helps
% explain why, despite the more conservative worst-case complexity, the method is competitive with
% random-gradient smoothing techniques in our empirical benchmarks.




% \end{remark}

\begin{remark}[Comparison with smooth stochastic direct search]
\label{rem:complexity_comparison_smooth}
To put the expected total tested-point complexity of Corollary \ref{cor:finalcompl} into perspective, consider the symmetric target regime where we seek an $(\varepsilon, \varepsilon)$-Goldstein stationary point (i.e., setting $\goldelta = \varepsilon$). The complexity bound algebraically factors as:
\[
   \mathcal O\Bigl( \varepsilon^{1-n}  \max\bigl\{\varepsilon^{-p},\ \varepsilon^{-\frac{p}{p-1}}\bigr\} \Bigr) 
   = \mathcal O\Bigl( \varepsilon^{1-n} \cdot \varepsilon^{-\frac{p}{\min(p-1, 1)}} \Bigr).
\]
This factorization isolates the theoretical cost of our direct nonsmooth approach. The second term,
\(\mathcal O(\varepsilon^{-p/\min(p-1,1)})\), is exactly the expected iteration-complexity
bound established by \cite{MR4359469} for stochastic directional direct search on smooth
\((C^{1,1})\) objective functions. In that setting, passing from iterations to tested points only
multiplies the bound by the cardinality of the positive spanning set, which is uniformly bounded
(for instance, of order \(n\) for standard positive bases). By contrast, in our nonsmooth random-polling
analysis, such a polling-set factor is replaced by
\(\mathcal O(\varepsilon^{1-n})\). This term represents the geometric cost of sampling enough
uniform random directions to hit a shrinking spherical cap aligned with a Goldstein descent
direction, without the deterministic alignment guarantee provided by the cosine measure of
positive spanning sets in smooth settings.
\end{remark}
It is also useful to make a comparison with randomized-smoothing
zeroth-order methods, such as those of Nesterov and Spokoiny~\cite{MR3627456}
and Kornowski and Shamir~\cite{JMLR:kornowskishamir}.
Those methods replace the original objective \(f\) by a smoothed surrogate
\(f_\mu\) and use random finite-difference or directional-derivative estimators
to approximate gradients of the smoothed function. This approach can lead to
sharper worst-case bounds in regimes where the smoothed problem is an
appropriate proxy for the original one. However, the resulting performance
depends on the choice of the smoothing radius \(\mu\) and on the stochastic
oracle structure. A larger value of \(\mu\) gives a smoother and more stable
surrogate but may introduce a larger bias between \(f_\mu\) and \(f\); a
smaller value of \(\mu\) reduces this bias but can make gradient estimators
more sensitive to stochastic noise. Thus, as \(\mu\) is reduced to improve the
approximation of \(f\) by \(f_\mu\), one may need additional sampling or
variance-control mechanisms to stabilize the estimate.\\
By contrast, the DSE method analyzed here does not optimize a smoothed
surrogate and does not attempt to reconstruct a gradient. It evaluates
candidate points and accepts steps through a sufficient-decrease test on
stochastic function estimates directly. The price paid for this direct
nonsmooth treatment is the worse cap-probability factor in the expected
complexity bound. This factor should however be understood as a worst-case
geometric penalty.
Moreover, for locally Lipschitz functions, points of nondifferentiability are
negligible by Rademacher's theorem  (see, e.g.,
\cite[Section~2.5]{Clarke1990}). Hence random trial points typically lie in
smooth pieces of the objective, where the set of successful directions may be
larger than the single certified cap used in the proof. Thus the true chance of
success can be substantially larger than what the conservative lower bound used
in the analysis says. The advantage in our framework is that no smoothing radius
has to be tuned. This helps explain why, despite the more conservative
worst-case complexity, the method is competitive with random-gradient smoothing
techniques in our empirical benchmarks.




% numerical experiments and discussion
% \section{Numerical results}
% \label{sec:numerics}

% We compare our direct search with extrapolation (DSE) against three baselines:
% the stochastic direct search (SDS) of~\cite{MR4758401}, StoMADS~\cite{MR4238147},
% and the zeroth-order Nesterov's random-gradient method proposed in \cite{MR3627456}.
% Unless specified otherwise, we choose the sufficient-decrease exponent as $\Dexp = 2$ for all methods,
% and we use the same budget and profiling protocol as in~\cite{MR4758401}.

% \subsection{DSE implementation and solvers compared}
% \label{subsec:impl}

% The DSE implementation used in the experiments is a practical variant of the
% theoretical scheme analyzed in previous sections.
% As in the SDS implementation from~\cite[Section 5]{MR4758401},
% we adopt a multi-phase strategy that improves robustness and speed in practice
% combining line searches along coordinates with line searches along dense directions.
% % :
% % \begin{itemize}
% % \item \textbf{Phase~I (coordinate line searches).} We perform independent line searches
% % along the $n$ coordinate directions and their opposites, each with its own stepsize.
% % \item \textbf{Phase~II (extrapolation on coordinates).} Once all coordinate stepsizes 
% % fall below a threshold $\bar\Delta_{1}$, we switch to extrapolation line searches on the
% % same $2n$ directions, as described by our acceptance rule.
% % \item \textbf{Phase~III (dense directions with extrapolation).} When all stepsizes
% % fall below a second threshold $\bar\Delta_{2}<\bar\Delta_{1}$, we switch to extrapolation
% % line searches along \emph{dense} directions drawn from a Sobol low-discrepancy sequence
% % to enhance uniform coverage of $\mathbb S^{n-1}$.
% % \end{itemize}
% SDS follows the implementation choices reported in~\cite[Section 5]{MR4758401}.
% For StoMADS we use the authors’ reference settings. For Nesterov's random-gradient we set parameters as discussed in the related paper.
% Amount of noise, oracle budget, starting point, initial stepsizes are fixed across all algorithms for fairness.

% \subsection{Data and Performance Profiles}
% \label{subsec:protocol}

% Following~\cite[Section 5]{MR4758401}, a run on problem $p$ is declared successful at
% iteration $k$ if
% \begin{equation}\label{eq:profile-tol}
%   f(x_k)\ \le\ f_L\;+\;\tau\bigl(f(x_0)-f_L\bigr),
% \end{equation}
% where $f_L$ is the best objective value achieved by any solver on $p$.
% We use the standard \emph{performance} and \emph{data} profiles from \cite{more2009benchmarking},
% with time counters $t_{p,s}$ defined as the number of function evaluations used by solver $s$
% to satisfy~\eqref{eq:profile-tol}. The performance profile of $s$ is
% \[
% \rho_s(\alpha)\;=\;\frac{1}{|{\mathcal P}|}\,
% \Bigl|\,\bigl\{\,p\in{\mathcal P}:\ 
% t_{p,s}\ \le\ \alpha\cdot \min_{s'} t_{p,s'}\ \bigr\}\Bigr|,
% \]
% and the data profile is
% \[
% d_s(\kappa)\;=\;\frac{1}{|{\mathcal P}|}\,
% \Bigl|\,\bigl\{\,p\in{\mathcal P}:\  t_{p,s}\ \le\ \kappa\,(n_p+1)\ \bigr\}\Bigr|,
% \]
% where $n_p$ is the dimension of problem $p$~\cite[Section~5]{MR4758401}.
% All profiles are computed with the \emph{true} objective values while the solvers operate on their (possibly stochastic) estimates, and the evaluation budget is fixed  as in~\cite[Section 5]{MR4758401}.
% We enforce a budget of $10^4(n_p+1)$ function evaluations per run.
% We report results for two tolerances $\tau\in\{10^{-2},10^{-4}\}$.\\
% To evaluate performances we use a standard nonsmooth DFO test set (unconstrained), aggregating all runs across problems and random seeds when building the profiles, as in~\cite[Section~5]{MR4758401}. The benchmark set comprises 96 instances of nonsmooth DFO problems (see \cite[Table 1]{MR4758401}). To stabilize comparisons, we solve each instance \emph{five} times per solver with different seeds, yielding $|P|=96 \times 5 = 480$ runs per solver.  
% % Profiles are built from the true objective values, and the evaluation budget is fixed as in \cite[Section 5]{MR4758401}. 
% % Table containing problems' names and dimensions can be found in Appendix, Table \ref{tab:rvz-bench}. 
% We simulated noise after averaging $p_k$ independent samples by adding to the objective $N(0, 1/p_k)$ distributed random variables. The remaining parameters were tuned with a basic grid search.

% % \subsection{Data and Performance Profiles}
% % \label{subsec:results}

% In Figure \ref{fig:profiles-grid}, we report data and performance profiles. It is easy to see that, for both tolerances, DSE (red, \texttt{dse}) outperforms the competitors for the whole range of budgets/ratios. More specifically, 

% \begin{itemize}
% \item \textbf{Data profiles.} The DSE curve lies strictly above those of SDS (\texttt{sds}),
% StoMADS (\texttt{sto}), and the random-gradient baseline (\texttt{rnd}) for essentially
% all budgets. At $\tau=10^{-2}$ the gap is pronounced, with DSE solving the largest
% fraction of problems at small budgets; Even at $\tau = 10^{-4}$, DSE continues to have a distinct advantage.
% \item \textbf{Performance profiles.} DSE attains the highest proportion of ``best-within-a-factor''
% wins across ratios $\alpha$. The separation from SDS is visible already near $\alpha\approx 1$,
% and remains throughout the range; For the majority of problems, StoMADS and the random-gradient method perform worse.
% \end{itemize}

% Overall, the profiles indicate that DSE is highly
% effective in practice and hhe multi-phase strategy helps to improve performances: the initial coordinate searches efficiently capture simple descent directions, while the Sobol-based dense directions give strong thorough exploration once stepsizes are small.



% % \subsection{Discussion}
% % \label{subsec:discussion}

% % The relative gains of DSE against SDS mirror the intuition from our analysis: dense directions increase the probability of finding sufficient decrease, while the extrapolation line search exploits curvature along promising rays. Compared to StoMADS, DSE benefits from (i) the dense search geometry in the small-stepsize regime and (ii) fewer unsuccessful polls thanks to extrapolation; both effects show up in the profiles. The random-gradient baseline underperforms on nonsmooth objectives,
% % as expected.\\
% % Finally, we stress that, as in~\cite[Section~5]{MR4758401}, our practical DSE implementation differs from the abstract algorithm used in the analysis. The convergence results still apply with minor refinements: it suffices that the initial warm-start phase is finite (a.s.) and that, from some index onward, the method employs dense directions on the sphere (e.g., i.i.d.\ uniform or a low-discrepancy sequence).


% \begin{figure}[H]
%   \centering

%   % ----- Row 1: tol = 1e-2 -----
%   \begin{subfigure}{0.4\textwidth}
%     \centering
%     \includegraphics[width=\linewidth]{images/dp_tol2.eps}
%     \caption{Data profile, $\tau=10^{-2}$}
%     \label{fig:data-2}
%   \end{subfigure}
%   \hspace{3em}
%   \begin{subfigure}{0.4\textwidth}
%     \centering
%     \includegraphics[width=\linewidth]{images/pp_tol2.eps}
%     \caption{Performance profile, $\tau=10^{-2}$}
%     \label{fig:perf-2}
%   \end{subfigure}

%   \vspace{0.75em}

%   % ----- Row 2: tol = 1e-4 -----
%   \begin{subfigure}{0.4\textwidth}
%     \centering
%     \includegraphics[width=\linewidth]{images/dp_tol4.eps}
%     \caption{Data profile, $\tau=10^{-4}$}
%     \label{fig:data-4}
%   \end{subfigure}
%   \hspace{3em}
%   \begin{subfigure}{0.4\textwidth}
%     \centering
%     \includegraphics[width=\linewidth]{images/pp_tol4.eps}
%     \caption{Performance profile, $\tau=10^{-4}$}
%     \label{fig:perf-4}
%   \end{subfigure}

%   \caption{Data (left) and performance (right) profiles. Top row: $\tau=10^{-2}$. Bottom row: $\tau=10^{-4}$.}
%   \label{fig:profiles-grid}
% \end{figure}



\section{Numerical Results}
\label{sec:numerics}

We compare our direct search with extrapolation (DSE) against four baselines:
the stochastic direct search (SDS) of~\cite{MR4758401}, StoMADS~\cite{MR4238147},
the zeroth-order gradient-smoothing method (GS) of Nesterov and Spokoiny~\cite{MR3627456},
and the Optimal Stochastic Nonsmooth Nonconvex Optimization Algorithm (OSNNOA)
of Kornowski and Shamir~\cite{JMLR:kornowskishamir}.
Unless specified otherwise, we choose the sufficient-decrease exponent as $\Dexp = 2$ for all methods,
and we use the same budget and profiling protocol as in~\cite{MR4758401}.
The implementation of the proposed method and the code used to reproduce the
numerical experiments are available on Github:  \url{https://github.com/APalmier99/DSE_StDFOpt.git}.


\subsection{DSE Implementation and Solvers Compared}
\label{subsec:impl}

The DSE implementation used in the experiments is a practical variant of the
theoretical scheme analyzed in previous sections.
As in the SDS implementation from~\cite[Section 5]{MR4758401},
we adopt a multi-phase strategy that improves robustness and speed in practice
combining line searches along coordinates with line searches along dense directions.
% :
% \begin{itemize}
% \item \textbf{Phase~I (coordinate line searches).} We perform independent line searches
% along the $n$ coordinate directions and their opposites, each with its own stepsize.
% \item \textbf{Phase~II (extrapolation on coordinates).} Once all coordinate stepsizes 
% fall below a threshold $\bar\Delta_{1}$, we switch to extrapolation line searches on the
% same $2n$ directions, as described by our acceptance rule.
% \item \textbf{Phase~III (dense directions with extrapolation).} When all stepsizes
% fall below a second threshold $\bar\Delta_{2}<\bar\Delta_{1}$, we switch to extrapolation
% line searches along \emph{dense} directions drawn from a Sobol low-discrepancy sequence
% to enhance uniform coverage of $\mathbb S^{n-1}$.
% \end{itemize}
SDS follows the implementation choices reported in~\cite[Section 5]{MR4758401}.
For StoMADS we use the authors' reference settings. For GS we set parameters as
discussed in~\cite{MR3627456}. For OSNNOA, we implemented the randomized-smoothing
scheme described in~\cite{JMLR:kornowskishamir}, using the proposed stochastic gradient
estimator and parameter choices as guidance.
Amount of noise, oracle budget, starting point, initial stepsizes are fixed across all algorithms for fairness following again the guidelines in~\cite[Section 5]{MR4758401}.

\subsection{Data and Performance Profiles}
\label{subsec:protocol}

Following~\cite{more2009benchmarking}, a run on problem $p$ is declared successful at
iteration $k$ if
\begin{equation}\label{eq:profile-tol}
  f(x_k)\ \le\ f_L\;+\;\tau\bigl(f(x_0)-f_L\bigr),
\end{equation}
where $f_L$ is the best objective value achieved by any solver on $p$.
We use the standard \emph{performance} and \emph{data} profiles,
with time counters $t_{p,s}$ defined as the number of function evaluations used by solver $s$
to satisfy~\eqref{eq:profile-tol}. The performance profile of $s$ is
\[
\rho_s(\alpha)\;=\;\frac{1}{|{\mathcal P}|}\,
\Bigl|\,\bigl\{\,p\in{\mathcal P}:\ 
t_{p,s}\ \le\ \alpha\cdot \min_{s'} t_{p,s'}\ \bigr\}\Bigr|,
\]
and the data profile is
\[
d_s(\kappa)\;=\;\frac{1}{|{\mathcal P}|}\,
\Bigl|\,\bigl\{\,p\in{\mathcal P}:\  t_{p,s}\ \le\ \kappa\,(n_p+1)\ \bigr\}\Bigr|,
\]
where $n_p$ is the dimension of problem $p$~\cite[Section~5]{MR4758401}.
All profiles are computed with the \emph{true} objective values while the solvers operate on their (possibly stochastic) estimates, and the evaluation budget is fixed  as in~\cite[Section 5]{MR4758401}.
We enforce a budget of $10^4(n_p+1)$ function evaluations per run.
We report results for two tolerances $\tau\in\{10^{-2},10^{-4}\}$.\\
To evaluate performances we use a standard nonsmooth DFO test set (unconstrained), aggregating all runs across problems and random seeds when building the profiles, as in~\cite[Section~5]{MR4758401}. The benchmark set comprises 96 instances of nonsmooth DFO problems (see \cite[Table 1]{MR4758401}). To stabilize comparisons, we solve each instance \emph{five} times per solver with different seeds, yielding $|P|=96 \times 5 = 480$ runs per solver.  
% Profiles are built from the true objective values, and the evaluation budget is fixed as in \cite[Section 5]{MR4758401}. 
% Table containing problems' names and dimensions can be found in Appendix, Table \ref{tab:rvz-bench}. 
We simulated noise after averaging $p_k$ independent samples by adding to the objective $N(0, 1/p_k)$ distributed random variables. The remaining parameters were tuned with a basic grid search.\\
% \subsection{Data and Performance Profiles}
% \label{subsec:results}
In Figure \ref{fig:profiles-grid}, we report data and performance profiles. It is easy to see that, for both tolerances, DSE outperforms the competitors for the whole range of budgets/ratios. More specifically, 
\begin{itemize}
\item \textbf{Data profiles.} The DSE curve (blue) lies strictly above those of SDS (red),
StoMADS (black),  OSNNOA (china blue) and GS (purple), for essentially
all budgets. At $\tau=10^{-2}$ the gap is pronounced, with DSE solving the largest
fraction of problems at small budgets; Even at $\tau = 10^{-4}$, DSE continues to have a distinct advantage.
\item \textbf{Performance profiles.} DSE attains the highest proportion of ``best-within-a-factor''
wins across ratios $\alpha$. The separation from SDS is visible already near $\alpha\approx 1$,
and remains throughout the range; For the majority of problems, StoMADS, GS, and OSNNOA perform worse.
\end{itemize}
Overall, the profiles indicate that DSE is highly
effective in practice and the multi-phase strategy helps to improve performances: the initial coordinate searches efficiently capture simple descent directions, while the dense directions give strong thorough exploration once stepsizes are small.
\begin{figure}[t]
  \centering

  % ----- Row 1: tol = 1e-2 -----
  \begin{subfigure}{0.4\textwidth}
    \centering
    \includegraphics[width=\linewidth]{images/dp_tol2.eps}
    \caption{Data profile, $\tau=10^{-2}$}
    \label{fig:data-2}
  \end{subfigure}
  \hspace{3em}
  \begin{subfigure}{0.4\textwidth}
    \centering
    \includegraphics[width=\linewidth]{images/pp_tol2.eps}
    \caption{Performance profile, $\tau=10^{-2}$}
    \label{fig:perf-2}
  \end{subfigure}

  \vspace{0.75em}

  % ----- Row 2: tol = 1e-4 -----
  \begin{subfigure}{0.4\textwidth}
    \centering
    \includegraphics[width=\linewidth]{images/dp_tol4.eps}
    \caption{Data profile, $\tau=10^{-4}$}
    \label{fig:data-4}
  \end{subfigure}
  \hspace{3em}
  \begin{subfigure}{0.4\textwidth}
    \centering
    \includegraphics[width=\linewidth]{images/pp_tol4.eps}
    \caption{Performance profile, $\tau=10^{-4}$}
    \label{fig:perf-4}
  \end{subfigure}

  \caption{Data (left) and performance (right) profiles. Top row: $\tau=10^{-2}$. Bottom row: $\tau=10^{-4}$.}
  \label{fig:profiles-grid}
\end{figure}


% conclusions 

\section{Conclusions and Future Directions}
\label{sec:conclusions}

We proposed and analyzed DSE, a stochastic direct-search method with extrapolation
for unconstrained nonsmooth zeroth-order optimization. The meth\-od combines random
polling directions with a stochastic sufficient-decrease line search and uses only noisy
function-value estimates. Under conditional tail and independence assumptions on the
oracle errors, we proved almost-sure convergence of refining subsequences to Clarke stationary
points.\\
We further established expected complexity bounds for reaching
\((r,\varepsilon)\)-Goldstein stationarity. The proof combines a merit-function
drift argument with a supermartingale stopping-time framework and a
spherical-cap probability estimate for the random polling directions. At the
outer-iteration level, we obtain an expected iteration complexity 
$\mathcal O\left(
        \max\left\{
            r^{-p},
            \varepsilon^{-p/(p-1)}
        \right\}
    \right).$
Moreover, after choosing $\bar{m}$, the maximum number of polling directions to be used at each iteration,  large enough to keep the
success probability uniformly bounded away from \(1/2\), the expected
tested-point complexity satisfies
$\mathcal O\left(
        \varepsilon^{1-n}
        \max\left\{
            r^{-p},
            \varepsilon^{-p/(p-1)}
        \right\}
    \right).$ The numerical results show that extrapolation can substantially improve the practical
performance of stochastic direct-search schemes on nonsmooth benchmark problems.
In particular, the proposed implementation compares favorably with SDS, StoMADS,
and randomized-smoothing baselines on the tested instances.\\
Several directions remain open. First, it would be interesting to sharpen the
dimension dependence in the nonsmooth complexity bound or to identify structural
conditions under which the conservative spherical-cap factor can be improved.  Second, extending the analysis to variance-reduction mechanisms,
common-random-number estimators, and constrained nonsmooth stochastic problems
would further broaden the applicability of the approach.



% We introduced and analyzed DSE, a direct-search method that combines a  line search with dense random directions to solve unconstrained problems with locally Lipschitz (possibly nonsmooth) objectives using only noisy zeroth-order evaluations. Under the estimation conditions of Section~3 (a reduction-tail control with exponent $\Dexp>1$ together with a selected-level $L^1$ bound), we proved almost-sure convergence to Clarke stationary points and derived an $\mathcal{O}(\varepsilon^{-\Dexp})$ expected oracle complexity bound for achieving $\varepsilon$-Clarke stationarity. The proof is based on three main components: a drift argument implying $\sum_k \Delta_k^\Dexp<\infty$ and $\Delta_k\to 0$ a.s.; tail bounds for the random line-search depth $H_k$, leading to the sampling schedule
% \[
% m_k=\Theta\!\big(\Delta_k^{-2\Dexp}\,[1+\log(1/\Delta_k)]^{2}\big),
% \]
% whose $\log$ factor is the mild price of working under minimal accuracy requirements; and a spherical-cap success guarantee on $\mathbb{S}^{n-1}$ showing that, with probability strictly larger than $1/2$, at least one of the sampled directions is sufficiently aligned with a Clarke-descent direction. These ingredients allow us to invoke the renewal-reward framework of \cite{MR4151319} and obtain the final expected complexity bound.
% \noindent
% Several extensions are natural. A primary objective is to reduce (or remove) the $\log^2(1/\Delta_k)$ overhead without resorting to fixed-confidence assumptions, for instance via sharper control of $H_k$ or a modified inner stopping rule. It would also be interesting to replace i.i.d.\ uniform directions with more structured dense constructions (e.g., low-discrepancy sequences) or adaptive/learned direction sets while preserving Clarke-stationarity guarantees. Finally, extending the framework to constrained problems and to more general noise models (e.g., heavy-tailed or dependent noise), while incorporating variance-reduction mechanisms beyond CRN, is also a natural next step.

% We analyzed a direct-search algorithm with a stochastic, run–until–failure line search along dense random directions, under minimal information: locally Lipschitz objectives and a mean–zero, finite-variance oracle. The core technical ingredients were (i) a selection-aware $L^1$ control for the reduction-estimator at the \emph{selected} inner level, which neutralizes the selection bias induced by the line-search depth; (ii) explicit tail bounds for the inner run length $H_k$; and (iii) a Cartis–Scheinberg supermartingale/stopping-time argument that yields expected evaluation complexity in a \emph{purely zeroth-order} fashion. On the convergence side, we established a positive-drift inequality, showed $\sum_k \Delta_k^p<\infty$ and $\Delta_k\to 0$ almost surely, and proved that every refining subsequence converges to a Clarke-stationary point. On the complexity side, we combined a geometric (spherical-cap) success mechanism for multiple dense directions with a capped inner routine to verify the renewal–reward hypotheses and obtain a clean bound on the expected stopping time to $\varepsilon$–Clarke stationarity. These results are obtained without variance–stepsize coupling and with no smoothness beyond local Lipschitz continuity. Additionally, we showed that simple sample averages enforce the required tail and post-selection controls, with a practical batch schedule
% \[
% m_k=\Theta\!\big(\Delta_k^{-2\Dexp}[1+\log(1/\Delta_k)]^2\big),
% \]
% where the benign $\log^2$ factor originates from summing the run–until–failure depth. Under stronger structural assumptions from the literature the $\log^2$ overhead disappears and the standard algebraic scalings are recovered. In addition, the reduced-sample-sizing framework based on \emph{reduction} tails applies directly in our setting and, with CRN and mild sample-path correlation, sharpens the batch size by a factor $\Delta_k^2$.
% Compared with recent stochastic line-search approaches that assume $L$–Lipschitz \emph{gradients}, single-level fixed-confidence accuracy, and variance–stepsize coupling, our analysis operates in a fully nonsmooth, gradient-free regime, tailors the proofs to the dependence created by the inner line search, and measures progress in terms of Clarke stationarity. The resulting guarantees are qualitatively different (zeroth order, Clarke-based) yet align with classical rates once stronger noise/accuracy structures are imposed.\\
% \emph{Limitations and open directions.} Our selection-aware analysis introduces a mild polylogarithmic overhead in the batching schedule; removing or tightening this factor without appealing to fixed-confidence assumptions is an interesting open problem. The success-probability lower bound currently uses multiple dense directions per iteration; adapting this to structured direction sets or learned direction distributions (while retaining Clarke-based guarantees) is another promising path. Extending the framework to constrained problems (e.g., via progressive barriers or stochastic feasibility tests), to non-i.i.d.\ or heavy-tailed noise, and to variance-reduction mechanisms beyond CRN are natural next steps. Finally, it would be valuable to explore adaptive choices of the number of trial directions $\bar{m}$ and the cap parameter that optimize the renewal–reward constants. In conclusion, direct-search with a stochastic run–until–failure line search admits rigorous convergence to Clarke stationarity and expected evaluation complexity under minimal, practically verifiable assumptions. The analysis hinges on controlling the post-selection error at the chosen inner level and on coupling dense random directions with a renewal–reward view of step-size dynamics. We believe this reduction-centric, selection-aware perspective can inform broader classes of stochastic zeroth-order methods in nonsmooth optimization.

%\pagebreak

% % appendix
\appendix
\section{Proofs of the Theoretical Results} \label{appendix:A}

% \subsection{Section \ref{sec:assumptions}: Assumptions and preliminaries}
\begin{proof}[Lemma \ref{lem: conditional mean bound}]
Work conditionally on $\mathcal F_{k-1}$ throughout. Introduce the normalized error
\[
Z_{k,i}\ :=\ \frac{|\bar E_{k,i}|}{\Delta_k^{\Dexp}}\ \ge 0 .
\]
Assumption~\ref{ass:quadratic tail} is exactly the conditional tail bound
\[
\mathbb{P}\!\left(Z_{k,i}\ \ge\ \alpha\ \middle|\ \mathcal F_{k-1}\right)\ \le\ \frac{\varepsilon_h}{\alpha^{\aexp}}
\qquad\forall\,\alpha\ge\varepsilon_h.
\]
By the layer-cake representation,
\[
\mathbb{E}\!\left[Z_{k,i}\ \middle|\ \mathcal F_{k-1}\right]
=\int_0^\infty \mathbb{P}\!\left(Z_{k,i}\ge t\ \middle|\ \mathcal F_{k-1}\right)\,dt.
\]
Fix any cutoff $c\ge\varepsilon_h$ and split the integral:
\[
\mathbb{E}[Z_{k,i}\mid\mathcal F_{k-1}]
\ \le\ \int_0^{c} 1\,dt
\ +\ \int_{c}^{\infty} \frac{\varepsilon_h}{t^{\aexp}}\,dt
\ =\ c+\varepsilon_h(\Dexp-1)c^{-1/(\Dexp-1)}.
\]
The right-hand side is minimized (over $c\ge\varepsilon_h$) at
\(
c^\star=\max\{\varepsilon_h^{(\Dexp-1)/\Dexp},\ \varepsilon_h\}
\)
so
\[
\min_{c\ge\varepsilon_h}\Big(c+\varepsilon_h(\Dexp-1)c^{-1/(\Dexp-1)}\Big)
=
\begin{cases}
\Dexp \varepsilon_h^{(\Dexp-1)/\Dexp}, & \varepsilon_h\le 1, \quad \textit{case A}\\[3pt]
\varepsilon_h \left[ 1+ (\Dexp-1)\varepsilon_h^{-1/(\Dexp-1)}\right], & \varepsilon_h>1 \quad \textit{ case B}.
\end{cases}
\]
Therefore
\(
\mathbb{E}[Z_{k,i}\mid\mathcal F_{k-1}]\le \mu_h
\)
with
% \(
% \mu_h=\min\{\textit{case A},\,\textit{case B}\}.
% \)
$\mu_h$ defined as above.
Multiplying back by $\Delta_k^{\Dexp}$ yields
\(
\mathbb{E}[|\bar E_{k,i}|\mid\mathcal F_{k-1}]\le \mu_h\,\Delta_k^{\Dexp}
\),
uniformly in $k$ and $i$, and the result is proved.
\end{proof}



\begin{proof}[Lemma \ref{lem:diff-L1-from-A44}]
Work conditionally on $\mathcal F_{k-1}$ throughout. By the previous lemma, we have
\[
\mathbb E\big[\,|\bar E_{k}^{(-1)}|\,\big|\,\mathcal F_{k-1}\big]\ \le\ \mu_h\,\Delta_k^{\Dexp},
\qquad
\mathbb E\big[\,|\bar E_{k,m}^{(i)}|\,\big|\,\mathcal F_{k-1}\big]\ \le\ \mu_h\,\Delta_k^{\Dexp}
\quad \forall i,\ \forall m.
\]
Using $E_{k,m}^{(i)}=\bar E_{k}^{(-1)}-\bar E_{k,m}^{(i)}$ and the triangle inequality,
\[
\mathbb E\big[\,|E_{k,m}^{(i)}|\,\big|\,\mathcal F_{k-1}\big]
\ \le\ \mathbb E\big[\,|\bar E_{k}^{(-1)}|\,\big|\,\mathcal F_{k-1}\big]
      + \mathbb E\big[\,|\bar E_{k,m}^{(i)}|\,\big|\,\mathcal F_{k-1}\big]
\ \le\ 2\,\mu_h\,\Delta_k^{\Dexp},
\]
which is \eqref{eq:diff-L1-core}. The signed bound \eqref{eq:diff-signed} follows since
$|\mathbb E[E_{k,m}^{(i)}\mid\mathcal F_{k-1}]|\le \mathbb E[|E_{k,m}^{(i)}|\mid\mathcal F_{k-1}]$ by Jensen inequality.
\end{proof}
\noindent
Finally, bounding $E^{\max}_k \leq \sum_{m}\sum_{i}|E_{k,m}^{(i)}|$ and applying the previous result, immediately yields Lemma \ref{ass:postL1} as well.



% \begin{remark}[Post-selection $L^1$ via tails of $H_k$]
% \label{rem:postsel_tail}
% For each level $i\ge 0$ define the (difference) estimation error
% \[
% E_{k,i}
% :=\big(\tilde f(X_k)-\tilde f(X_k^{(i)})\big)-\big(f(X_k)-f(X_k^{(i)})\big),
% \]
% and let $H_k$ be the random selected level from the inner line-search (success-until-first-failure rule with steps $\Delta_{k,i}=\gamma^{-i}\Delta_k$). Then, for a suitable batch schedule $m_k = O\left(\Delta_k^{-2\Dexp}(1+\log(1/\Delta_k))^2\right)$, it holds
% \begin{equation}
% \label{eq:PSL1}
% \me\!\left[\,|E_{k,H_k}|\,\middle|\,\F_{k-1}\right]\ \le\ c\,\Delta_k^{\Dexp},
% \end{equation}
% with a constant $c>0$ independent of $k$.
% \end{remark}
% \begin{proof}
% By the tower property and Cauchy--Schwarz,
% \begin{align}
% \me\!\left[\,|E_{k,H_k}|\,\middle|\,\F_{k-1}\right]
% &=\sum_{i\ge 0}\me\!\left[\,|E_{k,i}|\,\mathbf 1_{\{H_k=i\}}\,\middle|\,\F_{k-1}\right] \notag \\
% &\le \sum_{i\ge 0}\Big(\me\!\left[E_{k,i}^{2}\,\middle|\,\F_{k-1}\right]\Big)^{\!1/2}
% \Big(\Pb(H_k=i\,|\,\F_{k-1})\Big)^{\!1/2}. \label{eq:CS_selected}
% \end{align}
% Assume the additive oracle model $\tilde f=f+\zeta$ with $\me[\zeta\,|\F_{k-1}]=0$ and $\mathrm{Var}(\zeta\,|\F_{k-1})\le\sigma^2$, and let $m_k$ be the number of i.i.d. samples used to build the function estimates at each point (same $m_k$ for all inner points at iteration $k$). The two means in $E_{k,i}$ are based on independent batches, hence
% \begin{equation}
% \label{eq:Ekivar}
% \me\!\left[E_{k,i}^{2}\,\middle|\,\F_{k-1}\right]\ \le\ \frac{2\sigma^2}{m_k}
% \qquad\text{for all }i\ge 0.
% \end{equation}
% Define the tail constant
% \begin{equation}
% \label{eq:Ctail_def}
% C_{\text{tail}}(k)\ :=\ \sum_{i\ge 0}\sqrt{\Pb(H_k=i\, | \,\F_{k-1})}.
% \end{equation}
% Plugging \eqref{eq:Ekivar} into \eqref{eq:CS_selected} yields
% \begin{equation}
% \label{eq:PSL1_pre}
% \me\!\left[\,|E_{k,H_k}|\,\middle|\,\F_{k-1}\right]
% \ \le\ \sqrt{\tfrac{2\sigma^2}{m_k}}\;C_{\text{tail}}(k).
% \end{equation}
% By Lemma~10.1 (tail of the inner linesearch), with
% \[
% A_k\ :=\ \frac{2\mu_h\,\Delta_k^{\Dexp}+2f_{\max}}{\theta\,\Delta_k^{\Dexp}},
% \qquad \text{and}\qquad
% \Pb(H_k\ge i\, |\,\F_{k-1})\ \le\ \min\{1,\ A_k\,\gamma^{\Dexp i}\}\quad(i\ge 0),
% \]
% define $J_k:=\max\!\left\{0,\ \left\lfloor \tfrac{1}{\Dexp}\log_{1/\gamma}(A_k)\right\rfloor\right\}$ so that $A_k\gamma^{\Dexp J_k}\le 1$. Using $\Pb(H_k=i)\le\Pb(H_k\ge i)$ and splitting the sum at $J_k$,
% \begin{align}
% C_{\text{tail}}(k)
% &\le \sum_{i=0}^{J_k-1}1\;+\;\sum_{i\ge J_k}\sqrt{A_k}\,\gamma^{\frac{\Dexp}{2}i}
% \ \le\ J_k\ +\ \frac{\sqrt{A_k}\,\gamma^{\frac{\Dexp}{2}J_k}}{1-\gamma^{\Dexp/2}}
% \ \le\ J_k\ +\ \frac{1}{1-\gamma^{\Dexp/2}}.
% \label{eq:Ctail_bound}
% \end{align}
% Here $f_{\max}<\infty$ follows from compactness assumption \ref{ass:compact}.
% Moreover, the explicit bound
% \begin{equation}
% \label{eq:Jk_simple}
% J_k\ \le\ \frac{1}{\Dexp\log(1/\gamma)}\,
% \log\!\Big(\frac{2\mu_h}{\theta}+\frac{2f_{\max}}{\theta\,\Delta_k^{\Dexp}}\Big)\ +\ 1
% \end{equation}
% shows $J_k=O(\log(1/\Delta_k))$ as $\Delta_k\downarrow 0$.
% Combining \eqref{eq:PSL1_pre} and \eqref{eq:Ctail_bound}, a sufficient condition for \eqref{eq:PSL1} is
% \begin{equation}
% \label{eq:mk_exact}
% m_k\ \ge\ \frac{2\sigma^2}{c^2}\,\Big(J_k+\tfrac{1}{1-\gamma^{\Dexp/2}}\Big)^{\!2}\ \Delta_k^{-2 \Dexp}.
% \end{equation}
% % Using \eqref{eq:Jk_simple} to upper bound $J_k$, one obtains the simpler, uniform schedule
% % \begin{equation}
% % \label{eq:mk_simple}
% % m_k\ =\ m_0\,\Delta_k^{-2\Dexp}\Big(1+\log\tfrac{1}{\Delta_k}\Big)^2,
% % \qquad
% % m_0\ \ge\ \frac{2\sigma^2}{c^2}\,
% % \Big(1+\tfrac{1}{1-\gamma^{p/2}}+\tfrac{1}{p\log(1/\gamma)}\Big)^{\!2},
% % \end{equation}
% % which also guarantees \eqref{eq:PSL1} for all $k$.
% In particular, $m_k = O\left(\Delta_k^{-2\Dexp}(1+\log(1/\Delta_k))^2\right)$.
% \end{proof}

%%%%%%%%%%%%%%%%%%%%%%%%%%%%%%%%%%%%%%%%%%%%%%%%%%%%%%%%%%%%%%%

% \subsection{Section \ref{sec:clarke complexity}: Clarke Complexity}

%%%%%%%%%%%%%%%%%%%%%%%%%%%%%%%%%%%%%%%%%%%%%%%%%%%%%%%%%%%%%%%
% \subsection{Section \ref{sec:numerics}: Numerical results}
% \begin{table}[H]
%     \caption{Benchmark problems (name and dimension) used in the profiles.}
%     \label{tab:rvz-bench}
%     \centering
%     \begin{tabular}{c c|c c}
%     \textit{Name} & \textit{Dim} & \textit{Name} & \textit{Dim}\\
%     \hline
%     % ---- rows (left: name & dim) (right: name & dim) ----
%     crescent & 2            & watson            & 20 \\
%     cb2      & 2            & osborne 2         & 11 \\
%     charconn1 & 2           & shor              & 5 \\
%     charconn2 & 2           & colville 1        & 5 \\
%     demyanov-malozemov & 2  & hs 78             & 5 \\
%     dennis-woods & 2        & maxquad           & 10 \\
%     wong1    & 7            & gill              & 10 \\
%     wong2    & 10           & mxhilb            & 50 \\
%     wong3    & 20           & l1hilb            & 50 \\
%     elattar  & 6            & davidon 2         & 4 \\
%     goffin   & 50           & shelldual         & 15 \\
%     hald-madsen 1 & 2       & steiner 2         & 12 \\
%     lq       & 2            & transformer       & 6 \\
%     ql       & 2            & polak 6.10        & 1 \\
%     maxl     & 20           & wild1             & 20 \\
%     maxq     & 20           & wild2             & 20 \\
%     mifflin 1 & 2           & wild3             & 20 \\
%     mifflin 2 & 2           & wild19            & 20 \\
%     rosen-suzuki & 4        & wild11            & 20 \\
%     wf       & 2            & wild16            & 20 \\
%     spiral   & 2            & wild20            & 20 \\
%     evd52    & 3            & wild15            & 20 \\
%     kowalik-osborne & 4     & wild21            & 20 \\
%     oet5     & 4            & maxq              & \{10, 20, 30, 40\} \\
%     oet6     & 4            & l1hilb            & \{10, 20, 30, 40\} \\
%     gamma    & 4            & lq                & \{10, 20, 30, 40\} \\
%     exp      & 5            & cb3               & \{10, 20, 30, 40\} \\
%     pbc1     & 5            & cb32              & \{10, 20, 30, 40\} \\
%     evd61    & 6            & af                & \{10, 20, 30, 40\} \\
%     filter   & 9            & brown             & \{10, 20, 30, 40\} \\
%     polak 2  & 10           & mifflin2          & \{10, 20, 30, 40\} \\
%     polak 3  & 11           & crescent          & \{10, 20, 30, 40\} \\
%     polak 6  & 4            & crescent2         & \{10, 20, 30, 40\} \\
%     \hline     
%     \end{tabular}
% \end{table}

% As in \cite[Section 5 \& App.~A.2]{MR4758401}, the benchmark suite comprises 96 derivative-free nonsmooth problems (Table~\ref{tab:rvz-bench}). To stabilize comparisons, \emph{each problem is solved five times} per solver with different seeds, for a total of $96\times 5=480$ runs per solver.  Profiles are built from true objective values.

% \subsection{Complexity bound using Positive Spanning Sets}
% % ---------------------------------------------------------------------
% %  Lemma — Uniform probabilistic sufficient decrease
% %  using a fixed positive–spanning set
% % ---------------------------------------------------------------------

% As an alternative way to derive the complexity bound, we could introduce some quality measures for the directions used in each iteration. Polling directions should be chosen in such a way that their significant deterioration can be avoided asymptotically, i.e., in such a way to ensure that they never become close to losing the positive spanning property. We follow the same approach as in \cite{MR4359469} and use the cosine similarity measure to quantify the worst possible approximation of the positive spanning sets $\mathbb{D}_k$'s
% \begin{equation}
%     k(\mathbb{D}_k) \defeq \min_{v \in \mathbb{R}^n} \max_{d \in \mathbb{D}_k} \frac{v^Td}{\norm{v}\norm{d}}
% \end{equation}
% In order to avoid the aforementioned deterioration of polling directions, the positive spanning sets are required to satisfy the following assumption where the size of the directions does not tend to infinite or approach zero, and the cosine measure always stays positive.
% \begin{assumption} \label{ass : Dk}
% The following hold for all positive spanning sets $\mathbb{D}_k$ used for polling in Algorithm \ref{alg:dse2}. There exists a constant $k_{min} > 0$ such that $k(\mathbb{D}_k) > k_{min}$ for all $k$. There exist constants $d_{min} > 0$ and $d_{max} > 0$ such that $d_{min} \leq \norm{d} \leq d_{max}$ for all $d \in \mathbb{D}_k$.
% \end{assumption}
% In order to prove that Assumption \ref{ass:5}(2) holds, we first must link the stopping time with the unsuccessful iteration condition. Before proving the main result, we first report a technical lemma which will be used later.
% % \vspace{1em}
% % \fbox{
% %  \parbox{0.8\linewidth}{%
% %   \textcolor{red}{
% %   Non sono piu sicuro che il seguente lemma e la relativa dimostrazione funzionino. Bisogna reinquadrare un attimo il concetto di Beta-probabilistically eps-f accurate estimates. Posso ottenere questa assunzione a gratis dalle ipotesi che ho gia'?? Inoltre, nella dimostrazione bisogna sostituire $\gamma$ con $\theta$. Inoltre, bisogna poi anche capire se la violazione della sufficient decrease condition porti effettivamente a (diff) $ \geq (\gamma+2)$ etc... Vedi paper di Dzahini. Se sistemiamo questa poi abbiamo tutto.
% %   }
% % }
% % }\vspace{1em}

% % \begin{lemma} 
% % Assume that Assumption \ref{ass : Dk} hold and that $\delta^{k} \leq \delta_{\epsilon^{\prime}}$. Let $f_{0}^{k}$ and $f_{s}^{k}$ be $\varepsilon_{f}$-accurate estimates of $f\left(x^{k}\right)$ and $f\left(x^{k}+s^{k}\right)$ respectively. If $\norm{\partial_C f\left(x^{k}\right)}^{1 / \hat{p}}>\epsilon^{\prime}$, then
% % $$
% % f_{s}^{k}-f_{0}^{k} \leq-\gamma c \varepsilon_{f}\left(\delta^{k}\right)^{p}
% % $$

% % In particular, this means that the iteration $k$ of Algorithm 1 is successful.
% % \end{lemma}
% % \begin{proof}
% % Suppose $\dk \leq \delta_{\epsilon^{\prime}}$ and assume in contradiction that $\fks - \fko >-\gamma c \ef (\dk)^{p} \quad \forall s = \dk d, \quad d \in \mathbb{D}_k$ positive spanning set of iteration $k$. 
% % Choose $g \in \partial_Cf(\xk)$ and subsequently $d_k^* \in \mathbb{D}_k$ such that 
% % \begin{equation}
% %     \mathcal{K}(\mathbb{D}_k)\cdot \norm{g}\cdot\norm{d_k^*} \leq - \ip{g}{d_k^*}.
% % \end{equation}
% % Now,
% % \begin{equation}
% %     \fks - \fko > -\gamma c \ef (\dk)^p
% % \end{equation}
% % holds for $s = \dk d_k^*$ as well, so
% % \begin{equation}
% %     f(\xk+s) - f(\xk) = (f(\xk+s) - \fks) + (\fks - \fko) + (\fko - f(\xk)) 
% % \end{equation}
% % implies 
% % \begin{equation}
% %     f(\xk+s) - f(\xk) + (\gamma+2)c\ef(\dk)^{p-1} \geq 0, 
% % \end{equation}
% % which follows from the estimates being $\ef$-accurate.\\
% % From Lebourg mean value theorem for nonsmooth functions we have that $\exists\mu_k \in (0,1)$ such that 
% % \begin{equation}
% %     f(\xk+\dk d_k^*)-f(\xk) = \ip{\bar{g}}{\dk d_k^*}
% % \end{equation}
% % where $\bar{g} \in \partial_Cf(\xk+\mu_k\dk d_k^*)$, $\mu_k \in (0,1)$.
% % From the previous equation we get
% % \begin{equation}
% %     \dk \ip{\bar{g}}{d_k^*} + (\gamma+2)c\ef(\dk)^p \geq 0.
% % \end{equation}
% % Subtracting $-\ip{g}{d_k^*}$, with $g \in \partial_Cf(\xk + \mu_k\dk d_k^*)$ we get
% % \begin{equation}
% %     \mathcal{K}\norm{g}\norm{d_k^*} \leq -\ip{g}{d_k^*} \leq \ip{\dk \bar{g} -g}{d_k^*} + (\gamma+2)c\ef(\dk)^p
% % \end{equation}
% % from which we find
% % \begin{align}    
% %     \norm{g} &\leq \frac{1}{\mathcal{K}\norm{d_k^*}} \left( \ip{\bar{g} -g}{d_k^*} + (\gamma+2)c\ef(\dk)^{p-1} \right)\\
% %     &\leq \frac{1}{\mathcal{K}\norm{d_k^*}} \left( 2L_x\norm{d_k^*} + (\gamma+2)c\ef(\dk)^{p-1} \right)\\
% %     &\leq (\mathcal{K}d_{min})^{-1} \left( 2L_xd_{max} + (\gamma+2)c\ef(\dk)^{p-1} \right)\\
% % \end{align}
% % where $L_x$ is the locally Lipschitz constant for $f$ in the neighborhood of the Clarke stationary point $x$.


% % Then, dividing both sides of (36) by $\kappa\left(\mathbb{D}^{k}\right)\left\|d_{*}^{k}\right\|$ and using Assumption 6 and 5, lead to
% % $$
% % \begin{align*}
% % \left\|\nabla f\left(x^{k}\right)\right\| & \leq \kappa_{\min }^{-1}\left[L d_{\max } \delta^{k}+(\gamma+2) c \varepsilon_{f} d_{\min }^{-1}\left(\delta^{k}\right)^{p-1}\right] \\
% % & \leq \kappa_{\min }^{-1}\left(L d_{\max }+(\gamma+2) c \varepsilon_{f} d_{\min }^{-1}\right)\left(\delta^{k}\right)^{\min (p-1,1)} \tag{37}
% % \end{align*}
% % $$
% % where the inequality (37) follows from the fact that $\delta^{k} \leq \delta_{\epsilon^{\prime}}<1$. Now, recall that $\hat{p}=\min (p-1,1)$ and let $L_{1}:=\kappa_{\min }^{-1}\left(L d_{\max }+(\gamma+2) c \varepsilon_{f} d_{\min }^{-1}\right)$. Then, it follows from (37) that
% % $$
% % \begin{equation*}
% % \left\|\nabla f\left(x^{k}\right)\right\|^{1 / \hat{p}} \leq L_{1}^{1 / \hat{p}} \delta^{k} \leq L_{1}^{1 / \hat{p}} \delta_{\epsilon^{\prime}}=L_{1}^{1 / \hat{p}} \frac{\epsilon^{\prime}}{\zeta} \leq \epsilon^{\prime} \tag{38}
% % \end{equation*}

% % where the last inequality in (38) follows from (32), which achieves the proof.
% % \end{proof}


% % ---------- Geometry of the poll set ---------------------------------
% \begin{lemma}[Guaranteed acute direction]\label{lem:kappa}
% Let $\mathbb{D}=\{\pm e_1,\dots,\pm e_n\}\subset\mathbb S^{n-1}$.  For every non‑zero vector $g\in\mathbb R^n$ there exists $d\in \mathbb{D}$ such that
% \begin{equation}\label{eq:kappa}
%   -\langle g,d\rangle\;\ge\;\kappa\,\|g\|,
%   \qquad\text{where }\kappa:=1/\sqrt n.
% \end{equation}
% \end{lemma}
% \begin{proof}
% Since $\|g\|^2=\sum_{i=1}^n g_i^2$, there is an index $i^\star$ with $|g_{i^\star}|\ge\|g\|/\sqrt n$.  Choosing \[d=-\operatorname{sign}(g_{i^\star})e_{i^\star}\] gives $-\langle g,d\rangle=|g_{i^\star}|\ge\kappa\|g\|$.
% \end{proof}

% % ---------- Remarks ---------------------------------------------------
% \begin{remark}
% The constant $\kappa$ can be improved by replacing $\mathbb{D}$ with a regular simplex or any other positive–spanning set having a larger minimal angle to every vector in $\mathbb R^{n}$.
% \end{remark}

% % ---------- Main lemma (p = 2) --------------------------------------
% Key result is the following lemma, which follows the same idea of \cite{MR4359469}.
% \begin{lemma}[Probabilistic sufficient decrease with small radius]\label{lem:prob_SUCC_small_radius}
% Fix $\varepsilon > 0$ and let \[c_1 := \kappa \varepsilon - 2 L \,(>0).\]
% Define the radius and decrease constants
% \begin{align}
%   \Delta_{\varepsilon} := \min \Bigl\{ 1,\; \tfrac{\kappa \varepsilon}{4 L},\; \tfrac{c_1}{4 \varepsilon_f},\; \tfrac{c_1}{2S}\Bigr\}, \label{eq:delta_eps}
%   \quad
%   \theta := \tfrac{c_1}{2} \, \Delta_{\varepsilon}^{-1}. \label{eq:theta}
% \end{align}
% Consider an iteration $k < T_{\varepsilon}$ of Algorithm~\ref{alg:dse2} with polling radius $\Delta_k \le \Delta_{\varepsilon}$. Then
% \begin{enumerate}[label=(\alph*)]
%   \item (True decrease)  There exists $d_k \in D$ such that
%   \begin{equation}\label{eq:true_dec}
%      f(x_k + \Delta_k d_k) - f(x_k) \; \le \; - c_1 \, \Delta_k .
%   \end{equation}
%   \item (Noisy sufficient decrease)  Define
%   \begin{equation}\label{eq:good_event}
%      \mathcal G_k := \Bigl\{ |\tilde f(x_k) - f(x_k)| \le \varepsilon_f \Delta_k^{\Dexp} \Bigr\}
%                  \cap \Bigl\{ |\tilde f(x_k + \Delta_k d_k) - f(x_k + \Delta_k d_k)| \le \varepsilon_f \Delta_k^{\Dexp} \Bigr\} .
%   \end{equation}
%   On $\mathcal G_k$ the sufficient–decrease test is satisfied:
%   \begin{equation}\label{eq:suff_dec}
%      \tilde f(x_k + \Delta_k d_k) \; \le \; \tilde f(x_k) - \theta \, \Delta_k^{2} .
%   \end{equation}
% \end{enumerate}
% Consequently iteration $k$ is \emph{successful}, hence $\Delta_{k+1} \ge \gamma^{-1} \Delta_k$.  Conditional success probability is
% \begin{equation}\label{eq:prob_success}
%    \Pb \bigl( \text{iteration $k$ successful} \mid \mathcal F_{k-1} \bigr) \; \ge \; \beta^{2} ,
% \end{equation}
% so Assumption~5.1\,(2) holds with $q = \beta^{2}$ and $\lambda = -\ln \gamma$.
% \end{lemma}

% % ---------- Proof -----------------------------------------------------
% \begin{proof}
% Since $k < T_{\varepsilon}$, pick $g_k \in \partial_C f(x_k)$ with $\|g_k\| > \varepsilon$. Lemma~\ref{lem:kappa} yields a $d_k \in D$ such that $-\langle g_k, d_k \rangle \ge \kappa \varepsilon$. Let $s_k := \Delta_k d_k$.  The Lebourg mean‑value theorem gives
% \[
%    f(x_k + s_k) - f(x_k) = \langle \bar g_k, s_k \rangle
% \]
% for some $\bar g_k \in \partial_C f(x_k + \xi s_k)$ with $\xi \in (0,1)$.  Decompose:
% \[
%    \langle \bar g_k, d_k \rangle = \langle g_k, d_k \rangle + \langle \bar g_k - g_k, d_k \rangle .
% \]
% Using $-\langle g_k, d_k \rangle \ge \kappa \varepsilon$ and $\|\bar g_k - g_k\| \le 2 L$ (both norms individually $\le L$), we have
% \[
%    f(x_k + s_k) - f(x_k) \le -\kappa \varepsilon \, \Delta_k + 2 L \, \Delta_k = - c_1 \, \Delta_k ,
% \]
% which is \eqref{eq:true_dec}.\\
% On $\mathcal G_k$,
% \begin{align*}
%    \tilde f(x_k + s_k) - \tilde f(x_k) & = \bigl[f(x_k + s_k) - f(x_k)\bigr] + 2 \varepsilon_f \Delta_k^{2} \\
%    & \le - c_1 \, \Delta_k + 2 \varepsilon_f \Delta_k^{2} .
% \end{align*}
% Because $\Delta_k \le c_1/(4 \varepsilon_f)$ (definition of $\Delta_{\varepsilon}$), the quadratic term is at most $\tfrac{c_1}{2} \Delta_k$.  Hence
% \[
%    \tilde f(x_k + s_k) - \tilde f(x_k) \le - \tfrac{c_1}{2} \, \Delta_k .
% \]
% Since $\Delta_k \le \Delta_{\varepsilon}$, choosing $\theta \in [\tfrac{c_1}{4} \, \Delta_{\varepsilon}^{-1}, \; \tfrac{c_1}{2} \, \Delta_{\varepsilon}^{-1}]$ we obtain
% \[
%    - \tfrac{c_1}{2} \Delta_k \le - \theta \Delta_k^{2},
% \]
% which is precisely \eqref{eq:suff_dec}.\\
% Independence of oracle calls plus the accuracy assumption gives $\Pb(\mathcal G_k \mid \mathcal F_{k-1}) \ge \beta^{2}$.  Because all steps are deterministic once $\mathcal G_k$ occurs, we obtain \eqref{eq:prob_success}.
% \end{proof}



\section{Random Estimates Construction}
\label{appendix:random estimates}

Standard sampling strategies in stochastic derivative-free optimization often rely on averaging multiple independent realizations of the stochastic oracle. Under the classical assumptions commonly imposed on such sample-average estimators, one can establish the stronger oracle conditions typically assumed in the stochastic derivative-free optimization literature, which in turn imply Assumption~\ref{ass:quadratic tail}. The purpose of this appendix is to make this connection explicit and derive the corresponding batch-size requirements.
% In this appendix we show how the tail-bound given in
% Assumption~\ref{ass:quadratic tail} can be obtained from a standard tool from the literature, i.e., a 
% sample-average estimator. Specifically, we consider
% an unbiased stochastic oracle with finite conditional moments and derive the
% corresponding batch-size requirement ensuring that Assumption~\ref{ass:quadratic tail}
% is satisfied. 
% We then assume the following noise oracle
% \[
% \mathbb E[F(X_{k,m}^{(i)},\zeta)
% % \saraedit{\tilde f_k^{(i)}}
% \mid \mathcal F_{k-1}]=f(X_{k,m}^{(i)}),
% \qquad
% \mathbb E\left[
% \left|F(X_{k,m}^{(i)},\zeta)-f_{k,m}^{(i)}\right|^r
% % \saraedit{\tilde f_k^{(i)}-f_k^{(i)}}\right|^r
% \mid \mathcal F_{k-1}
% \right]
% \le
% \sigma_r^r,
% \]
% for some \(\sigma_r>0\), i.e., the estimates are \emph{unbiased} and the \emph{conditional $r$-th moment is finite}, with $r=r(p)=p/(p-1)$.
% % This is crucial in order to obtain, with sufficiently large probability, sufficiently accurate function estimates.
% To build the estimator, for every iteration $k$, every inner step $i$, and every direction $m$, we average \(w^{(i)}_{k}\) independent samples
% \begin{equation}\label{eq:estimates-construction}
%   % \bar f^{(i)}_{k}(x) \;:=\; \frac{1}{w^{(i)}_{k}}\sum_{j=1}^{w^{(i)}_{k}}\tilde f(x, \zeta_j).
%    \tilde f^{(i)}_{k,m}:=\frac{1}{w^{(i)}_{k}}\sum_{j=1}^{w^{(i)}_{k}}F(X^{(i)}_{k,m},\zeta_j)
% \end{equation}
% With this \emph{sample-average} estimator, %using $\tilde f^{(i)}_k = \bar f^{(i)}_k$, 
% Assumption~\ref{ass:quadratic tail} is satisfied by using 
% $w^{(i)}_k=\mathcal{O}(\Delta_k^{-2p})$ samples per point. Indeed, the following holds.
% We show here how sample-average estimates can be constructed so as to satisfy
% Assumption~\ref{ass:quadratic tail}. The argument is based on Rosenthal's
% inequality and requires a finite conditional $r$-th moment of the oracle noise.
\begin{proposition}[Rosenthal batch size for Assumption~\ref{ass:quadratic tail}]
\label{prop:rosenthal-ass33}
Let \(p\in(1,2]\) and set
\[
    r:=\frac{p}{p-1}\ge2.
\]
Let \(\mathcal H\) be a sigma-algebra, let \(X\) be an
\(\mathcal H\)-measurable random point, and let \(\Delta>0\) be an
\(\mathcal H\)-measurable random variable.
Suppose that, conditionally on \(\mathcal H\), the oracle samples
\[
    F(X,\zeta_1),\ldots,F(X,\zeta_W)
\]
are independent and satisfy, for some \(\sigma>0\),
\[
    \mathbb E\!\left[
        F(X,\zeta_j)-f(X)
        \,\middle|\,
        \mathcal H
    \right]
    =0, 
    \qquad 
    \mathbb E\!\left[
        \left|F(X,\zeta_j)-f(X)\right|^r
        \,\middle|\,
        \mathcal H
    \right]
    \le \sigma^r,
\]
for every \(j=1,\ldots,W\), with \(W\) a positive, integer-valued,
\(\mathcal H\)-measurable random variable.\\
Define the sample-average estimator
\[
    \tilde f(X)
    :=
    \frac{1}{W}
    \sum_{j=1}^{W}F(X,\zeta_j).
\]
Then there exists a constant \(C_r>0\), depending only on \(r\), such that,
if
\[
    W
    \ge
    \left(
        \frac{C_r\sigma^r}{\varepsilon_h}
    \right)^{2/r}
    \Delta^{-2p},
\]
then
\[
    \mathbb P\!\left(
        |\tilde f(X)-f(X)|
        \ge
        \alpha\Delta^p
        \,\middle|\,
        \mathcal H
    \right)
    \le
    \frac{\varepsilon_h}{\alpha^{p/(p-1)}}
    \qquad
    \text{for every }\alpha>0.
\]
Consequently, a batch size
 $W=\mathcal O(\Delta^{-2p})$
is sufficient to satisfy Assumption~\ref{ass:quadratic tail}.
% Let $p\in(1,2]$ and set
% \[
%     r:=\frac{p}{p-1}\ge 2.
% \]
% Assume that the stochastic oracle is conditionally unbiased and has uniformly bounded
% conditional $r$-th moment on the set of queried points; namely, for some
% $\sigma_r>0$,
% \[
%     \mathbb E\!\left[
%         F(x,\zeta)-f(x)
%         \,\middle|\,
%         \mathcal F_{k-1}
%     \right]=0,
%     \qquad
%     \mathbb E\!\left[
%         \left|F(x,\zeta)-f(x)\right|^r
%         \,\middle|\,
%         \mathcal F_{k-1}
%     \right]
%     \le
%     \sigma_r^r
% \]
% for every point $x$ queried at iteration $k$.\\
% For each queried point, define the sample-average estimator
% \[
%     \tilde f_{k}^{}(x)
%     :=
%     \frac{1}{w_k^{}}
%     \sum_{j=1}^{w_k^{}}
%     F(x,\zeta_j),
% \]
% where the samples are conditionally independent given $\mathcal F_{k-1}$.\\
% Then there exists a constant $C_r>0$, depending only on $r$, such that if
% \[
%     w_k^{}
%     \ge
%     \left(
%         \frac{C_r\sigma_r^r}{\varepsilon_h}
%     \right)^{2/r}
%     \Delta_{k}^{-2p},
% \]
% then
% \[
% \mathbb P\!\left(
%     |\f_{k}^{}(x)-f(x)|
%     \ge
%     \alpha \Delta_k^p
%     \,\middle|\,
%     \mathcal F_{k-1}
% \right)
% \le
% \frac{\varepsilon_h}{\alpha^{p/(p-1)}}
% \qquad
% \forall \alpha>0.
% \]
% In particular, Assumption~\ref{ass:quadratic tail} is satisfied by considering
% \[
%     w_k^{}=\mathcal O(\Delta_{k}^{-2p}).
% \]
% samples.
\end{proposition}
\begin{proof}
Define
\[
    \xi_j:=F(X,\zeta_j)-f(X),
    \qquad j=1,\ldots,W.
\]
Then
\[
    \tilde f(X)-f(X)
    =
    \frac{1}{W}\sum_{j=1}^{W}\xi_j.
\]
Conditionally on \(\mathcal H\), the random variables
\(\xi_1,\ldots,\xi_W\) are independent and centered, and satisfy
\[
    \mathbb E\!\left[
        |\xi_j|^r
        \,\middle|\,
        \mathcal H
    \right]
    \le\sigma^r.
\]
Since \(r\ge2\), the conditional version of Rosenthal's inequality gives
\[
\begin{aligned}
    \mathbb E\!\left[
        \left|
        \sum_{j=1}^{W}\xi_j
        \right|^r
        \,\middle|\,
        \mathcal H
    \right]
    % &\le
    % C_r\left[
    %     \sum_{j=1}^{w}
    %     \mathbb E\!\left[
    %         |\xi_j|^r
    %         \,\middle|\,
    %         \mathcal H
    %     \right]
    %     +
    %     \left(
    %         \sum_{j=1}^{w}
    %         \mathbb E\!\left[
    %             |\xi_j|^2
    %             \,\middle|\,
    %             \mathcal H
    %         \right]
    %     \right)^{r/2}
    % \right]
    % \\
    % &\le
    \le C_r\sigma^r W^{r/2},
\end{aligned}
\]
where \(C_r>0\) depends only on \(r\).
Therefore,
\[
    \mathbb E\!\left[
        |\tilde f(X)-f(X)|^r
        \,\middle|\,
        \mathcal H
    \right]
    \le
    C_r\sigma^r W^{-r/2}.
\]
By conditional Markov's inequality,
\[
\begin{aligned}
    &\mathbb P\!\left(
        |\tilde f(X)-f(X)|
        \ge
        \alpha\Delta^p
        \,\middle|\,
        \mathcal H
    \right)
    \le
    \frac{
        C_r\sigma^r W^{-r/2}
    }{
        \alpha^r\Delta^{pr}
    }.
\end{aligned}
\]
If
\[
    W
    \ge
    \left(
        \frac{C_r\sigma^r}{\varepsilon_h}
    \right)^{2/r}
    \Delta^{-2p},
\]
then
\[
    C_r\sigma^r W^{-r/2}\Delta^{-pr}
    \le\varepsilon_h.
\]
Hence
\[
    \mathbb P\!\left(
        |\tilde f(X)-f(X)|
        \ge
        \alpha\Delta^p
        \,\middle|\,
        \mathcal H
    \right)
    \le
    \frac{\varepsilon_h}{\alpha^r}.
\]
Since \(r=p/(p-1)\), the claim follows.
% Define
% \[
%     \xi_j(x):=F(x,\zeta_j)-f(x),
%     \qquad j=1,\ldots,w_k.
% \]
% Then
% \[
%     \tilde f_k(x)-f(x)
%     =
%     \frac{1}{w_k}
%     \sum_{j=1}^{w_k}\xi_j(x).
% \]
% Conditionally on $\mathcal F_{k-1}$, the random variables
% $\xi_1(x),\ldots,\xi_{w_k}(x)$ are independent, centered, and satisfy
% \[
%     \mathbb E\!\left[
%         |\xi_j(x)|^r
%         \,\middle|\,
%         \mathcal F_{k-1}
%     \right]
%     \le
%     \sigma_r^r.
% \]
% Since $r\ge2$, the conditional version of Rosenthal's inequality gives
% \[
%     \mathbb E\!\left[
%         \left|
%         \frac{1}{w_k}
%         \sum_{j=1}^{w_k}\xi_j(x)
%         \right|^r
%         \,\middle|\,
%         \mathcal F_{k-1}
%     \right]
%     \le
%     C_r\sigma_r^r w_k^{-r/2},
% \]
% where $C_r>0$ depends only on $r$.\\
% Now apply Markov's inequality to the nonnegative random variable
% $|\tilde f_k(x)-f(x)|^r$. For every $\alpha>0$,
% \[
% \begin{aligned}
% \mathbb P\!\left(
%     |\tilde f_k(x)-f(x)|
%     \ge
%     \alpha\Delta_k^p
%     \,\middle|\,
%     \mathcal F_{k-1}
% \right)
% &=
% \mathbb P\!\left(
%     |\tilde f_k(x)-f(x)|^r
%     \ge
%     \alpha^r\Delta_k^{pr}
%     \,\middle|\,
%     \mathcal F_{k-1}
% \right)
% \\
% &\le
% \frac{
%     \mathbb E\!\left[
%         |\tilde f_k(x)-f(x)|^r
%         \,\middle|\,
%         \mathcal F_{k-1}
%     \right]
% }{
%     \alpha^r\Delta_k^{pr}
% }
% \\
% &\le
% \frac{
%     C_r\sigma_r^r w_k^{-r/2}
% }{
%     \alpha^r\Delta_k^{pr}
% }.
% \end{aligned}
% \]
% If
% \[
%     w_k
%     \ge
%     \left(
%         \frac{C_r\sigma_r^r}{\varepsilon_h}
%     \right)^{2/r}
%     \Delta_k^{-2p},
% \]
% then
% \[
%     C_r\sigma_r^r w_k^{-r/2}\Delta_k^{-pr}
%     \le
%     \varepsilon_h.
% \]
% Therefore,
% \[
% \mathbb P\!\left(
%     |\tilde f_k(x)-f(x)|
%     \ge
%     \alpha\Delta_k^p
%     \,\middle|\,
%     \mathcal F_{k-1}
% \right)
% \le
% \frac{\varepsilon_h}{\alpha^r}.
% \]
% Finally, since $r=p/(p-1)$, this becomes
% \[
% \mathbb P\!\left(
%     |\tilde f_k(x)-f(x)|
%     \ge
%     \alpha\Delta_k^p
%     \,\middle|\,
%     \mathcal F_{k-1}
% \right)
% \le
% \frac{\varepsilon_h}{\alpha^{p/(p-1)}}.
% \]
% Thus Assumption~\ref{ass:quadratic tail} is satisfied.
\end{proof}
\begin{remark}[Application to the algorithm]
At iteration \(k\), Proposition~\ref{prop:rosenthal-ass33} is applied with
\[
    \mathcal H=\mathcal F_{k-1},
    \qquad
    \Delta=\Delta_k.
\]
For the base-point estimate, take \(X=X_k\), while for the trial-point estimate associated with direction \(m\) and extrapolation level \(i\), take $X=X_{k,m}^{(i)}.$
Thus, define
\[
    \tilde f_k^{(-1)}
    :=
    \frac{1}{W_k^{(-1)}}
    \sum_{j=1}^{W_k^{(-1)}}
    F\bigl(X_k,\zeta_{k,j}^{(-1)}\bigr)
\]
and
\[
    \tilde f_{k,m}^{(i)}
    :=
    \frac{1}{W_{k,m}^{(i)}}
    \sum_{j=1}^{W_{k,m}^{(i)}}
    F\bigl(X_{k,m}^{(i)},\zeta_{k,m,j}^{(i)}\bigr).
\]
Here,
\[
    \left\{\zeta_{k,j}^{(-1)}\right\}_{j}
    \quad\text{and}\quad
    \left\{\zeta_{k,m,j}^{(i)}\right\}_{m,i,j}
\]
denote conditionally independent copies of the generic oracle random variable \(\zeta\), given \(\mathcal F_{k-1}\).\\
These estimates satisfy Assumption~\ref{ass:quadratic tail} whenever
\[
    W_k^{(-1)}
    =
    \mathcal O(\Delta_k^{-2p}),
    \qquad
    W_{k,m}^{(i)}
    =
    \mathcal O(\Delta_k^{-2p}),
\]
uniformly over \(m\in[1:\bar m]\) and \(i\in[0:\Imax]\).\\
Moreover, using fresh conditionally independent batches for the base point and for all potential trial points ensures that
Assumption~\ref{ass:indep} is also satisfied.
\end{remark}
\begin{remark}
For $p=2$, one has $r=2$, and the result reduces to the usual finite-variance
case with
$W_k=\mathcal O(\Delta_k^{-4}).$
For $p\in(1,2)$, one has $r>2$, so a moment assumption stronger than finite
variance is needed in order to obtain the tail exponent required by
Assumption~\ref{ass:quadratic tail}.
\end{remark}



% \section{Random Estimates Construction} \label{appendix:random estimates}


% The following results show how many samples are needed to build $\beta$-probabilistically $\varepsilon_f$-accurate estimates, based on sample average as in \eqref{eq:estimates-construction}.
% \begin{proposition}[Chebyshev batch size depending on \(\Delta_k\)]
% \label{prop:cheb-delta}
% Let \(\Dexp>1\), pick a confidence level \(\beta\in(1/2,1)\) and an
% accuracy factor \(\varepsilon_f>0\).
% If
% \begin{equation}\label{eq:cheb-delta}
%    m_k
%    \;\;\ge\;\;
%    \frac{\sigma^{2}}
%         {(1-\beta)\,\varepsilon_f^{2}}\, \Delta_k^{-2 \Dexp},
% \end{equation}
% then
% \[
%    \Pb\bigl(
%        |\bar f_{m_k}(x_k)-f(x_k)|
%        \le\varepsilon_f\,\Delta_k^{\Dexp}
%        \,\big|\,
%        \F_{k-1}\bigr)
%    \;\ge\;
%    \beta .
% \]
% Namely, the estimates are $\beta$-probabilistically $\varepsilon_f$-accurate.
% \end{proposition}
% \begin{proof}
% Chebyshev's inequality with threshold
% \(t=\varepsilon_f\Delta_k^{\Dexp}\) yields
% \[
%    \Pb\Bigl(
%      |\bar f_{m_k}-f|>t
%      \,\big|\,
%      \F_{k-1}\Bigr)
%    \;\le\;
%    \frac{\sigma^{2}}{m_k\,\varepsilon_f^{2}\,\Delta_k^{2\Dexp}}.
% \]
% Demanding that the right-hand side does not exceed \(1-\beta\) gives
% \eqref{eq:cheb-delta}.
% \end{proof}
% Equation~\eqref{eq:cheb-delta} shows that, for fixed noise level
% \(\sigma^{2}\), shrinking \(\varepsilon_f\) or tightening the confidence
% \(1-\beta\) both increase the required batch size at rate
% \(\varepsilon_f^{-2}\) or \((1-\beta)^{-1}\), respectively.
% \begin{remark}[Rosenthal bound for finite $r$-th moment]
% Assume instead that the noise only has a finite conditional
% $r$-th moment for some $r>2$, i.e.\
% $\me\bigl[\,|\zeta|^{\,r}\!\mid\F_{k-1}\bigr]\le\sigma^{\,r}$
% with $\sigma>0$.  
% Rosenthal’s inequality (conditional version) yields
% \[
%   \Pb\bigl(|\bar f_{m_k}-f|>\varepsilon_f\Delta_k^{\Dexp}\;|\;\F_{k-1}\bigr)
%   \;\le\;
%   \frac{C_r\,\sigma^{\,r}}
%        {m_k^{\,r/2}\,\varepsilon_f^{\,r}\D_k^{\Dexp\,r}},
% \]
% where $C_r$ is an explicit constant, depending on $r$.
% Insisting that this probability is $\le 1-\beta$ gives the
% sample-size rule
% \[
%    m_k
%    \;\ge\;
%    \Bigl[
%      \frac{C_r\,\sigma^{\,r}}
%           {(1-\beta)\,\varepsilon_f^{\,r}}
%    \Bigr]^{\!2/r}\,
%    \Delta_k^{-2\Dexp}.
% \]
% Because the exponent $2/r$ is $<1$, Rosenthal can require
% \emph{fewer} samples than Chebyshev.
% \end{remark}
% \begin{remark}[Sub-Gaussian tail reduces \(m_k\)] \label{rem:subgauss-tails}
% If the noise is \emph{sub-Gaussian}, i.e.\
% \(\me[e^{\lambda\zeta}\mid\F_{k-1}]\le
%  \exp(\lambda^{2}\sigma^{2}/2)\) for all \(\lambda\in\R\),
% then a Chernoff bound gives
% \[
%    \Pb\bigl(|\bar f_{m_k}-f|>t \;|\;\F_{k-1}\bigr)
%    \;\le\;2\exp\!\bigl(-m_k t^{2}/(2\sigma^{2})\bigr).
% \]
% Replacing \(t=\varepsilon_f\Delta_k^{\Dexp}\) and solving for \(m_k\)
% yields the sharper rule
% \[
%    m_k
%    \;\ge\;
%    \frac{2\sigma^{2}}{\varepsilon_f^{2}}\,
%    \ln\!\Bigl(\tfrac{2}{1-\beta}\Bigr)\,
%    \Delta_k^{-2\Dexp},
% \]
% where the factor \(1/(1-\beta)\) in \eqref{eq:cheb-delta} is upgraded to
% \(\ln\!\bigl(2/(1-\beta)\bigr)\).
% \end{remark}
% \begin{proof}
% Let $\bar\zeta_k:=\frac1{m_k}\sum_{j=1}^{m_k}\zeta_{k,j}$ be the mini-batch noise.
% By conditional independence and the sub-Gaussian mgf bound,
% \[
% \EE\!\big[e^{\lambda \bar\zeta_k}\mid\F_{k-1}\big]
% =\prod_{j=1}^{m_k}\EE\!\Big[e^{\frac{\lambda}{m_k}\zeta_{k,j}}\ \Big|\ \F_{k-1}\Big]
% \ \le\ \prod_{j=1}^{m_k}\exp\!\Big(\frac{\lambda^2\sigma^2}{2m_k^2}\Big)
% =\exp\!\Big(\frac{\lambda^2\sigma^2}{2m_k}\Big).
% \]
% Thus $\bar\zeta_k$ is (conditionally) sub-Gaussian with proxy $\sigma^2/m_k$. By Chernoff’s method,
% for any $t>0$,
% \[
% \Pb(\bar\zeta_k\ge t\mid\F_{k-1})
% =\Pb(e^{\lambda \bar\zeta_k}\ge e^{\lambda t}\mid\F_{k-1})
% \ \le\ e^{-\lambda t}\,\EE[e^{\lambda \bar\zeta_k}\mid\F_{k-1}]
% \ \le\ \exp\!\Big(-\lambda t+\frac{\lambda^2\sigma^2}{2m_k}\Big).
% \]
% Minimizing the RHS over $\lambda>0$ gives $\lambda^\star = m_k t/\sigma^2$ and
% \[
% \Pb(\bar\zeta_k\ge t\mid\F_{k-1})\ \le\ \exp\!\Big(-\,\frac{m_k t^2}{2\sigma^2}\Big).
% \]
% By symmetry, the same bound holds for $\Pb(\bar\zeta_k\le -t\mid\F_{k-1})$, hence
% \[
% \Pb\big(|\bar\zeta_k|>t\mid\F_{k-1}\big)\ \le\ 2\exp\!\Big(-\,\frac{m_k t^2}{2\sigma^2}\Big).
% \]
% Since $\bar f_{m_k}-f(X_k)=\bar\zeta_k$, the tail bound follows. Setting $t=\varepsilon_f\Delta_k^{\,\Dexp}$ and
% requiring the two-sided tail to be at most $1-\beta$,
% \[
% 2\exp\!\Big(-\,\frac{m_k\,\varepsilon_f^2\Delta_k^{2\Dexp}}{2\sigma^2}\Big)\ \le\ 1-\beta
% \quad\Longleftrightarrow\quad
% m_k \ \ge\ \frac{2\sigma^{2}}{\varepsilon_f^{2}}\,
% \ln\!\Big(\frac{2}{1-\beta}\Big)\,\Delta_k^{-2\Dexp}.
% \]
% This yields the stated batch-size rule (take the ceiling to make $m_k\in\mathbb N$).
% \end{proof}
% \begin{remark}[Sampling cost grows like $\mathcal O(\Delta_k^{-2\Dexp})$]
% \label{rem:batch-scaling}
% The adaptive mini-batch rules in Chebyshev, Chernoff, Rosenthal, all take the generic form
% \[
%    m_k \;=\;
%    C\cdot\;
%    \Delta_k^{-2\Dexp},
%    \qquad \Dexp>1,
% \]
% with $C$ constant. Thus, the random oracle merely changes the \emph{constant} factor.
% % Hence, even in the stochastic setting, the number of function evaluations required by one poll step scales with exactly the same
% % exponent $2p$ that appears in deterministic direct search. 
% \end{remark}
% %--------------------------------------------------------------------------
% % Remark: How to satisfy Assumption via a batch-size floor
% %--------------------------------------------------------------------------
% % \begin{remark}[Enforcing \eqref{eq:postL1_assumption} with a batch-size floor]\label{rem:floor_policy}
% % Assumption~\ref{ass:postL1} is guaranteed if the line-search \emph{accepts} level $i$ only after the
% % per-point mini-batch size $m_{k,i}$ reaches a deterministic floor $M_k=M_k(\Delta_k)$ depending on $\Delta_k$, i.e. accept level $i$ only if
% % \begin{equation}
% % \begin{aligned}\label{eq:accept_with_floor}
% % & \tilde f(X_k)-\tilde f(X_k^{(i)})\ \ge\ \theta\,\Delta_{k,i}^{\Dexp}\\
% % & m_{k,i}\ \ge\ M_k(\Delta_k).
% % \end{aligned}
% % \end{equation}
% % Assume the oracle returns $\tilde f(x)=f(x)+\zeta$ with $\me[\zeta\mid\F_{k-1}]=0$ and
% % $\operatorname{Var}(\zeta\mid\F_{k-1})\le\sigma^2<\infty$, and that the two sample means at
% % $X_k$ and $X_k^{(i)}$ use independent batches of size $m_{k,i}$. On the \underline{last} acceptance event $A_i:=\{H_k=i\}$ we
% % have $m_{k,i}\ge M_k(\Delta_k)$ deterministically and
% % \[
% % \me\!\left[E_{k,i}\mid\F_{k-1},A_i\right]=0,
% % \qquad
% % \operatorname{Var}\!\left(E_{k,i}\mid\F_{k-1},A_i\right)=\frac{2\sigma^2}{m_{k,i}}
% % \ \le\ \frac{2\sigma^2}{M_k}.
% % \]
% % By Jensen/Cauchy--Schwarz,
% % \begin{equation}\label{eq:L1_from_floor}
% % \me\!\left[\big|E_{k,i}\big|\,\middle|\,\F_{k-1},A_i\right]
% % \ \le\
% % \sqrt{\frac{2\sigma^2}{M_k}}.
% % \end{equation}
% % Hence, choosing
% % \begin{equation}\label{eq:floor_choice}
% % M_k\ \ge\ \frac{2\,\sigma^2}{c^2}\,\Delta_k^{-2\Dexp}
% % \end{equation}
% % immediately yields the target bound \eqref{eq:postL1_assumption}. In particular, the sampling
% % cost is $\Theta(\Delta_k^{-2\Dexp})$.
% % \end{remark}

% \begin{remark}[Post-selection $L^1$ via tails of $H_k$]
% \label{rem:postsel_tail}
% For each level $i\ge 0$ define the (difference) estimation error
% \[
% E_{k,i}
% :=\big(\tilde f(X_k)-\tilde f(X_k^{(i)})\big)-\big(f(X_k)-f(X_k^{(i)})\big),
% \]
% and let $H_k$ be the random selected level from the inner line-search (success-until-first-failure rule with steps $\Delta_{k,i}=\gamma^{-i}\Delta_k$). Then, for a suitable batch schedule $m_k = O\left(\Delta_k^{-2\Dexp}(1+\log(1/\Delta_k))^2\right)$, it holds
% \begin{equation}
% \label{eq:PSL1}
% \me\!\left[\,|E_{k,H_k}|\,\middle|\,\F_{k-1}\right]\ \le\ c\,\Delta_k^{\Dexp},
% \end{equation}
% with a constant $c>0$ independent of $k$.
% \end{remark}
% \begin{proof}
% By the tower property and Cauchy--Schwarz,
% \begin{align}
% \me\!\left[\,|E_{k,H_k}|\,\middle|\,\F_{k-1}\right]
% &=\sum_{i\ge 0}\me\!\left[\,|E_{k,i}|\,\mathbf 1_{\{H_k=i\}}\,\middle|\,\F_{k-1}\right] \notag \\
% &\le \sum_{i\ge 0}\Big(\me\!\left[E_{k,i}^{2}\,\middle|\,\F_{k-1}\right]\Big)^{\!1/2}
% \Big(\Pb(H_k=i\,|\,\F_{k-1})\Big)^{\!1/2}. \label{eq:CS_selected}
% \end{align}
% Assume the additive oracle model $\tilde f=f+\zeta$ with $\me[\zeta\,|\F_{k-1}]=0$ and $\mathrm{Var}(\zeta\,|\F_{k-1})\le\sigma^2$, and let $m_k$ be the number of i.i.d. samples used to build the function estimates at each point (same $m_k$ for all inner points at iteration $k$). The two means in $E_{k,i}$ are based on independent batches, hence
% \begin{equation}
% \label{eq:Ekivar}
% \me\!\left[E_{k,i}^{2}\,\middle|\,\F_{k-1}\right]\ \le\ \frac{2\sigma^2}{m_k}
% \qquad\text{for all }i\ge 0.
% \end{equation}
% Define the tail constant
% \begin{equation}
% \label{eq:Ctail_def}
% C_{\text{tail}}(k)\ :=\ \sum_{i\ge 0}\sqrt{\Pb(H_k=i\, | \,\F_{k-1})}.
% \end{equation}
% Plugging \eqref{eq:Ekivar} into \eqref{eq:CS_selected} yields
% \begin{equation}
% \label{eq:PSL1_pre}
% \me\!\left[\,|E_{k,H_k}|\,\middle|\,\F_{k-1}\right]
% \ \le\ \sqrt{\tfrac{2\sigma^2}{m_k}}\;C_{\text{tail}}(k).
% \end{equation}
% By Lemma~10.1 (tail of the inner linesearch), with
% \[
% A_k\ :=\ \frac{2\mu_h\,\Delta_k^{\Dexp}+2f_{\max}}{\theta\,\Delta_k^{\Dexp}},
% \qquad \text{and}\qquad
% \Pb(H_k\ge i\, |\,\F_{k-1})\ \le\ \min\{1,\ A_k\,\gamma^{\Dexp i}\}\quad(i\ge 0),
% \]
% define $J_k:=\max\!\left\{0,\ \left\lfloor \tfrac{1}{\Dexp}\log_{1/\gamma}(A_k)\right\rfloor\right\}$ so that $A_k\gamma^{\Dexp J_k}\le 1$. Using $\Pb(H_k=i)\le\Pb(H_k\ge i)$ and splitting the sum at $J_k$,
% \begin{align}
% C_{\text{tail}}(k)
% &\le \sum_{i=0}^{J_k-1}1\;+\;\sum_{i\ge J_k}\sqrt{A_k}\,\gamma^{\frac{\Dexp}{2}i}
% \ \le\ J_k\ +\ \frac{\sqrt{A_k}\,\gamma^{\frac{\Dexp}{2}J_k}}{1-\gamma^{\Dexp/2}}
% \ \le\ J_k\ +\ \frac{1}{1-\gamma^{\Dexp/2}}.
% \label{eq:Ctail_bound}
% \end{align}
% (Here $f_{\max}<\infty$ follows from compactness of the search region / level set.)
% Moreover, the explicit envelope
% \begin{equation}
% \label{eq:Jk_simple}
% J_k\ \le\ \frac{1}{\Dexp\log(1/\gamma)}\,
% \log\!\Big(\frac{2\mu_h}{\theta}+\frac{2f_{\max}}{\theta\,\Delta_k^{\Dexp}}\Big)\ +\ 1
% \end{equation}
% shows $J_k=O(\log(1/\Delta_k))$ as $\Delta_k\downarrow 0$.
% Combining \eqref{eq:PSL1_pre} and \eqref{eq:Ctail_bound}, a sufficient condition for \eqref{eq:PSL1} is
% \begin{equation}
% \label{eq:mk_exact}
% m_k\ \ge\ \frac{2\sigma^2}{c^2}\,\Big(J_k+\tfrac{1}{1-\gamma^{\Dexp/2}}\Big)^{\!2}\ \Delta_k^{-2 \Dexp}.
% \end{equation}
% % Using \eqref{eq:Jk_simple} to upper bound $J_k$, one obtains the simpler, uniform schedule
% % \begin{equation}
% % \label{eq:mk_simple}
% % m_k\ =\ m_0\,\Delta_k^{-2p}\Big(1+\log\tfrac{1}{\Delta_k}\Big)^2,
% % \qquad
% % m_0\ \ge\ \frac{2\sigma^2}{c^2}\,
% % \Big(1+\tfrac{1}{1-\gamma^{p/2}}+\tfrac{1}{p\log(1/\gamma)}\Big)^{\!2},
% % \end{equation}
% % which also guarantees \eqref{eq:PSL1} for all $k$.
% In particular, $m_k = O\left(\Delta_k^{-2\Dexp}(1+\log(1/\Delta_k))^2\right)$.
% \end{proof}
% \input{sections/appendix_delta}

%\pagebreak
% ---------- End -------------------------------------------------------

\bibliographystyle{plain}
\bibliography{sdse}
	
\end{document}
